% paralleles.tex — f163v, vv. 80–122 : loci paralleli. Généré par .travail/f163v_lp_redac/rediger.py
% à partir des collectes vérifiées (paralleles_80-101.json, paralleles_102-122.json) et des parallèles
% d'Austin–Olson 2004 retrouvés dans les sources ; règles : edition/CONVENTIONS.md § 4
\lp{80 ἐπείπερ]}{Ar.~\textit{Ach.}~495 ἐπειδήπερ αὐτὸς αἱρεῖ, λέγε ; Ar.~\textit{Nub.}~1412 τύπτειν τ᾽, ἐπειδήπερ γε τοῦτ᾽ ἐστ᾽ εὐνοεῖν τὸ τύπτειν ; Ar.~\textit{Thesm.}~762 ἀλλ᾽ ἐπειδήπερ πάρει ; Eur.~\textit{Hel.}~1237 πρός νύν σε γονάτων τῶνδ᾽, ἐπείπερ εἶ φίλος ; Soph.~\textit{OT}~1003 ἐπείπερ εὔνους ἦλθον, ἐξελυσάμην ; Suid.~αι~302 Αἱρῇ· βούλει. ἐπείπερ αὐτὸς αἱρῇ λέγειν ; \textit{idem} Cramer, \textit{AP} IV~103.12 (βούλῃ … λέγειν. Ἀριστοφάνης) ; Ar.~\textit{Vesp.}~1129 ἀλλ᾽ οὖν πεπειράσθω γ᾽, ἐπειδήπερ γ᾽ ἅπαξ ; Ar.~\textit{Av.}~1360 οὐδέν γ᾽. ἐπειδήπερ γὰρ ἦλθες ὦ μέλε ; Eur.~\textit{Hec.}~1286 ἐπείπερ οὕτω καὶ λίαν θρασυστομεῖ}
\lp{80 ἐπείπερ ἐστὶ (ἐπεὶ τρίτη ᾽στὶ R)]}{Ar.~\textit{Vesp.}~1507 μὰ τὸν Δί᾽ οὐδέν γ᾽ ἄλλο πλήν γε καρκίνους}
\lp{80 Θεσμοφορίων]}{Com. adesp. fr.~1132.1–2~K.-A. (P.Oxy.~3540) … ιων ἡ τρίτη — … μητρος τροφός ; Ar.~fr.~331~K.-A. (Thesm. II) δαίμων περὶ τὴν Δήμητραν, ἣν προλογίζουσαν ἐν ταῖς ἑτέραις Θεσμοφοριαζούσαις ἐποίησεν ; Is.~3.80 ἠναγκάζετο ἂν ὑπὲρ τῆς γαμετῆς γυναικὸς καὶ θεσμοφόρια ἑστιᾶν τὰς γυναῖκας ; Σ Luc.~\textit{DMeretr.}~2.1 Θεσμοφορία ἑορτὴ Ἑλλήνων μυστήρια περιέχουσα, τὰ δὲ αὐτὰ καὶ Σκιρροφορία καλεῖται ; Hsch.~κ~472 Καλλιγένειαν· οὐ τὴν γῆν, ἀλλὰ τὴν Δήμητραν ; P.Oxy.~1802 fr.~3 ii ὅθεν καὶ μελίσσας ἔκτοτε — κληθῆναι τὰς θεσμοφοριοιζούσας {κληθῆναι} γυναῖκας ; Ar.~\textit{Thesm.}~299–300 καὶ τῷ Πλούτῳ καὶ τῇ Καλλιγενείᾳ — καὶ τῇ Κουροτρόφῳ ; Hsch.~κ~4816 Κωλιάς· Ἀφροδίτης ἐπὶ Κωλιάδος ἐστὶν ἱερὸν ἐν τῇ Ἀττικῇ … ἔστι δὲ καὶ Δήμητρος ἱερὸν αὐτόθι πολύστυλον ; Hsch.~σ~1825 στήνια· ἑορτὴ Ἀθήνησι· καὶ διασκώπτουσι καὶ λοιδοροῦσιν}
\lp{80 Θεσμοφορίων ἡ μέση]}{\textit{Test.} Hsch.~τ~1440 τρίτη Θεσμοφορίων· ζητεῖται πῶς ἅμα μὲν λέγει τρίτην Θεσμοφορίων εἶναι, ἅμα δὲ μέσην, τεττάρων οὐσῶν ἡμερῶν ; Ar.~\textit{Av.}~1519–1520 ἀλλ᾽ ὡσπερεὶ Θεσμοφορίοις νηστεύομεν / ἄνευ θυηλῶν ; Ar.~\textit{Thesm.}~375–376 ἐκκλησίαν ποιεῖν ἕωθεν τῇ μέσῃ / τῶν Θεσμοφορίων, ᾗ μάλισθ᾽ ἡμῖν σχολή ; Ar.~\textit{Thesm.}~949 Παύσων σέβεται καὶ νηστεύει ; Ar.~\textit{Thesm.}~984 νηστεύωμεν δὲ πάντως ; Alciphr.~2.37 ἡ μὲν οὖν Ἄνοδος κατὰ τὴν πρώτην γέγονεν ἡμέραν, ἡ Νηστεία δὲ τὸ τήμερον εἶναι παρ᾽ Ἀθηναίοις ἑορτάζεται, τῇ Καλλιγενείᾳ δὲ εἰς τὴν ἐπιοῦσαν θύουσιν ; Ath.~7.307f ἄνδρες φίλοι, μὴ καὶ ἡμεῖς ἄγομεν Θεσμοφορίων τὴν μέσην, ὅτι δίκην κεστρέων νηστεύομεν ; Σ Ar.~\textit{Thesm.}~80 ἑνδεκάτη Πυανεψιῶνος ἄνοδος. δωδεκάτη νηστεία, ἐν ᾗ σχολὴν ἄγουσαι ὑπόκεινται αἱ γυναῖκες ἐκκλησιάζουσαι περὶ Εὐριπίδου. τρισκαιδεκάτη καλλιγένεια ; Σ Ar.~\textit{Thesm.}~80 (Callimaque) τοῦτο τῶν ζητουμένων ἐστί, πῶς καὶ τρίτην καὶ μέσην εἶπεν … δεκάτη ἐν Ἁλιμοῦντι Θεσμοφόρια ἄγεται, ὥστε τρίτην μὲν ἀπὸ δεκάτης ιβ΄ εἶναι, μέσην δὲ μὴ συναριθμουμένης τῆς δεκάτης. τοῦτο δὲ αἰνιγματῶδες κατὰ Καλλιμάχου ἄν τις φαίη ; Zen.~4.52 Κεστρεὺς νηστεύει: παροιμία ἐπὶ τῶν πάνυ λαιμάργων ; \textit{idem} Diogenian.~5.53 (om. παροιμία … πάνυ ; add. ὑποκρινομένων δὲ νηστεύειν· ἄπληστον γὰρ τὸ ζῶον), Plu.~\textit{Prov.Alex.}~1.8 (ἐπὶ τῶν δικαιοπραγούντων μὲν, οὐδὲν δὲ πλέον ἐκ τῆς δικαιοσύνης ἐπιφερομένων pro παροιμία … λαιμάργων), Greg. Cypr.~2.90 (om. παροιμία … πάνυ ; add. ἄπληστον δὲ τὸ ζῶον) ; Σ Ar.~\textit{Thesm.}~80 (ὅπου γε λιμώττουσιν) καὶ ὅπου γε λιμώττουσιν, ἀστεϊζόμενοι τὴν μέσην τῶν Θεσμοφορίων ἄγειν φασίν, ἐπεὶ αὕτη ἡ νηστεία ; Ar.~\textit{Thesm.}~570 τὸν σησαμοῦνθ᾽ ὃν κατέφαγες, τοῦτον χεσεῖν ποιήσω ; Ar.~fr.~159~K.-A. ἆρ᾽ ἔνδον ἀνδρῶν κεστρέων ἀποικία; / ὡς μὲν γάρ εἰσι νήστιδες, γιγνώσκετε ; Diph.~fr.~53~K.-A. οὗτοι δεδειπνήκασιν· ὁ δὲ τάλας ἐγὼ / κεστρεὺς ἂν εἴην ἕνεκα νηστείας ἄκρας ; \textit{h.Cer.}~49–50,~200–201 οὐδέ ποτ᾽ ἀμβροσίης καὶ νέκταρος ἡδυπότοιο / πάσσατ᾽ ἀκηχεμένη … ἀλλ᾽ ἀγέλαστος ἄπαστος ἐδητύος ἠδὲ ποτῆτος / ἧστο πόθῳ μινύθουσα βαθυζώνοιο θυγατρός ; Plu.~\textit{De\,Is.\,et\,Os.}~378e καὶ γὰρ Ἀθήνησι νηστεύουσιν αἱ γυναῖκες ἐν Θεσμοφορίοις χαμαὶ καθήμεναι ; Plu.~\textit{Dem.}~30.4 ἐν ᾗ τὴν σκυθρωποτάτην τῶν Θεσμοφορίων ἡμέραν ἄγουσαι παρὰ τῇ θεῷ νηστεύουσιν αἱ γυναῖκες ; Paus.~1.31.1 Ἀλιμουσίοις ⟨μὲν⟩ Θεσμοφόρου Δήμητρος καὶ Κόρης ἐστὶν ἱερόν ; Hsch.~α~5234 ἄνοδος· ἀνάβασις. ἡ ἑνδεκάτη τοῦ Πυανεψιῶνος, ὅτε αἱ γυναῖκες ἀνέρχονται εἰς Θεσμοφόριον, οὕτω καλεῖται ; Hsch.~κ~2384 κεστρεῖς· τοὺς κεχηνότας καὶ πεινῶντας κεστρεῖς λέγουσι. καὶ τοὺς Ἀθηναίους οὕτως ἔλεγον καὶ προσηγόρευον· τὸ γὰρ ζῷον αὐτὸ λαίμαργόν τέ ἐστι καὶ ἄπληστον ; Hsch.~τ~1374 τριήμερος· Θεσμοφόρια ὑπὸ Λακώνων ; Phot.~θ~134 Θεσμοφορίων ἡμέραι ιδ΄· δεκάτῃ Θεσμοφόρια, ἑνδεκάτῃ κάθοδος, δωδεκάτῃ νηστεία, τρισκαιδεκάτῃ καλλιγένεια ; Suid.~ν~345 Νηστεία: ἐκ τοῦ νηστεύω. καὶ Νήστιμος, ἡ τῆς νηστείας ἡμέρα}
\lp{81 τοῦτ᾽ αὐτὸ γάρ τοι]}{Ar.~\textit{Lys.}~46 ταῦτ᾽ αὐτὰ γάρ τοι κἄσθ᾽ ἃ σώσειν προσδοκῶ ; Ar.~\textit{Eq.}~180 δι᾽ αὐτὸ γάρ τοι τοῦτο καὶ γίγνει μέγας ; Ar.~\textit{Ran.}~73–74 τοῦτο γάρ τοι καὶ μόνον / ἔτ᾽ ἐστὶ λοιπὸν ἀγαθόν, εἰ καὶ τοῦτ᾽ ἄρα ; Eur.~\textit{Andr.}~906 τοῦτ᾽ αὐτὸ καὶ νοσοῦμεν· εὖ μ᾽ ὑπηγάγου ; Ar.~\textit{Nub.}~1499 ἀπολεῖς ἀπολεῖς. — τοῦτ᾽ αὐτὸ γὰρ καὶ βούλομαι ; Ar.~\textit{Thesm.}~171 ἅπασ᾽ ἀνάγκη· ταῦτα γάρ τοι γνοὺς ἐγὼ ; Eur.~\textit{El.}~261 τοῦτ᾽ αὐτὸ ταρβῶν, πρὸς δὲ καὶ σώφρων ἔφυ ; Eur.~\textit{Supp.}~1067 τοῦτ᾽ αὐτὸ χρῄζω, πάντας Ἀργείους μαθεῖν ; Soph.~\textit{OT}~1013 τοῦτ᾽ αὐτό, πρέσβυ, τοῦτό μ᾽ εἰσαεὶ φοβεῖ}
\lp{81 κἀπολεῖν με προσδοκῶ]}{Ar.~\textit{Nub.}~1136 θείς μοι πρυτανεῖ᾽ ἀπολεῖν μέ φησι κἀξολεῖν ; Ar.~\textit{Thesm.}~77 εἴτ᾽ ἔστ᾽ ἔτι ζῶν εἴτ᾽ ἀπόλωλ᾽ Εὐριπίδης ; Hdt.~7.156.2 προσδοκῶντας ἀπολέεσθαι διὰ τοῦτο ; Diogenian.~1.52 Αἲξ τὴν μάχαιραν: ἐπὶ τῶν ἑαυτοῖς λανθανόντως κακόν τι ἐπισπωμένων ; Diogenian.~7.52 Προμέρου κύνες: τοῦτόν φασιν οἰκέτην ὄντα βασιλικὸν μισῆσαι καθ᾽ ὑπερβολὴν τῶν τραγῳδιῶν τὸν ποιητὴν Εὐριπίδην, ἐπαφεῖναί τε αὐτῷ κύνας ἀγρίους, οἵ τοῦτον κατεθοινήσαντο ; J.~\textit{AJ}~5.59 ὡς ὑπὸ μὲν τούτων ἀπολεῖσθαι προσδοκᾶν, ὑπὸ δὲ τῶν ἐπ᾽ ὀλέθρῳ τοῦ Χαναναίων γένους στρατευσάντων σωθήσεσθαι ; Hsch.~δ~2106 δοκάζει· μένει. ἐπιτηρεῖ. δοκεῖ. προσδοκᾷ ; App. Prov.~5.9 Ὗς ὑπέροπλον δραμεῖται: ἐπὶ τῶν ἑαυτοὺς εἰς ὄλεθρον ἐμβαλλόντων}
\lp{81 προσδοκῶ]}{Ar.~\textit{Thesm.}~846 ἰλλὸς γεγένημαι προσδοκῶν· ὁ δ᾽ οὐδέπω ; Eur.~\textit{Supp.}~99 ἡμῶν δ᾽ ἀκούειν· προσδοκῶ τι γὰρ νέον ; Aesch.~\textit{Ag.}~674–675 Μενέλεων γὰρ οὖν / πρῶτόν τε καὶ μάλιστα προσδόκα μολεῖν}
\lp{82 αἱ γὰρ γυναῖκες ἐπιβεβουλεύκασί μοι]}{Satyr.~\textit{Vit.\,Eur.} (P.Oxy.~1176) fr.~39 col. x~23–33 αἱ δὲ γυναῖκες ἐπισυνέστησαν αὐτῷ τοῖς Θεσμοφορίοις καὶ ἀθρόαι παρῆσαν ἐπὶ τὸν τόπον ἐν ᾧ [σ]χολάζων [ἐ]τύγχανεν ; Zen.~4.91 Λήμνιον κακόν: παροιμία, ἣν διαδοθῆναι φασὶν ἀπὸ τῶν παρανομηθέντων εἰς τοὺς ἄνδρας ἐν Λήμνῳ ὑπὸ τῶν γυναικῶν ; Eur.~\textit{IT}~1032 δειναὶ γὰρ αἱ γυναῖκες εὑρίσκειν τέχνας ; Eur.~\textit{Med.}~407–409 πρὸς δὲ καὶ πεφύκαμεν / γυναῖκες, ἐς μὲν ἔσθλ᾽ ἀμηχανώταται, / κακῶν δὲ πάντων τέκτονες σοφώταται ; Diogenian.~4.4 Γυναικὶ μὴ πίστευε, μηδ᾽ ἂν ἀποθάνῃ: ὅτι δεῖ γυναιξὶ μὴ πιστεύειν ; App. Prov.~1.19 Ἄλλος ἐχθρὸς πολέμιος καὶ ἄλλη γυνὴ ἡ πολεμία ; Pl.~\textit{Smp.}~203b ἡ οὖν Πενία ἐπιβουλεύουσα διὰ τὴν αὑτῆς ἀπορίαν παιδίον ποιήσασθαι ἐκ τοῦ Πόρου ; D.S.~20.37.6 τῶν γυναικῶν τινας ἐκόλασεν ὡς ἐπιβεβουλευκυίας}
\lp{82 ἐπιβεβουλεύκασί μοι]}{Ar.~\textit{Thesm.}~335–337 εἴ τις ἐπιβουλεύει τι τῷ δήμῳ κακὸν / τῷ τῶν γυναικῶν ἢ ᾽πικηρυκεύεται / Εὐριπίδῃ ; Soph.~\textit{OT}~618–619 ὅταν ταχύς τις οὑπιβουλεύων λάθρᾳ / χωρῇ ; Luc.~\textit{Fug.}~3 ἆρα μὴ οἱ ἰδιῶται αὖθις ἐπιβεβουλεύκασί σοι ; Hsch.~ε~4316 ἐπιβουλεύων· λοιδορῶν, ἐφυβρίζων ; Ar.~\textit{Pax}~407 ὑμῖν ἐπιβουλεύοντε πολὺν ἤδη χρόνον ; Ar.~\textit{Plu.}~570 ἐπιβουλεύουσί τε τῷ πλήθει καὶ τῷ δήμῳ πολεμοῦσιν ; Aesch.~\textit{Sept.}~29 νυκτηγορεῖσθαι κἀπιβουλεύσειν πόλει ; Soph.~\textit{Aj.}~726 τὸν τοῦ μανέντος κἀπιβουλευτοῦ στρατοῦ ; D.~33.3 ὁ μὲν οὖν μετὰ τούτου μοι ἐπιβεβουλευκὼς καὶ τὸν ἀγῶνα τουτονὶ κατεσκευακὼς ; Hsch.~δ~2163 δολοπεύει· ἐπιβουλεύει, ἐνεδρεύει ; Hsch.~ε~4650 ἐπιβουλεύει· ὁρμᾷ. ζημιοῖ}
\lp{83 κἀν Θεσμοφόροιν]}{Ar.~\textit{Thesm.}~277–278 ὡς τὸ τῆς ἐκκλησίας / σημεῖον ἐν τῷ Θεσμοφορείῳ φαίνεται ; Ar.~\textit{Eccl.}~442–443 κοὔτε τἀπόρρητ᾽ ἔφη / ἐκ Θεσμοφόροιν ἑκάστοτ᾽ αὐτὰς ἐκφέρειν ; D.~21.9 λέγων, ὥσπερ ἠκούσατε, ποιεῖν τὴν ἐκκλησίαν ἐν Διονύσου ; Σ Ar.~\textit{Thesm.}~83 κἀν Θεσμοφόροιν· Ἀττικὴ σύνταξις καὶ δυϊκή, ἀντὶ τοῦ ἐν τῷ τῶν Θεσμοφόρων ἱερῷ, Δήμητρος καὶ Περσεφόνης ; Hsch.~δ~2036 δίωγμα· θυσία τις Ἀθήνησιν ἐν ἀπορρήτῳ τελούμενη ὑπὸ τῶν γυναικῶν ἐν τοῖς Θεσμοφορίοις. τὸ αὐτὸ καὶ ἀποδίωγμα ὕστερον ἐκλήθη ; Diogenian.~2.38 Ἀττικοὶ τὰ Ἐλευσίνια: ἐπὶ τῶν καθ᾽ ἑαυτοὺς συλλεγομένων καὶ πραττόντων τι ἀπόῤῥητον. Κρύφα γὰρ κἀκεῖνα ἐτέλουν ; Ar.~\textit{Thesm.}~880 ληροῦντι λῆρον; Θεσμοφορεῖον τουτογί ; Ar.~\textit{Plu.}~411 ἐγώ, κατακλίνειν αὐτὸν εἰς Ἀσκληπιοῦ ; Hsch.~ζ~145 ζημία· θυσία τις ἀποδιδομένη ὑπὲρ τῶν γινομένων ἐν Θεσμοφορίοις ; Hsch.~θ~383 θεσμοφόρια· τὰ μυστήρια τῆς Δήμητρος. ἢ ὄνομα ἑορτῆς ; Hsch.~π~4128 πρυτανεῖον· τρία Ἀθήνησι συσσίτια· πρυτανεῖα, θεσμοφορεῖον, πρυτανεῖον ; Hsch.~κ~3098 κνέωρον· φυτόν τι, ὃ τοῖς Θεσμοφορίοις ὑποστόρνυται, καὶ ὃ εἰς κάθαρσιν χρῶνται. καὶ γυναικεῖον μόριον (com. ad fr. 369 K.–A.)}
\lp{83 Θεσμοφόροιν]}{Ar.~\textit{Thesm.}~286–287 δέσποινα πολυτίμητε Δήμητερ φίλη / καὶ Φερρέφαττα ; Ar.~\textit{Thesm.}~297 εὔχεσθε ταῖν Θεσμοφόροιν ; Hsch.~θ~384 θεσμοφόρος· ὁ θεσμοὺς φέρων, ἤτοι νόμους}
\lp{83 μέλλουσι περί μου τήμερον]}{Ar.~\textit{Thesm.}~181–182 μέλλουσί μ᾽ αἱ γυναῖκες ἀπολεῖν τήμερον / τοῖς Θεσμοφορίοις, ὅτι κακῶς αὐτὰς λέγω ; Ar.~\textit{Thesm.}~377–378 καὶ χρηματίζειν πρῶτα περὶ Εὐριπίδου, / ὅ τι χρὴ παθεῖν ἐκεῖνον}
\lp{83 τήμερον]}{Hsch.~τ~789 τήμερον· Ἀττικοὶ τήμερος καὶ τήμερα λέγουσιν ; Moer. τ~6 Hansen τήμερον Ἀττικοί· σήμερον Ἕλληνες ; Ar.~\textit{Thesm.}~76 τῇδε θἠμέρᾳ κριθήσεται ; Eur.~\textit{Alc.}~20–21 τῇδε γάρ σφ᾽ ἐν ἡμέρᾳ / θανεῖν πέπρωται καὶ μεταστῆναι βίου ; Eur.~\textit{Hipp.}~21–22 ἃ δ᾽ εἰς ἔμ᾽ ἡμάρτηκε, τιμωρήσομαι / Ἱππόλυτον ἐν τῇδ᾽ ἡμέρᾳ ; Soph.~\textit{Aj.}~753–755 εἶρξαι κατ᾽ ἦμαρ τοὐμφανὲς τὸ νῦν τόδε / Αἴανθ᾽ ὑπὸ σκηναῖσι ; Suid.~τ~509 Τήμερος: ὁ σημερινός. καὶ ἔστι τεταγμένον ἐπὶ σώματος. τὸ δὲ τήμερον ἐπὶ χρόνου λέγεται ; [Eur.]~\textit{Rh.}~683 θανῇ γὰρ σήμερον δράσας κακῶς ; Hsch.~σ~503 σήμερον· αὐτίκα, ταχύ (Η 30 . .)}
\lp{84 ἐκκλησιάζειν]}{Ar.~\textit{Thesm.}~302–303 ἐκκλησίαν τήνδε καὶ σύνοδον τὴν νῦν — κάλλιστα καὶ ἄριστα ποῆσαι ; Ar.~\textit{Thesm.}~329–330 τελέως δ᾽ / ἐκκλησιάσαιμεν Ἀθηναίων / εὐγενεῖς γυναῖκες ; Ar.~\textit{Thesm.}~372–373 ἔδοξε τῇ βουλῇ τάδε / τῇ τῶν γυναικῶν ; Suid.~σ~325 ὅτε δὲ ἔμελλε γίνεσθαι ἐκκλησία, σημεῖον ἐτίθετο. οὕτως οὖν καὶ τῶν γυναικῶν μελλουσῶν ἐν τοῖς Θεσμοφορίοις ἐκκλησιάζειν, σημεῖον ἐτέθη ; Diogenian.~4.1 Γυνὴ στρατηγεῖ, καὶ, Γυνὴ στρατεύεται: ἐπὶ τῶν δειλῶν ; Ar.~\textit{Thesm.}~658 τὴν πύκνα πᾶσαν καὶ τὰς σκηνὰς καὶ τὰς διόδους διαθρῆσαι ; Ar.~\textit{Thesm.}~1145 δῆμός τοί σε καλεῖ γυναικῶν ; Is.~8.19 αἵ τε γυναῖκες αἱ τῶν δημοτῶν μετὰ ταῦτα προὔκριναν αὐτὴν μετὰ τῆς Διοκλέους γυναικὸς τοῦ Πιθέως ἄρχειν εἰς τὰ Θεσμοφόρια}
\lp{84 ἐκκλησιάζειν ἐπ᾽ ὀλέθρῳ]}{Eur.~\textit{Or.}~884–887 ἐπεὶ δὲ πλήρης ἐγένετ᾽ Ἀργείων ὄχλος, / κῆρυξ ἀναστὰς εἶπε· Τίς χρῄζει λέγειν, / πότερον Ὀρέστην κατθανεῖν ἢ μὴ χρεών, / μητροκτονοῦντα}
\lp{84 ἐπ᾽ ὀλέθρῳ]}{Diogenian.~4.5 Γυναικῶν ὄλεθροι: ἐπὶ τῶν οἰκτρῶς ἀναιρουμένων ; \textit{idem} Suid.~γ~499 (Γυναικὸς ὄλεθρος ; add. παρόσον πολλοὶ δεινῶς ὑπὸ γυναικῶν ἀπώλοντο) ; Eur.~\textit{IA}~1237 πόθεν ἦλθ᾽ ἐπ᾽ ὀλέθρῳ τὠμῷ, πάτερ ; Soph.~\textit{Ant.}~1291–1292 σφάγιον ἐπ᾽ ὀλέθρῳ / γυναικεῖον ἀμφικεῖσθαι μόρον ; \textit{Il.}~7.402 ὡς ἤδη Τρώεσσιν ὀλέθρου πείρατ᾽ ἐφῆπται ; \textit{idem} Diogenian.~6.96 (om. ὡς ἤδη Τρώεσσιν ; add. ἐπιδέδεται καὶ ηὐτρέπισται τὸ τοῦ ὀλέθρου ἡμῶν πέρας) ; Macar.~4.85 Καθ᾽ αὑτοῦ Βελλεροφόντης: ἐπὶ τῶν ἐπ᾽ ὀλέθρῳ οἰκείῳ διακονούντων ; Eur.~\textit{IA}~886 ὦ θύγατερ, ἥκεις ἐπ᾽ ὀλέθρῳ καὶ σὺ καὶ μήτηρ σέθεν ; Eur.~\textit{IT}~601 οὔκουν δίκαιον ἐπ᾽ ὀλέθρῳ τῷ τοῦδ᾽ ἐμὲ / χάριν τίθεσθαι ; Eur.~\textit{Phoen.}~533–534 πολλοὺς δ᾽ ἐς οἴκους καὶ πόλεις εὐδαίμονας / ἐσῆλθε κἀξῆλθ᾽ ἐπ᾽ ὀλέθρῳ τῶν χρωμένων ; Lys.~12.60 μισθωσάμενοι δὲ πάντας ἀνθρώπους ἐπ᾽ ὀλέθρῳ τῆς πόλεως ; Hsch.~ε~6280 ἐς ὄλεθρον· εἰς Ἅιδην ; Hsch.~κ~2677 Κιλίκιος ὄλεθρος· πονηρός. πονηροὶ γὰρ οἱ Κίλικες ; Hsch.~ο~516 ὄλεθρος βαθύς· ἔνιοι ⟨παρὰ Ἀριστοφάνει (fr. 332, 3 K.–A.)⟩ ἐπὶ κόσμου γυναικείου· ⟨ἄλλοι ἐπὶ δαπάνης … ὄλεθρος γὰρ ὁ θάνατος⟩}
\lp{84 τιὴ τί δή]}{Ar.~\textit{Vesp.}~1155 παράθου γε μέντοι καὶ κρεάγραν. — τιὴ τί δή ; Ar.~\textit{Pax}~1018 ἀλλ᾽ οὐ θέμις. — τιὴ τί δή ; Hsch.~τ~856 τίη; διατί; ἱνατί ; \textit{idem} Suid.~τ~567 (δή pro ἱνατί) ; Eust.~631 (ad \textit{Il.}~6.145) Τὸ δὲ τιὴ ὡς Ἀττικῶς ἔχει τὴν διὰ τοῦ η ἐπέκτασιν, καθὰ καὶ τὸ ἐπειή … Ὅτι δὲ τὸ τίη, ὃ κεῖται καὶ παρὰ τῷ κωμικῷ, παροξύνουσιν οἱ παλαιοὶ τεχνικοί ; Ar.~\textit{Eq.}~126 τὸν περὶ σεαυτοῦ χρησμὸν ὀρρωδῶν; — τιή ; Ar.~\textit{Eq.}~731 ὑπὸ τουτουὶ καὶ τῶν νεανίσκων. — τιή ; Ar.~\textit{Nub.}~784 ὁτιὴ τί; ναὶ πρὸς τῶν θεῶν ὦ Σώκρατες ; Ar.~\textit{Pax}~927 ἀλλ᾽ ὑὶ παχείᾳ καὶ μεγάλῃ; — μὴ μή. — τιή ; Ar.~\textit{Ran.}~1473 ἔκρινα νικᾶν Αἰσχύλον. τιὴ γὰρ οὔ ; Ar.~\textit{Eccl.}~796 θάρρει, καταθήσεις, κἂν ἔνης ἔλθῃς. — τιή ; Ar.~\textit{Eccl.}~1086 χαλεπαί γ᾽ ἂν ἦστε γενόμεναι πορθμῆς. — τιή ; Ar.~\textit{Plu.}~136 ὅτι τί δή}
\lp{85 ὁτιή]}{Suid.~ο~733 Ὁτίη τί: ἀντὶ τοῦ διατί. τὸ δὲ τί περιττόν. Ἀριστοφάνης Νεφέλαις ; Ar.~\textit{Nub.}~755 οὐκ ἂν ἀποδοίην τοὺς τόκους. — ὁτιὴ τί δή}
\lp{85 τραγῳδῶ]}{Ar.~\textit{Pax}~147–148 εἶτα χωλὸς ὢν Εὐριπίδῃ / λόγον παράσχῃς καὶ τραγῳδία γένῃ ; Ar.~\textit{Thesm.}~149–150 χρὴ γὰρ ποιητὴν ἄνδρα πρὸς τὰ δράματα / ἃ δεῖ ποιεῖν πρὸς ταῦτα τοὺς τρόπους ἔχειν ; Ar.~\textit{Thesm.}~450–451 νῦν δ᾽ οὗτος ἐν ταῖσιν τραγῳδίαις ποιῶν / τοὺς ἄνδρας ἀναπέπεικεν οὐκ εἶναι θεούς ; Hsch.~δ~1105 διακωμῳδῶ· διατραγῳδῶ ; Ar.~\textit{Nub.}~1091 τί δαί; τραγῳδοῦσ᾽ ἐκ τίνων ; Ar.~\textit{Ran.}~368 κωμῳδηθεὶς ἐν ταῖς πατρίοις τελεταῖς ταῖς τοῦ Διονύσου ; D.~19.189 ποῦ δ᾽ ἅλες; ποῦ τράπεζα; ποῦ σπονδαί; ταῦτα γὰρ τραγῳδεῖ περιιών ; Hsch.~κ~4843 κωμῳδούμενοι· παιζόμενοι παρὰ τὴν τραγῳδίαν ; Hsch.~τ~1238 τραγῳδεῖν· χορεύειν ; \textit{idem} Phot.~τ~414 (add. Ἀριστοφάνης) ; Phot.~τ~413 τραγῳδία· δεινοπάθεια ; Diogenian.~5.13 Ἢ χρὴ τραγῳδεῖν πάντας, ἢ μελαγχολᾶν: ἐπὶ τῶν μήτε λυπεῖσθαι μήτε χαίρειν δυναμένων}
\lp{85 τραγῳδῶ καὶ κακῶς αὐτὰς λέγω]}{Ar.~\textit{Ach.}~630–631 διαβαλλόμενος δ᾽ ὑπὸ τῶν ἐχθρῶν ἐν Ἀθηναίοις ταχυβούλοις, / ὡς κωμῳδεῖ τὴν πόλιν ἡμῶν καὶ τὸν δῆμον καθυβρίζει ; Ar.~\textit{Pax}~751 οὐκ ἰδιώτας ἀνθρωπίσκους κωμῳδῶν οὐδὲ γυναῖκας ; Diph.~fr.~74.3–5~K.-A. οὐκ ἄν ποτε / Εὐριπίδης γυναῖκα σώσει᾽. οὐχ ὁρᾷς / ἐν ταῖς τραγῳδίαισιν αὐτὰς ὡς στυγεῖ ; [X.]~\textit{Ath.}~2.18 κωμῳδεῖν δ᾽ αὖ καὶ κακῶς λέγειν τὸν μὲν δῆμον οὐκ ἐῶσιν, ἵνα μὴ αὐτοὶ ἀκούωσι κακῶς}
\lp{85 κακῶς αὐτὰς λέγω]}{Ar.~\textit{Lys.}~368–369 οὐκ ἔστ᾽ ἀνὴρ Εὐριπίδου σοφώτερος ποιητής· / οὐδὲν γὰρ οὕτω θρέμμ᾽ ἀναιδές ἐστιν ὡς γυναῖκες ; Ar.~\textit{Thesm.}~181–182 μέλλουσί μ᾽ αἱ γυναῖκες ἀπολεῖν τήμερον / τοῖς Θεσμοφορίοις, ὅτι κακῶς αὐτὰς λέγω ; Ar.~\textit{Thesm.}~386–392 προπηλακιζομένας … ὑπὸ / Εὐριπίδου τοῦ τῆς λαχανοπωλητρίας / καὶ πολλὰ καὶ παντοῖ᾽ ἀκουούσας κακά. … / ποῦ δ᾽ οὐχὶ διαβέβληχ᾽, ὅπουπερ ἔμβραχυ / εἰσὶν θεαταὶ καὶ τραγῳδοὶ καὶ χοροί, / τὰς μοιχοτρόπους, τὰς ἀνδρεραστρίας καλῶν ; Ar.~\textit{Thesm.}~546–548 ἐπίτηδες εὑρίσκων λόγους, ὅπου γυνὴ πονηρὰ / ἐγένετο, Μελανίππας ποιῶν Φαίδρας τε· Πηνελόπην δὲ / οὐπώποτ᾽ ἐποίησ᾽, ὅτι γυνὴ σώφρων ἔδοξεν εἶναι ; Ar.~\textit{Thesm.}~1160–1167 ἐφ᾽ ᾧτ᾽ ἀκοῦσαι μηδὲν ὑπ᾽ ἐμοῦ μηδαμὰ / κακὸν τὸ λοιπόν … οὐδὲν μή ποτε / κακῶς ἀκούσητ᾽ ; Eur.~\textit{Hipp.}~616–617 ὦ Ζεῦ, τί δὴ κίβδηλον ἀνθρώποις κακὸν / γυναῖκας ἐς φῶς ἡλίου κατῴκισας ; Eur.~\textit{Hipp.}~664–666 ὄλοισθε. μισῶν δ᾽ οὔποτ᾽ ἐμπλησθήσομαι / γυναῖκας, οὐδ᾽ εἴ φησί τίς μ᾽ ἀεὶ λέγειν· / ἀεὶ γὰρ οὖν πώς εἰσι κἀκεῖναι κακαί ; Eur.~fr.~494.1–2~Kannicht μάτην ἄρ᾽ ἐς γυναῖκας ἐξ ἀνδρῶν ψόγος / ψάλλει κενὸν τόξευμα καὶ κακῶς λέγει· / αἱ δ᾽ εἴσ᾽ ἀμείνους ἀρσένων, ἐγὼ λέγω ; Satyr.~\textit{Vit.\,Eur.} (P.Oxy.~1176) fr.~39 col. x~1–10 ἀπήχθοντ᾽ αὐτῷ πάντες, οἱ μὲν ἄνδρες διὰ τὴν δυ[σ]ομιλίαν, α[ἱ δ]ὲ γυναῖκε[ς δ]ιὰ τοὺς ψόγους τοὺς ἐν τοῖς ποιήμασιν ; Antip. Tars.~\textit{SVF} III fr.~63 (Stob.~4.22) καὶ ὁ Εὐριπίδης εἰς ταῦτα ἀποβλέψας καὶ ἀποθέμενος τὴν ἐν τῷ γράφειν μισογυνίαν ; Ath.~13.557e εἰπόντος Σοφοκλεῖ τινος ὅτι μισογύνης ἐστὶν Εὐριπίδης, ἔν γε ταῖς τραγῳδίαις, ἔφη ὁ Σοφοκλῆς· ἐπεὶ ἔν γε τῇ κλίνῃ φιλογύνης ; Vit. Eur.~\textit{§}~6 (Σ Eur. I p.~6 Schwartz) λέγουσι δὲ καὶ ὅτι αἱ γυναῖκες διὰ τοὺς ψόγους, οὓς ἐποίει εἰς αὐτὰς διὰ τῶν ποιημάτων, τοῖς Θεσμοφορίοις ἐπέστησαν αὐτῷ βουλόμεναι ἀνελεῖν. ἐφείσαντο δὲ αὐτοῦ πρῶτον μὲν διὰ τὰς Μούσας, ἔπειτα δὲ βεβαιωσαμένου μηκέτι αὐτὰς κακῶς ἐρεῖν ; Gell.~15.20.6–7 mulieres fere omnes in maiorem modum exosus fuisse dicitur … eius odii in mulieres Aristophanes quoque meminit ἐν ταῖς προτέραις Θεσμοφοριαζούσαις ; Ar.~\textit{Ach.}~502–503 οὐ γάρ με νῦν γε διαβαλεῖ Κλέων ὅτι / ξένων παρόντων τὴν πόλιν κακῶς λέγω ; Ar.~\textit{Lys.}~283 τασδὶ δὲ τὰς Εὐριπίδῃ θεοῖς τε πᾶσιν ἐχθρὰς ; Ar.~\textit{Lys.}~475 ἢν μή τις ὥσπερ σφηκιὰν βλίττῃ με κἀρεθίζῃ ; Ar.~\textit{Thesm.}~539 γυνὴ γυναῖκας οὖσα μὴ κακῶς λέγειν τὸ λοιπόν ; Ar.~\textit{Thesm.}~785–787 ἡμεῖς τοίνυν ἡμᾶς αὐτὰς εὖ λέξωμεν παραβᾶσαι, / καίτοι πᾶς τις τὸ γυναικεῖον φῦλον κακὰ πόλλ᾽ ἀγορεύει, / ὡς πᾶν ἐσμὲν κακὸν ἀνθρώποις κἀξ ἡμῶν ἐστιν ἅπαντα ; Ar.~\textit{Thesm.}~962–964 εἰ δέ τις / προσδοκᾷ κακῶς ἐρεῖν / ἐν ἱερῷ γυναῖκά μ᾽ οὖσαν ἄνδρας, οὐκ ὀρθῶς φρονεῖ ; Ar.~\textit{Ran.}~1043–1044 ἀλλ᾽ οὐ μὰ Δί᾽ οὐ Φαίδρας ἐποίουν πόρνας οὐδὲ Σθενεβοίας, / οὐδ᾽ οἶδ᾽ οὐδεὶς ἥντιν᾽ ἐρῶσαν πώποτ᾽ ἐποίησα γυναῖκα ; Hsch.~α~2014 αἰξωνεύεσθαι· κακηγορεῖν (cf. Men. fr. 200 K.–A.) ; Hsch.~α~7434 ἀρρηνεῖν· λοιδορεῖν. καὶ ⟨ἐπὶ⟩ γυναικὶ· πρὸς ἄνδρα διαφέρεσθαι ; Hsch.~κ~301 κακήγορος· κακῶς λέγων, κακολόγος}
\lp{86 νὴ τὸν Ποσειδῶ]}{Ar.~\textit{Ran.}~276 νὴ τὸν Ποσειδῶ ᾽γωγε, καὶ νυνί γ᾽ ὁρῶ ; Ar.~\textit{Ach.}~798 ναὶ τὸν Ποτειδᾶν καί κ᾽ ἄνις γα τῶ πατρός ; Ar.~\textit{Eq.}~1201 νὴ τὸν Ποσειδῶ· καὶ σὺ γὰρ τοὺς ἐκ Πύλου ; Ar.~\textit{Ran.}~183 νὴ τὸν Ποσειδῶ κἄστι γ᾽ ὁ Χάρων οὑτοσί ; Ar.~\textit{Ran.}~1074 νὴ τὸν Ἀπόλλω, καὶ προσπαρδεῖν γ᾽ ἐς τὸ στόμα τῷ θαλάμακι}
\lp{86 νὴ τὸν Ποσειδῶ καὶ δικαίαν ἂν πάθοις]}{Ar.~\textit{Ach.}~560–561 νὴ τὸν Ποσειδῶ, καὶ λέγει γ᾽ ἅπερ λέγει / δίκαια πάντα κοὐδὲν αὐτῶν ψεύδεται ; Ar.~\textit{Lys.}~403 νὴ τὸν Ποσειδῶ τὸν ἁλυκὸν δίκαιά γε}
\lp{86 καὶ δικαίαν ἂν πάθοις]}{Ar.~\textit{Nub.}~1379 ἀλλ᾽ αὖθις αὖ τυπτήσομαι. — νὴ τὸν Δί᾽ ἐν δίκῃ γ᾽ ἄν ; Aesch.~\textit{Cho.}~313–314 δράσαντι παθεῖν, / τριγέρων μῦθος τάδε φωνεῖ ; Arist.~\textit{EN}~1132b25–27 εἴ κε πάθοι τά τ᾽ ἔρεξε, δίκη κ᾽ ἰθεῖα γένοιτο ; Ar.~\textit{Lys.}~68 ὁ γοῦν ἀνάγυρός μοι κεκινῆσθαι δοκεῖ ; Luc.~\textit{Herm.}~58 Οἶμαι ἔγωγε καὶ δίκαιά γ᾽ ἂν πάθοιμι ; Aesop.~\textit{Fab.}~82.1 Chambry θνῄσκων δὲ ἔλεγε· Δίκαια πάσχω, τὸν πονηρὸν οἰκτείρας ; Diogenian.~1.25 Ἀνάγυρον κινεῖς: ἐπὶ τῶν κινούντων τινὰ ἐπὶ τῷ ἑαυτῶν κακῷ ; \textit{idem} Zen.~2.55 (Εἴρηται οὖν ἡ παροιμία ἐντεῦθεν pro Ἀνάγυρον κινεῖς ; ἐπὶ κακῷ ἑαυτῶν) ; Ar.~\textit{Ach.}~323 οὐκ ἀκουσόμεσθα δῆτα. — δεινά τἄρα πείσομαι ; Pl.~\textit{Ap.}~41e–42a καὶ ἐὰν ταῦτα ποιῆτε, δίκαια πεπονθὼς ἐγὼ ἔσομαι ὑφ᾽ ὑμῶν αὐτός τε καὶ οἱ ὑεῖς ; X.~\textit{Hell.}~7.3.12 ἔγνωσαν δίκαια τὸν Εὔφρονα πεπονθέναι ; Plu.~\textit{Lyc.}~15.2 οἱ δὲ περιϊόντες ᾖδον εἰς αὑτοὺς ᾠδήν τινα πεποιημένην, ὡς δίκαια πάσχοιεν, ὅτι τοῖς νόμοις ἀπειθοῦσι ; Σ Eur.~\textit{Or.}~538 ἔπαθε μὲν δίκαια, καθὸ ἀπέθανε, καὶ δικαίαν τιμωρίαν δέδωκε}
\lp{86 καὶ δικαίαν ἂν πάθοις (ellipse de πληγήν)]}{Ar.~\textit{Nub.}~972 ἐπετρίβετο τυπτόμενος πολλὰς ὡς τὰς Μούσας ἀφανίζων ; Ar.~\textit{Av.}~1018 φθαίης ἄν· ἐπίκεινται γὰρ ἐγγὺς αὑταιί ; Aesch.~\textit{Ag.}~1385–1386 καὶ πεπτωκότι / τρίτην ἐπενδίδωμι ; Soph.~\textit{El.}~1416 παῖσον, εἰ σθένεις, διπλῆν ; Soph.~\textit{OC}~544 δευτέραν ἔπαισας, ἐπὶ νόσῳ νόσον ; Theoc.~14.34–35 πὺξ ἐπὶ κόρρας / ἤλασα, κἄλλαν αὖθις ; Hdt.~3.64.3 ὥς οἱ καιρίῃ ἔδοξε τετύφθαι ; D.~19.197 ξαίνει κατὰ τοῦ νώτου πολλάς ; Ter.~\textit{Haut.}~402 diu etiam duras dabit}
\lp{86 καὶ δικαίαν ἂν πάθοις (γνῶμαι d’Euripide)]}{Ar.~\textit{Thesm.}~275–276 μέμνησο τοίνυν ταῦθ᾽, ὅτι ἡ φρὴν ὤμοσεν, / ἡ γλῶττα δ᾽ οὐκ ὀμώμοκ᾽ ; Eur.~\textit{Alc.}~691 χαίρεις ὁρῶν φῶς· πατέρα δ᾽ οὐ χαίρειν δοκεῖς ; \textit{idem} Ar.~\textit{Thesm.}~194 ; Eur.~\textit{Hipp.}~612 ἡ γλῶσσ᾽ ὀμώμοχ᾽, ἡ δὲ φρὴν ἀνώμοτος ; \textit{idem} Arist.~\textit{Rh.}~3.15.8 (add. ὥσπερ Εὐριπίδης πρὸς Ὑγιαίνοντα ἐν τῇ ἀντιδόσει κατηγοροῦντα ὡς ἀσεβής, ὅς γ᾽ ἐποίησε κελεύων ἐπιορκεῖν)}
\lp{86 δικαίαν]}{EM~24.48 Ἀθήνη Ἀθηναία, ὡς δίκη δικαία ; Hsch.~δ~1792 δίκαια· ἔννομα}
\lp{87 ἀτάρ]}{Ar.~\textit{Thesm.}~207 ἀτὰρ ἡ πρόφασίς γε νὴ Δί᾽ εἰκότως ἔχει ; Hsch.~α~8011 ἀτὰρ δέ· πλήν, ὅμως δέ}
\lp{87 ἀτὰρ τίν᾽]}{Eur.~\textit{Hec.}~258 ἀτὰρ τί δὴ σόφισμα τοῦθ᾽ ἡγούμενοι ; Ar.~\textit{Thesm.}~215 ἀτὰρ τί μέλλεις δρᾶν μ᾽ ; Eur.~\textit{Heracl.}~879 ἀτὰρ τί κεύθων Ἰόλεως σοφόν ποτε ; Ar.~\textit{Nub.}~30 ἀτὰρ τί χρέος ἔβα με μετὰ τὸν Πασίαν ; Ar.~\textit{Eccl.}~358 ἀτὰρ τί δράσω}
\lp{87 τίν᾽ … μηχανὴν ἔχεις]}{Ar.~\textit{Thesm.}~765–766 ἄγε δὴ τίς ἔσται μηχανὴ σωτηρίας; / τίς πεῖρα, τίς ἐπίνοι᾽ ; Eur.~\textit{Hel.}~1033–1034 τοὐνθένδε δὴ σὲ τοὺς λόγους φέροντα χρὴ / κοινὴν ξυνάπτειν μηχανὴν σωτηρίας ; Pl.~\textit{R.}~415c τοῦτον οὖν τὸν μῦθον ὅπως ἂν πεισθεῖεν, ἔχεις τινὰ μηχανήν ; Dio Chrys.~59.10 εἰ δ᾽ οὖν ἔχεις τινὰ μηχανήν ; Hdt.~3.85.1 νῦν ὦν εἴ τινα ἔχεις σοφίην, μηχανῶ ; Hdt.~3.85.2 εἰ τοίνυν τι τοιοῦτον ἔχεις σόφισμα, ὥρη μηχανᾶσθαι}
\lp{87 ἐκ ταύτης]}{Aesch.~\textit{Pers.}~788 πῶς ἂν ἐκ τούτων ἔτι / πράσσοιμεν ὡς ἄριστα Περσικὸς λεώς ; Ar.~\textit{Ran.}~185 τίς εἰς ἀναπαύλας ἐκ κακῶν καὶ πραγμάτων ; Antiatt.~94.7~Bekker (Cratin. fr.~37~K.-A.) Ἐξ αὐτῆς: ἀντὶ τοῦ παραχρῆμα. Κρατῖνος Δηλιάσιν}
\lp{87 ἐκ ταύτης … μηχανήν]}{Aesch.~\textit{Supp.}~457–459 ἔχω στρόφους ζώνας τε, συλλαβὰς πέπλων … ἐκ τῶνδε τοίνυν, ἴσθι, μηχανὴ καλή ; Ar.~\textit{Vesp.}~346 ἀλλ᾽ ἐκ τούτων ὥρα τινά σοι ζητεῖν καινὴν ἐπίνοιαν ; [Arist.]~\textit{Mu.}~396b33–34 καὶ ἐκ τούτων μηχανησαμένη τῷ παντὶ σωτηρίαν ; Hld.~9.6 εἰ μὴ φθαίη τις ἐκ θεῶν καὶ ἐξ Ὑδάσπου μηχανὴ σωτηρίας}
\lp{87 μηχανήν]}{Ar.~\textit{Ach.}~738 ἀλλ᾽ ἔστι γάρ μοι Μεγαρικά τις μαχανά ; Ar.~\textit{Thesm.}~927 ἢν μὴ προλίπωσ᾽ αἱ μυρίαι με μηχαναί ; Ar.~\textit{Thesm.}~928 αὕτη μὲν ἡ μήρινθος οὐδὲν ἔσπασεν ; \textit{idem} Diogenian.~3.35 (ἔσπακεν ; add. ἐπὶ τῶν μηδὲν ἑλκυσάντων), Eust.~760.45 (ad \textit{Il.}~9.423) (add. ὕστερον δὲ παροιμία ὁμοία τοιαύτη ἐξεφωνήθη et ἀντὶ τοῦ αὕτη ἡ μηχανὴ οὐδὲν ἤνυσεν) ; Antiph.~fr.~189.13–16~K.-A. ⟨ἔπει⟩θ᾽ ὅταν μηθὲν δύνωντ᾽ εἰπεῖν ἔτι, / κομιδῆι δ᾽ ἀπειρήκωσιν ἐν τοῖς δράμασιν, / αἴρουσιν ὥσπερ δάκτυλον τὴν μηχανήν, / καὶ τοῖς θεωμένοισιν ἀποχρώντως ἔχει ; Eur.~\textit{Andr.}~66 τί δρῶσι; ποίας μηχανὰς πλέκουσιν αὖ ; Aesch.~\textit{PV}~59 δεινὸς γὰρ εὑρεῖν κἀξ ἀμηχάνων πόρον ; Aesch.~\textit{PV}~469–471 τοιαῦτα μηχανήματ᾽ ἐξευρὼν τάλας / βροτοῖσιν, αὐτὸς οὐκ ἔχω σόφισμ᾽ ὅτῳ / τῆς νῦν παρούσης πημονῆς ἀπαλλαγῶ ; Aesch.~fr.~373~Radt δεινοὶ πλέκειν τοι μηχανὰς Αἰγύπτιοι ; \textit{idem} Zen.~3.37 (add. ἐπὶ τῶν σφόδρα κακουργοτάτων· τοιοῦτοι γὰρ οἱ Αἰγύπτιοι) ; Pl.~\textit{Cra.}~425d ὥσπερ οἱ τραγῳδοποιοὶ ἐπειδάν τι ἀπορῶσιν ἐπὶ τὰς μηχανὰς καταφεύγουσι θεοὺς αἴροντες ; Arist.~\textit{Po.}~1454a37–b1 φανερὸν οὖν ὅτι καὶ τὰς λύσεις τῶν μύθων ἐξ αὐτοῦ δεῖ τοῦ μύθου συμβαίνειν, καὶ μὴ ὥσπερ ἐν τῇ Μηδείᾳ ἀπὸ μηχανῆς ; Poll.~4.128–129 ἡ μηχανὴ δὲ θεοὺς δείκνυσι καὶ ἥρως τοὺς ἐν ἀέρι Βελλεροφόντας ἢ Περσέας, [καὶ κεῖται κατὰ τὴν ἀριστερὰν πάροδον, ὑπὲρ τὴν σκηνὴν τὸ ὕψος]. ὃ δ᾽ ἐστὶν ἐν τραγῳδίᾳ μηχανή, [τοῦτο] [καλοῦσιν] ἐν κωμῳδίᾳ κράδην ; AB~208.9–12 Ἀπὸ μηχανῆς: μηχανή ἐστι παρὰ τοῖς κωμικοῖς ἐγκυκλήματός τι εἶδος ἀπὸ συνθήκης πρὸς ὃ φέρεται ὁ εἰς τὴν σκηνὴν δείξεως χάριν θεοῦ ἢ ἄλλου τινὸς ἥρωος ; Suid.~α~1081 Ἀλλ᾽ ἔστιν ἡμῖν Μεγαρική τις μηχανή: ἀντὶ τοῦ πονηρά, πανοῦργος μηχανή. διεβάλλοντο γὰρ ἐπὶ πονηρίᾳ οἱ Μεγαρεῖς ; Zen.~1.93 Ἂν ἡ λεοντῆ μὴ ἐξίκηται, τὴν ἀλωπεκῆν πρόσαψον: ἂν μὴ φανερῶς δύνῃ βλάψαι, πανουργίᾳ χρῆσαι. Ἂν μὴ κατὰ ῥώμην τὸ προκείμενον ἐξανύοιτο, μηχανῇ καὶ τέχνῃ περαινέσθω τὸ λειπόμενον ; \textit{idem} Tosi n°~2223 (ὅπου γὰρ ἡ λεοντῆ μὴ ἐφικνεῖται, προσραπτέον ἐκεῖ τὴν ἀλωπεκῆν), Erasm.~\textit{Adag.}~3.5.81 (Si leonina pellis non satis est, vulpina addenda … Cum viribus res obtineri non potest, dolus adhibendus) ; Diogenian.~2.84 Ἀπὸ μηχανῆς θεὸς ἐπιφανείς: ἐπὶ τῶν ἀπροςδοκήτως ἐπ᾽ ὠφελείᾳ τινὶ ἀναφαινομένων. Ἐν γὰρ ταῖς τραγῳδίαις θεοὶ ἐξ ἀφανοῦς ἐπὶ τῆς σκηνῆς ἐφαίνοντο ; \textit{idem} Men.~fr.~213~K.-A. (Σ \textit{Pl.\,Clit.}~407a) (παροιμία ἀπὸ μηχανῆς θεὸς ἐπεφάνης ; ἐπ᾽ ὠφελείαι ἢ σωτηρίαι φαινομένων ; ἐν γὰρ ταῖς τραγωιδίαις ἐξ ἀφανοῦς θεοὶ ; add. Μένανδρος Θεοφορουμένηι καὶ Κεκρυφάλωι), Tosi n°~1203 (Deus ex machina), Erasm.~\textit{Adag.}~1.1.68 (Θεὸς ἀπὸ μηχανῆς ἐπιφανείς, id est Deus ex improviso apparens. In eos dicitur, quibus in rebus perplexis praeter spem exoritur aliquis, qui salutem adferat) ; Ar.~\textit{Ach.}~391 ἀλλ᾽ ἐξάνοιγε μηχανὰς τὰς Σισύφου ; Ar.~\textit{Ach.}~445 δώσω· πυκνῇ γὰρ λεπτὰ μηχανᾷ φρενί ; Ar.~\textit{Eq.}~756 νῦν δή σε πάντα δεῖ κάλων ἐξιέναι σεαυτοῦ ; Ar.~\textit{Eq.}~759 κἀκ τῶν ἀμηχάνων πόρους εὐμήχανος πορίζειν ; Ar.~\textit{Nub.}~479–480 ἵν᾽ αὐτὸν εἰδὼς ὅστις ἐστὶ μηχανὰς / ἤδη ᾽πὶ τούτοις πρὸς σὲ καινὰς προσφέρω ; Ar.~\textit{Vesp.}~149 ἐνταῦθά νυν ζήτει τιν᾽ ἄλλην μηχανήν ; Ar.~\textit{Pax}~174 ὦ μηχανοποιὲ πρόσεχε τὸν νοῦν ὡς ἐμέ ; Ar.~\textit{Thesm.}~1132 ἀλλ᾽ ἄλλην τινὰ / τούτῳ πρέπουσαν μηχανὴν προσοιστέον ; Ar.~\textit{Ran.}~1429 καὶ πόριμον αὑτῷ τῇ πόλει δ᾽ ἀμήχανον ; Ar.~fr.~192~K.-A. (Daedalus) ὁ μηχανοποιὸς ὁπότε βούλει τὸν τροχὸν / ἐᾶν †κἀνεκάς† λέγε «χαῖρε φέγγος ἡλίου» ; Cratin.~fr.~320~K.-A. τούτου δ᾽ ἁπάσας ἀποτεμῶ τὰς μηχανάς ; Alex.~fr.~131.9~K.-A. ἀπὸ μηχανῆς πωλοῦντες ὥσπερ οἱ θεοί ; Eur.~\textit{Andr.}~85 πολλὰς ἂν εὕροις μηχανάς· γυνὴ γὰρ εἶ ; Eur.~\textit{Hel.}~813 ἐς ἄπορον ἥκεις· δεῖ δὲ μηχανῆς τινος ; Eur.~\textit{Heracl.}~1002 ἔχθραν πατρῴαν, πάντα κινῆσαι πέτρον ; \textit{idem} Hsch.~π~389 (πάντα κινήσας πέτρον· παροιμία) ; Eur.~\textit{Hipp.}~480–481 ἦ τἄρ᾽ ἂν ὀψέ γ᾽ ἄνδρες ἐξεύροιεν ἄν, / εἰ μὴ γυναῖκες μηχανὰς εὑρήσομεν ; Eur.~\textit{IA}~413 ἐγὼ δ᾽ ἐπ᾽ ἄλλας εἶμι μηχανάς τινας ; Eur.~\textit{Med.}~260–261 ἤν μοι πόρος τις μηχανή τ᾽ ἐξευρεθῇ / πόσιν δίκην τῶνδ᾽ ἀντιτείσασθαι κακῶν ; Eur.~\textit{Med.}~278 ἐχθροὶ γὰρ ἐξιᾶσι πάντα δὴ κάλων ; Eur.~\textit{Phoen.}~890 μί᾽ ἔστιν ἄλλη μηχανὴ σωτηρίας ; \textit{Il.}~9.423 ὄφρ᾽ ἄλλην φράζωνται ἐνὶ φρεσὶ μῆτιν ἀμείνω ; Hdt.~3.152 καίτοι πάντα σοφίσματα καὶ πάσας μηχανὰς ἐπεποιήκεε ἐς αὐτοὺς Δαρεῖος ; [Pl.]~\textit{Clit.}~407a ὥσπερ ἐπὶ μηχανῆς τραγικῆς θεός ; Plu.~\textit{Lys.}~25.1–2 ὥσπερ ἐν τραγῳδίᾳ μηχανὴν αἴρων ἐπὶ τοὺς πολίτας, λόγια πυθόχρηστα καὶ χρησμοὺς συνετίθει καὶ κατεσκεύαζεν ; Luc.~\textit{Hermot.}~86 ἀνέσπασας ἐπιστάς, τὸ τῶν τραγῳδῶν τοῦτο, θεὸς ἐκ μηχανῆς ἐπιφανείς ; Luc.~\textit{Philops.}~29 καὶ τὸ τοῦ λόγου, θεὸν ἀπὸ μηχανῆς ἐπεισκυκληθῆναί μοι τοῦτον ᾤμην ὑπὸ τῆς Τύχης ; Alciphr.~3.4 τίς δαίμων ἢ θεὸς ἀπὸ μηχανῆς ἐρρύσατό μ᾽ ἐν ἀκαρεῖ ; Hsch.~α~3649 ἀμήχανον· ἄπορον … πρὸς ὃν μηχανὴν οὐκ ἔστιν … εὑρεῖν, ἢ δυσχερές, ἢ ἐνδεές ; Hsch.~μ~1301 μηχαναί· τέχναι. ἐπίνοιαι. βουλεύματα. ἅμαξαι. συμπεράσματα ; Hsch.~μ~1303 μηχανορράφος· κατασκευαστής, ἐπινοητὴς κακῶν ; \textit{idem} Phot.~μ~433 (ἐπινοητὴς κακῶν, κατασκευαστής) ; Phot.~μ~432 μηχανικόν· Ξενοφῶν (HG. 3, 1, 8) τὸν πανοῦργον καὶ δεινόν ; Suid.~α~3438 Ἀπὸ μηχανῆς: ἐπὶ τῶν παραδόξων καὶ παραλόγων. οἱ γὰρ τῶν τραγῳδιῶν ποιηταὶ … εἰώθασι θεοὺς εἰσάγειν, οὐκ ἐπ᾽ αὐτῆς τῆς σκηνῆς ὁρμωμένους, ἀλλ᾽ ἐξ ὕψους ὑπό τινος μηχανῆς ; Suid.~α~3504 λέγεται δὲ ἄπορος καὶ ὁ μὴ δυνάμενος μηχανήν τινα εὑρεῖν, ὁ ἐνδεής ; Zen.~3.77 (Ἐν ἀμηχανίῃ) Ἐν γὰρ ἀμηχανίῃ καὶ Καρκίνος ἔμμορε τιμῆς ; Zen.~5.62 Πάντα κάλων σεῖε: παροιμία ἐπὶ τῶν πάσῃ προθυμίᾳ χρωμένων ; \textit{idem} Hsch.~π~388 (σείει), Tosi n°~1285 (Plenis velis) (σείειν ; om. παροιμία … χρωμένων) ; Zen.~5.63 Πάντα λίθον κίνει: παροιμία ; \textit{idem} Hsch.~π~390 (κινεῖν ; add. ἧς μέμνηται Ἀρίσταρχος) ; Plu.~\textit{Prov.Alex.}~2.16 Κράδης ῥαγείσης: κράδη νῦν οὐχ ὁ σύκινος κλάδος, ἀλλ᾽ ἡ ἀγκυρὶς, ἀφ᾽ ἧς οἱ ὑποκριταὶ ἐν ταῖς τραγικαῖς σκηναῖς ἐξαρτῶνται, θεοῦ μιμούμενοι ἐπιφάνειαν, ζωστῆρσι καὶ ταινίαις κατειλημμένοι. ἐπὶ τῶν προσφανέντων αἰφνιδίως καὶ ἀσχημονούντων ; App. Prov.~3.21 Θρᾳκεία παρεύρεσις: ἐπὶ τῶν δολίως τὶ μηχανωμένων καὶ χαλεποῖς ἔργοις λόγους εὐπροςώπους ἐξευρισκόντων ; Phot.~α~646 αἶρε δάκτυλον· τίθεται ἐπὶ τῶν ἀπαγορευόντων ἐν τῇ ἀγωνίᾳ· ἐπαίρεσθαι γὰρ τὴν χεῖρα σύμβολον τοῦ νενικῆσθαι ; \textit{idem} Synag.~α~564~Cunningham, Suid.~αι~281 (add. παροιμία et καὶ Πολύβιος· ὁ δὲ προσπεσὼν αἴρεται νίκην εὐτυχεστάτην), Apostol.~1.75 (ἐπὶ τῶν νενικηκότων. ἐπαίρουσι γὰρ τὴν χεῖρα …) ; Ar.~\textit{Vesp.}~365–366 ἀλλὰ καὶ νῦν ἐκπόριζε / μηχανὴν ὅπως τάχισθ᾽ ; Ar.~\textit{Lys.}~111 ἐθέλοιτ᾽ ἂν οὖν, εἰ μηχανὴν εὕροιμ᾽ ἐγώ ; Eur.~\textit{Alc.}~221 ἔξευρε μηχανάν τιν᾽ Ἀδμήτῳ κακῶν ; Eur.~\textit{IT}~112 ἄγαλμα πάσας προσφέροντε μηχανάς ; Eur.~\textit{Ion}~833 οἳ συντιθέντες τἄδικ᾽ εἶτα μηχαναῖς ; [Eur.]~\textit{Rh.}~141 κἂν μέν τιν᾽ ἐχθρῶν μηχανὴν πυθώμεθα ; [Eur.]~\textit{Rh.}~578 ἴσως ἐφ᾽ ἡμῖν μηχανὴν στήσων τινά ; Aesch.~\textit{Ag.}~1609 πᾶσαν συνάψας μηχανὴν δυσβουλίας ; Aesch.~\textit{Sept.}~209 πρύμνηθεν ηὗρε μηχανὴν σωτηρίας ; D.~40.59 Τιμοκράτης δὲ μόνος, ὥσπερ ἀπὸ μηχανῆς, μαρτυρεῖ ; Arist.~\textit{Metaph.}~985a18–20 Ἀναξαγόρας τε γὰρ μηχανῇ χρῆται τῷ νῷ πρὸς τὴν κοσμοποιίαν, καὶ ὅταν ἀπορήσῃ διὰ τίν᾽ αἰτίαν ἐξ ἀνάγκης ἐστί, τότε παρέλκει αὐτόν ; Luc.~\textit{Merc.Cond.}~1 ἤ τιν᾽ ἄλλον ἐκ μηχανῆς θεὸν ἐπὶ τῷ καρχησίῳ καθεζόμενον ; Aristid.~\textit{Or.}~46, p.~190 αὐτὸς ὥσπερ θεός τις τῶν ἀπὸ μηχανῆς χεῖρα ὑπερέσχε ; Hsch.~α~6583 ἀποροῦντες· ἀμηχανοῦντες ; Hsch.~β~1339 βυσσοδομεύει· ἐνθυμεῖται. μηχανᾶται. ἀποκεκρυμμένως βουλεύεται, καὶ ἐκ βάθους ; Hsch.~π~3034 πόριμον· εὐπόριστον. πορευτόν ; Phryn.~\textit{PS} p.~57.21~Bekker Ὀνεύειν: τὸ περιάγειν τὰ μηχανήματα, οἷς τούς τε λίθους καὶ τὰ ἄλλα ἕλκουσιν. ὄνος δ᾽ ἐστὶν αὐτὸ τὸ μηχάνημα τὸ διὰ τροχιλίων καὶ σχοινίων ; Suid.~δ~107 Δαίδαλον: ποικίλον κατασκεύασμα. ἀπὸ Δαιδάλου τινὸς μηχανοποιοῦ, ὃς κατασκευάσας διάφορα μηχανήματα γέγονε περιβόητος ; App. Prov.~4.32 Φασὶ γὰρ ὅτι νεανίσκος ὁ Ἀρτέμων ἐγένετο περιμάχητος γυναιξίν· ἄλλοι δὲ ὅτι μηχανοποιὸς ἐγένετο σοφώτατος κατὰ τοὺς Περικλέους χρόνους}
\lp{88 Ἀγάθωνα]}{Ar.~\textit{Thesm.}~29–30 ἐνταῦθ᾽ Ἀγάθων ὁ κλεινὸς οἰκῶν τυγχάνει / ὁ τραγῳδοποιός. — ποῖος οὗτος Ἁγάθων ; Ar.~\textit{Ran.}~83–84 Ἀγάθων δὲ ποῦ ᾽στιν; — ἀπολιπών μ᾽ ἀποίχεται, / ἀγαθὸς ποιητὴς καὶ ποθεινὸς τοῖς φίλοις ; Ar.~fr.~178~K.-A. (Σ \textit{Cyrill.}) Ἀγαθώνειος, ει· αὔλησίς τις μαλακὴ παρ᾽ Ἀριστοφάνει ἐν Γηρυτάδηι. Ἀγάθων γὰρ ὁ τραγικὸς ἐπὶ μαλακίαι διεκωμωιδεῖτο ; \textit{idem} Σ Luc.~\textit{Rh.\,Pr.}~11 (Ἀγάθων τραγωιδίας ποιητὴς εἰς μαλακίαν σκωπτόμενος ; add. Ἀριστοφάνης Γηρυτάδηι B), Hsch.~α~281 ≅ Phot.~α~83 (Ἀγαθώνειον αὔλησιν· τὴν μαλακήν. Ἀγάθων γὰρ ὁ τραγικὸς ἐπὶ μαλακίαι διεβάλλετο) ; Pl.~\textit{Smp.}~174b ἵνα καὶ τὴν παροιμίαν διαφθείρωμεν μεταβαλόντες, ὡς ἄρα καὶ Ἀγάθων᾽ ἐπὶ δαῖτας ἴασιν αὐτόματοι ἀγαθοί ; Plu.~\textit{Quaest.\,conv.}~2.1.12 (634b–d) διὸ καὶ τῶν κωμικῶν ἔνιοι τὴν πικρίαν ἀφαιρεῖν δοκοῦσι τῷ σκώπτειν ἑαυτούς, ὡς Ἀριστοφάνης εἰς τὴν φαλακρότητα καὶ τὴν Ἀγάθωνος ἀπόλειψιν ; Luc.~\textit{Rh.Pr.}~11 πάναβρόν τινα Σαρδανάπαλλον ἢ Κινύραν ἢ αὐτὸν Ἀγάθωνα, τὸν τῆς τραγῳδίας ἐπέραστον ἐκεῖνον ποιητήν ; Ael.~\textit{VH}~13.4 εἶτα συγκλιθέντα αὐτῷ Ἀγάθωνα τὸν τῆς τραγῳδίας ποιητὴν περιλαβὼν κατεφίλει, τετταράκοντα ἐτῶν που γεγονότα ; Philostr.~\textit{VS}~1.9.1 καὶ Ἀγάθων δὲ ὁ τῆς τραγῳδίας ποιητής, ὃν ἡ κωμῳδία σοφόν τε καὶ καλλιεπῆ οἶδε, πολλαχοῦ τῶν ἰάμβων γοργιάζει ; Ath.~5.187c χλευάζει τε τὰ ἰσόκωλα τὰ Ἀγάθωνος καὶ τὰ ἀντίθετα ; Σ Ar.~\textit{Ran.}~83 Οὗτος τραγῳδίας ποιητὴς (Ἀθηναῖος, υἱὸς Τισαμενοῦ). κωμῳδεῖται δὲ εἰς θηλύτητα. (ἦν δὲ καὶ κωμῳδοποιὸς τοῦ Σωκρατικοῦ διδασκαλείου.) ; Suid.~α~124 Ἀγάθων· ὄνομα κύριον. τραγικὸς δὲ ἦν· διεβέβλητο δὲ ἐπὶ μαλακίᾳ. Ἀριστοφάνης· Ἀγάθων δὲ ποῦ ᾽στιν; … οὗτος ὁ Ἀγάθων ἀγαθὸς ἦν τὸν τρόπον, ποθεινὸς τοῖς φίλοις καὶ τὴν τράπεζαν λαμπρός. … ἐκωμῳδεῖτο δὲ εἰς θηλύτητα ; Zen.~2.19 Αὐτόματοι δ᾽ ἀγαθοὶ ἀγαθῶν ἐπὶ δαῖτας ἴενται ; \textit{idem} Tosi n°~516 (ἵενται) ; Pl.~\textit{Prt.}~315d–e νέον τι ἔτι μειράκιον, ὡς μὲν ἐγᾦμαι καλόν τε κἀγαθὸν τὴν φύσιν, τὴν δ᾽ οὖν ἰδέαν πάνυ καλός. ἔδοξα ἀκοῦσαι ὄνομα αὐτῷ εἶναι Ἀγάθωνα ; Pl.~\textit{Smp.}~173a ὅτε τῇ πρώτῃ τραγῳδίᾳ ἐνίκησεν Ἀγάθων ; X.~\textit{Symp.}~8.32 καίτοι Παυσανίας γε ὁ Ἀγάθωνος τοῦ ποιητοῦ ἐραστὴς ; Arist.~\textit{EN}~1140a19 καθάπερ καὶ Ἀγάθων φησὶ τέχνη τύχην ἔστερξε καὶ τύχη τέχνην ; Plu.~\textit{Reg.\,apophth.}~177a–b τοῦ δὲ Εὐριπίδου τὸν καλὸν Ἀγάθωνα περιλαμβάνοντος ἐν τῷ συμποσίῳ καὶ καταφιλοῦντος ἤδη γενειῶντα ; Ael.~\textit{VH}~2.21 Ἀγάθωνος ἤρα τοῦ ποιητοῦ Παυσανίας ὁ ἐκ Κεραμέων. καὶ τοῦτο μὲν διατεθρύληται ; Ael.~\textit{VH}~14.13 ἀλλὰ σύ γε, γενναῖε, λέληθας σεαυτὸν τὸν Ἀγάθωνα ἐκ τοῦ Ἀγάθωνος ἀφανίζων ; Ath.~5.185a κατὰ τὸν καλὸν Ἀγάθωνα ; Ath.~10.454c τὸ δ᾽ αὐτὸ πεποίηκε καὶ Ἀγάθων ὁ τραγῳδιοποιὸς ἐν τῷ Τηλέφῳ ; Σ Ar.~\textit{Ran.}~84 οὗτος δὲ ἀγαθὸς ἦν τὸν τρόπον καὶ τὴν τράπεζαν λαμπρός. καὶ φασὶν ὅτι τὸ Πλάτωνος συμπόσιον ἐν ἑστιάσει αὐτοῦ γέγραπται ; Arist.~\textit{Po.}~1451b21 οἷον ἐν τῷ Ἀγάθωνος Ἀνθεῖ ; Plu.~\textit{Quaest.\,conv.}~2.1.5 (632a–d) καὶ Σωκράτην πάλιν Ἀλκιβιάδης ἔσκωπτεν εἰς ζηλοτυπίαν τὴν περὶ Ἀγάθωνος}
\lp{88 Ἀγάθωνα πεῖσαι]}{Eur.~fr.~170~Kannicht οὐκ ἔστι Πειθοῦς ἱερὸν ἄλλο πλὴν λόγος, / καὶ βωμὸς αὐτῆς ἔστ᾽ ἐν ἀνθρώπου φύσει ; \textit{idem} Ar.~\textit{Ran.}~1391 (om. v. 2) ; Suid.~αι~192 Αἵματι κλαίειν: καθ᾽ ὑπερβολήν, οὐ δάκρυσιν. ἐφ᾽ ὧν μὴ δύναιντο πεῖσαι πάντα πράττοντες, οὕτως ἔλεγον οἱ ἀρχαῖοι ; Hsch.~π~2019 †πέσσαι†· πεῖσαι. ἢ πάσχειν. καὶ νέμειν, καὶ τεχνᾶσθαι τῇ διανοίᾳ ; Diogenian.~4.21 Δῶρα θεοὺς πείθει, καὶ αἰδοίους βασιλῆας: ἐπὶ τῶν διὰ δῶρα}
\lp{88 τραγῳδοδιδάσκαλον]}{Cratin.~fr.~276~K.-A. ἴτω δὲ καὶ τραγῳδίας / ὁ Κλεομάχου διδάσκαλος, / παρατιλτριῶν ἔχων χορὸν / λυδιστὶ τιλλουσῶν μέλη / πονηρά ; Arist.~\textit{Po.}~1449a5 οἱ μὲν ἀντὶ τῶν ἰάμβων κωμῳδοποιοὶ ἐγένοντο, οἱ δὲ ἀντὶ τῶν ἐπῶν τραγῳδοδιδάσκαλοι ; Harp. s.v.~διδάσκαλος ἰδίως διδασκάλους λέγουσι τοὺς ποιητὰς τῶν διθυράμβων ἢ τῶν κωμῳδιῶν ἢ τῶν τραγῳδιῶν ; Ar.~\textit{Ach.}~628 ἐξ οὗ γε χοροῖσιν ἐφέστηκεν τρυγικοῖς ὁ διδάσκαλος ἡμῶν ; Ar.~\textit{Eq.}~507 εἰ μέν τις ἀνὴρ τῶν ἀρχαίων κωμῳδοδιδάσκαλος ἡμᾶς ; Ar.~\textit{Pax}~737 κωμῳδοδιδάσκαλος ἀνθρώπων καὶ κλεινότατος γεγένηται ; Ar.~\textit{Av.}~912–913 οὔκ, ἀλλὰ πάντες ἐσμὲν οἱ διδάσκαλοι / Μουσάων θεράποντες ὀτρηροί ; Ar.~\textit{Av.}~1403–1404 ταυτὶ πεποίηκας τὸν κυκλιοδιδάσκαλον, / ὃς ταῖσι φυλαῖς περιμάχητός εἰμ᾽ ἀεί ; Ar.~fr.~348~K.-A. μήτε Μούσας ἀνακαλεῖν ἑλικοβοστρύχους / μήτε Χάριτας βοᾶν εἰς χορὸν Ὀλυμπίας· / ἐνθάδε γάρ εἰσιν, ὥς φησιν ὁ διδάσκαλος ; Pl.~\textit{Ion}~536a ὁρμαθὸς πάμπολυς ἐξήρτηται χορευτῶν τε καὶ διδασκάλων καὶ ὑποδιδασκάλων ; Isoc.~12.168 τίς γὰρ οὐκ οἶδεν, ἢ τίς οὐκ ἀκήκοε τῶν τραγῳδοδιδασκάλων Διονυσίοις ; Hsch.~δ~1488 διδάσκαλος ⟨χοροῦ⟩· πᾶς ὁ ἐνεργῶν τι περὶ τὴν δραματοποιίαν καὶ ⟨τὴν⟩ τῶν κύκλιων χορῶν παρασκευήν ; Hsch.~κ~4842 κωμῳδοί· κωμικοί, τραγῳδοί, καὶ χορευταὶ τραγῳδοί, οἱ τραγικοὶ ποιηταὶ διδάσκαλοι ; AB~281.25 Μελετών: οἶκος ἐν ᾧ οἱ τραγῳδοὶ ἐμελέτων ; Greg. Cypr.~3.85 Τραγικὸς πίθηκος: ἐπὶ τῶν παρ᾽ ἀξίαν σεμνυνομένων ; Ar.~\textit{Pax}~829 ψυχὰς δύ᾽ ἢ τρεῖς διθυραμβοδιδασκάλων ; Luc.~\textit{Cal.}~1 τοιγάρτοι μυρίας ἤδη τοῖς τραγῳδοδιδασκάλοις ἀφορμὰς εἰς τὰ δράματα ; Ath.~15.699b Ἀλέξανδρος ὁ Αἰτωλὸς ὁ τραγῳδοδιδάσκαλος ; Procop.~\textit{Anecd.}~1.4 κἀν τοῖς τραγῳδοδιδασκάλοις τετάξομαι ; Hsch.~δ~2230 δορύαλλος· τὸ τῶν γυναικῶν μόριον, ἀπὸ τοῦ δέρειν, ἐφ᾽ ὕβρει τοῦ τραγῳδοποιοῦ Δορύλλου}
\lp{89 ἐς Θεσμοφόροιν]}{Ar.~\textit{Lys.}~2 ἢ ᾽ς Πανὸς ἢ ᾽πὶ Κωλιάδ᾽ ἢ ᾽ς Γενετυλλίδος ; Ar.~\textit{Vesp.}~123 νύκτωρ κατέκλινεν αὐτὸν εἰς Ἀσκληπιοῦ ; Theon Prog.~5 ποσταία γυνὴ ἀπ᾽ ἀνδρὸς καθαρὰ εἰς τὸ θεσμοφορεῖον κάτεισιν ; D.L.~9.43 ἐν τῇ τῶν θεσμοφόρων ἑορτῇ μέλλοι τεθνήξεσθαι καὶ τῇ θεῷ τὸ καθῆκον αὐτὴ οὐ ποιήσειν}
\lp{89 ἐς Θεσμοφόροιν ἐλθεῖν]}{App. Prov.~1.88 Γυνὴ εἰς Ἡρακλέους οὐ φοιτᾶ: πρὸς τοὺς ἀναξίους τινῶν πράξεων, παρόσον Ἡρακλῆς ἐδούλευσεν Ὀμφάλῃ}
\lp{89 τί δράσοντ᾽]}{Eur.~\textit{Alc.}~537 ὡς δὴ τί δράσων τόνδ᾽ ὑπορράπτεις λόγον ; Eur.~\textit{HF}~1136 τί φῄς; τί δράσας ; [Eur.]~\textit{Rh.}~99 σὺ δ᾽ ὡς τί δράσων πρὸς τάδ᾽ ὁπλίζῃ χέρας ; Soph.~\textit{Trach.}~160 ἀλλ᾽ ὥς τι δράσων εἷρπε κοὐ θανούμενος}
\lp{89 τί δράσοντ᾽; εἰπέ μοι]}{Eur.~\textit{El.}~624 τί δρῶνθ᾽; ὁρῶ γὰρ ἐλπίδ᾽ ἐξ ἀμηχάνων ; Ar.~\textit{Eq.}~741 εἰπέ νυν τί δρῶν ; Ar.~\textit{Nub.}~900 εἰπὲ τί ποιῶν ; Ar.~\textit{Pax}~210 τοῦ δ᾽ οὕνεχ᾽ ἡμᾶς ταῦτ᾽ ἔδρασαν; εἰπέ μοι ; Ar.~\textit{Eccl.}~865 τί δράσεις; εἰπέ μοι ; Eur.~\textit{El.}~599 λέξον, τί δρῶν ἂν φονέα τεισαίμην πατρός ; Ar.~\textit{Eq.}~1361 τοῦτον τί δράσεις, εἰπέ, τὸν ξυνήγορον ; Eur.~\textit{Cyc.}~138 σὺ δ᾽ ἀντιδώσεις, εἰπέ μοι, χρυσὸν πόσον ; Eur.~\textit{IA}~381 εἰπέ μοι, τί δεινὰ φυσᾷς αἱματηρὸν ὄμμ᾽ ἔχων}
\lp{89–91 τί δράσοντ᾽ … πότερα φανερὸν ἢ λάθρᾳ]}{Eur.~\textit{Hec.}~876–878 πῶς οὖν; τί δράσεις; πότερα φάσγανον χερὶ / λαβοῦσα γραίᾳ φῶτα βάρβαρον κτενεῖς, / ἢ φαρμάκοισιν ἢ ᾽πικουρίᾳ τινί ; Soph.~\textit{El.}~1195 τί δρῶσα; πότερα χερσὶν ἢ λύμῃ βίου}
\lp{90 ἐκκλησιάσοντ᾽]}{Ar.~\textit{Vesp.}~32 ἐκκλησιάζειν πρόβατα συγκαθήμενα ; Ar.~\textit{Eccl.}~161–162 ὡς ἐγὼ / ἐκκλησιάσουσ᾽ οὐκ ἂν προβαίην τὸν πόδα / τὸν ἕτερον ; Isoc.~8.2 ἥκομεν γὰρ ἐκκλησιάσοντες περὶ πολέμου καὶ εἰρήνης ; Hsch.~α~5933 ἀπελλάζειν· ἐκκλησιάζειν. Λάκωνες ; Ar.~\textit{Plu.}~171 ἐκκλησία δ᾽ οὐχὶ διὰ τοῦτον γίγνεται ; \textit{idem} Suid.~ε~471 (add. Ἐκκλησίαν· συναγωγὴν ὄχλου. Ἀριστοφάνης Πλούτῳ ; πλοῦτον ; γίνεται) ; Hsch.~α~720 ἀγορεύειν· λέγειν ἐν ἐκκλησίᾳ ; Hsch.~ε~1442 ἐκκλησία· σύνοδος, συναγωγή, πανήγυρις ; Suid.~ε~469 Ἐκκλησιάζων: εἰς ἐκκλησίαν συγκαλῶν ; App. Prov.~2.15 ἱστόρηται γὰρ, ὡς τῶν ἐν τῇ Ἀσίᾳ πόλεων ἐκκλησιαζουσῶν καὶ ψήφους καταβαλλομένων}
\lp{90–91 ἐκκλησιάσοντ᾽ ἐν ταῖς γυναιξὶ … λέξονθ᾽ ὑπὲρ ἐμοῦ]}{Ar.~\textit{Thesm.}~184–187 ἐὰν γὰρ ἐγκαθεζόμενος λάθρᾳ / ἐν ταῖς γυναιξίν, ὡς δοκῶν εἶναι γυνή, / ὑπεραποκρίνῃ μου, σαφῶς σώσεις ἐμέ. / μόνος γὰρ ἂν λέξειας ἀξίως ἐμοῦ}
\lp{90 ἐν ταῖς γυναιξί]}{Ar.~\textit{Ach.}~630 διαβαλλόμενος δ᾽ ὑπὸ τῶν ἐχθρῶν ἐν Ἀθηναίοις ταχυβούλοις ; Ar.~\textit{Vesp.}~1185 μῦς καὶ γαλᾶς μέλλεις λέγειν ἐν ἀνδράσιν ; Ar.~fr.~161~K.-A. ἐν τοῖσι συνδείπνοις ἐπαινῶν Αἰσχύλον ; Soph.~\textit{Aj.}~293 γύναι, γυναιξὶ κόσμον ἡ σιγὴ φέρει ; \textit{idem} Arist.~\textit{Pol.}~1260a29–31 (διὸ δεῖ, ὥσπερ ὁ ποιητὴς εἴρηκε περὶ γυναικός, οὕτω νομίζειν ἔχειν περὶ πάντων· γυναικὶ pro γύναι, γυναιξὶ ; add. ἀλλ᾽ ἀνδρὶ οὐκέτι τοῦτο), Suid.~β~210, Tosi n°~295 (Γυναιξὶ κόσμον ἡ σιγὴ φέρει) (om. γύναι) ; Pl.~\textit{Grg.}~464d ὥστ᾽ εἰ δέοι ἐν παισὶ διαγωνίζεσθαι ὀψοποιόν τε καὶ ἰατρόν, ἢ ἐν ἀνδράσιν οὕτως ἀνοήτοις ὥσπερ οἱ παῖδες ; Ar.~\textit{Nub.}~432–433 ἐν τῷ δήμῳ γνώμας οὐδεὶς νικήσει πλείονας ἢ σύ. / μὴ ᾽μοί γε λέγειν ; Pl.~\textit{Grg.}~521e κρινοῦμαι γὰρ ὡς ἐν παιδίοις ἰατρὸς ἂν κρίνοιτο κατηγοροῦντος ὀψοποιοῦ ; Pl.~\textit{Menex.}~235d εἰ μὲν γὰρ δέοι Ἀθηναίους ἐν Πελοποννησίοις εὖ λέγειν ἢ Πελοποννησίους ἐν Ἀθηναίοις, ἀγαθοῦ ἂν ῥήτορος δέοι τοῦ πείσοντος καὶ εὐδοκιμήσοντος ; Arist.~\textit{Rh.}~1367b7–9 σκοπεῖν δὲ καὶ παρ᾽ οἷς ὁ ἔπαινος· ὥσπερ γὰρ ὁ Σωκράτης ἔλεγεν, οὐ χαλεπὸν Ἀθηναίους ἐν Ἀθηναίοις ἐπαινεῖν ; \textit{idem} Arist.~\textit{Rh.}~1415b30–32 (ὃ γὰρ λέγει Σωκράτης ἐν τῷ ἐπιταφίῳ, ἀληθές, ὅτι οὐ χαλεπὸν Ἀθηναίους ἐν Ἀθηναίοις ἐπαινεῖν, ἀλλ᾽ ἐν Λακεδαιμονίοις) ; NT~1 Cor.~14.34 Αἱ γυναῖκες ἐν ταῖς ἐκκλησίαις σιγάτωσαν ; Eur.~\textit{Heracl.}~476–477 γυναικὶ γὰρ σιγή τε καὶ τὸ σωφρονεῖν / κάλλιστον, εἴσω θ᾽ ἥσυχον μένειν δόμων}
\lp{90 χἂν δέῃ]}{Ar.~\textit{Lys.}~887 χἂ δυσκολαίνει πρὸς ἐμὲ καὶ βρενθύεται ; Th.~1.22.1 ὡς δ᾽ ἂν ἐδόκουν ἐμοὶ ἕκαστοι περὶ τῶν αἰεὶ παρόντων τὰ δέοντα μάλιστ᾽ εἰπεῖν ; Ar.~\textit{Ran.}~161 οἵ σοι φράσουσ᾽ ἁπαξάπανθ᾽ ὧν ἂν δέῃ ; Soph.~\textit{OC}~12–13 μανθάνειν γὰρ ἥκομεν / ξένοι πρὸς ἀστῶν, ἃν δ᾽ ἀκούσωμεν τελεῖν ; Plaut.~\textit{Mil.}~3 ut, ubi usus veniat, contra conserta manu ; X.~\textit{An.}~1.3.6 ἐγὼ σὺν ὑμῖν ἕψομαι καὶ ὅ τι ἂν δέῃ πείσομαι}
\lp{91 λέξονθ᾽ ὑπὲρ ἐμοῦ]}{Ar.~\textit{Vesp.}~951 ὑπεραποκρίνεσθαι κυνός, λέξω δ᾽ ὅμως ; Ar.~\textit{Ach.}~355 ἐμοῦ ᾽θέλοντος ὑπὲρ ἐπιξήνου λέγειν ; Ar.~\textit{Thesm.}~466–471 τὸ μὲν ὦ γυναῖκες ὀξυθυμεῖσθαι σφόδρα / Εὐριπίδῃ … ὅμως δ᾽ ἐν ἀλλήλαισι χρὴ δοῦναι λόγον ; X.~\textit{An.}~6.6.18 συμπέμψατε μέντοι μοι ὑμῶν αὐτῶν ἑλόμενοι πρὸς Κλέανδρον οἵτινες, ἄν τι ἐγὼ παραλίπω, καὶ λέξουσιν ὑπὲρ ἐμοῦ ; Ar.~\textit{Lys.}~210 λεγέτω δ᾽ ὑπὲρ ὑμῶν μί᾽ ἅπερ ἂν κἀγὼ λέγω}
\lp{91 πότερα φανερὸν ἢ λάθρᾳ]}{Eur.~\textit{Ba.}~817–818 ἀλλ᾽ ἐξιχνεύσουσίν σε, κἂν ἔλθῃς λάθρᾳ. / ἀλλ᾽ ἐμφανῶς· καλῶς γὰρ ἐξεῖπας τάδε ; \textit{Od.}~19.299 ἤδη δὴν ἀπεών, ἢ ἀμφαδὸν ἦε κρυφηδόν ; D.~20.74 ὅσῳ τὸ φανερῶς τοῦ λάθρᾳ κρεῖττον ; Ar.~\textit{Vesp.}~1017–1018 τὰ μὲν οὐ φανερῶς ἀλλ᾽ ἐπικουρῶν κρύβδην ἑτέροισι ποιηταῖς ; Th.~6.34.2 ἐθελήσειαν ἡμῖν ἤτοι κρύφα γε ἢ φανερῶς ἢ ἐξ ἑνός γέ του τρόπου ἀμῦναι ; And.~4.21 τοὺς μὲν λάθρᾳ ἀδικοῦντας κολάζοντες, τοὺς δὲ φανερῶς ἀσελγαίνοντας θαυμάζοντες ; Pl.~\textit{Lg.}~838b μήτε φανερῶς μήτε λάθρᾳ ; Lys.~8.5 τὸν αὐτὸν ἄνδρα λάθρᾳ μὲν ἐλοιδορεῖτε, φανερῶς δὲ φίλον ἐνομίζετε ; Zen.~5.59 Ὁ Σκύθης τὸν ἵππον: ἐπὶ τῶν κρύφα τινὸς ἐφιεμένων, φανερῶς δὲ ἀπωθουμένων καὶ διαπτυόντων αὐτὸ εἴρηται ἡ παροιμία ; Pl.~\textit{Smp.}~182d λέγεται κάλλιον τὸ φανερῶς ἐρᾶν τοῦ λάθρᾳ ; App. BC~3.13.91 οὔτε φανερῶς οὔτε λάθρᾳ ζητοῦντες εὕρισκον ; Philo Spec. leg.~3.87 φανερῶς ἢ λάθρᾳ}
\lp{91 φανερόν]}{Ar.~\textit{Thesm.}~431 ταῦτ᾽ ἐγὼ φανερῶς λέγω ; Arist.~\textit{Pol.}~1332a25–27 ὥσπερ εἰ τοῦ κιθαρίζειν λαμπρὸν καὶ καλῶς ; Soph.~\textit{Phil.}~345 λέγοντες, εἴτ᾽ ἀληθὲς εἴτ᾽ ἄρ᾽ οὖν μάτην ; Hsch.~ε~4105 ἐπάϊστος· φανερός. κατάφωρος}
\lp{91 λάθρᾳ]}{Ar.~\textit{Ran.}~1168 λάθρᾳ γὰρ ἦλθεν οὐ πιθὼν τοὺς κυρίους ; Ar.~\textit{Vesp.}~347 ἥτις σε λάθρᾳ τἀνδρὸς τουδὶ καταβῆναι δεῦρο ποιήσει ; Ar.~\textit{Thesm.}~482 εἶτα καταβαίνω λάθρᾳ ; Hsch.~ε~3273 ἐν παραβύστῳ· ἐν μυστηρίῳ … ἐν κρύπτῷ, λεληθότως … λάθρα ; Hsch.~κ~1084 καταδύμεναι· λάθρα ὑπεισελθεῖν ; Hsch.~κ~3073 κλωπᾶσθαι· ἐπιθυμεῖν. θέλειν. καὶ τὸ λάθρα καὶ ἀψοφητὶ ἰέναι καὶ πράσσειν, καὶ κλέπτεσθαι ; Hsch.~λ~130 λαίθαργοι· κύνες κρύφα δάκνοντες ; Zen.~4.90 ὁ προςσαίνων μὲν, λάθρα δὲ δάκνων ; Hsch.~κ~4242 κρύβδην· λάθρα ; Hsch.~λ~103 λάθρᾳ· κρυφίως, ἠρεμα, ἡσύχως}
\lp{92 λάθρᾳ]}{Zen.~1.41 Ἄϊδος κυνῆ: πρὸς τοὺς ἐπικρύπτοντας ἑαυτοὺς διά τινων μηχανημάτων ; Ar.~\textit{Ach.}~390 σκοτοδασυπυκνότριχά τιν᾽ Ἄιδος κυνῆν ; AB~194 s.v.~ὑπὸ μάλης Ὑπὸ μάλης: ἀντὶ τοῦ κρύφα, ὃ λέγουσι λάθρᾳ καὶ ἀφανῶς. τοῦτ᾽ ἔστιν ὑπὸ μασχάλης ; AB~292 s.v.~παράβυστον Παράβυστον: τὸ λάθρᾳ γινόμενον. ἦν δὲ καὶ παράβυστον δικαστήριον Ἀθήνησιν, ὃ λάθρᾳ ἔκρινεν ; Diogenian.~2.40 Ἑστίᾳ θύειν: ἐπὶ τῶν κρύφα πραττόντων. Οἱ γὰρ Ἑστίᾳ θύοντες οὐδὲν ἐξέφερον τῆς θυσίας}
\lp{92 λάθρᾳ, στολὴν γυναικὸς ἠμφιεσμένον]}{Eur.~\textit{Ba.}~821–823 στεῖλαί νυν ἀμφὶ χρωτὶ βυσσίνους πέπλους. / τί δὴ τόδ᾽; ἐς γυναῖκας ἐξ ἀνδρὸς τελῶ; / μή σε κτάνωσιν, ἢν ἀνὴρ ὀφθῇς ἐκεῖ ; Plu.~\textit{Caes.}~10.1 ὁ Κλώδιος οὔπω γενειῶν καὶ διὰ τοῦτο λήσειν οἰόμενος ἐσθῆτα καὶ σκευὴν ψαλτρίας ἀναλαβὼν ἐχώρει, νέᾳ γυναικὶ τὴν ὄψιν ἐοικώς ; Ael.~fr.~44 Βάττος ὁ Κυρήνην κτίσας τῆς Θεσμοφόρου τὰ μυστήρια ἐγλίχετο μαθεῖν … ἀθρόαι ὑφ᾽ ἑνὶ συνθήματι ἐπὶ τὸν Βάττον ᾖξαν, ἵνα αὐτὸν ἀφέλωνται τοῦ ἔτι εἶναι ἄνδρα ; Lys.~6.51 οὗτος γὰρ ἐνδὺς στολήν, μιμούμενος τὰ ἱερὰ ἐπεδείκνυε τοῖς ἀμυήτοις καὶ εἶπε τῇ φωνῇ τὰ ἀπόρρητα ; Euseb.~\textit{HE}~7.32.11 λεληθότως γυναικείαν στολὴν ἀμπισχομένους νύκτωρ ; Polyaen. Exc.~28 αὐτὸς δὲ νέους στολίσας γυναικείαν στολὴν μετὰ ἐγχειριδίων προῆγεν ὡς πρὸς τὴν ἑορτὴν τῆς Δήμητρος ; Hdt.~4.78.4–5 τὴν στολὴν ἀποθέμενος τὴν Σκυθικὴν λάβεσκε ἂν Ἑλληνίδα ἐσθῆτα … μή τίς μιν Σκυθέων ἴδοι ἔχοντα ταύτην τὴν στολήν ; Hdt.~6.134.2 ἕρκος θεσμοφόρου Δήμητρος ὑπερθορεῖν, οὐ δυνάμενον τὰς θύρας ἀνοῖξαι ; Plu.~\textit{Cic.}~28.1 εἰς τὴν οἰκίαν αὐτοῦ παρεισῆλθε κρύφα, λαβὼν ἐσθῆτα καὶ σκευὴν ψαλτρίας· ἔθυον γὰρ αἱ γυναῖκες τὴν ἀπόρρητον ἐκείνην καὶ ἀθέατον ἀνδράσι θυσίαν ; Plu.~\textit{Sol.}~8.4–5 καταλαβὼν αὐτόθι πάσας τὰς γυναῖκας τῇ Δήμητρι τὴν πάτριον θυσίαν ἐπιτελούσας … τῶν δὲ νεωτέρων τοὺς μηδέπω γενειῶντας ἐνδύμασι καὶ μίτραις καὶ ὑποδήμασι τοῖς ἐκείνων σκευασαμένους ; Arr.~\textit{Anab.}~1.25.9 ὁ Ἀμφοτερὸς στολὴν ἐνδὺς ἐπιχώριον, ὡς μὴ γνώριμος εἶναι κατὰ τὴν ὁδόν, λανθάνει ἐφικόμενος παρὰ Παρμενίωνα ; Paus.~4.17.1 τὰς γυναῖκας ἀγούσας ἑορτὴν … τὸν δὲ Ἀριστομένην τύπτουσαι ταῖς δᾳσὶ ζῶντα αἱροῦσιν}
\lp{92 στολήν]}{Ar.~\textit{Thesm.}~136 ποδαπὸς ὁ γύννις; τίς πάτρα; τίς ἡ στολή ; Poll.~4.115 καὶ σκευὴ μὲν ἡ τῶν ὑποκριτῶν στολή — ἡ δ᾽ αὐτὴ καὶ σωμάτιον ἐκαλεῖτο ; Eur.~\textit{Hel.}~1382–1383 πέπλους δ᾽ ἀμείψασ᾽ ἀντὶ ναυφθόρου στολῆς / ἐγώ νιν ἐξήσκησα ; Sapph.~fr.~57 Voigt τίς δ᾽ ἀγροΐωτις θέλγει νόον … / ἀγροΐωτιν ἐπεμμένα στόλαν … / οὐκ ἐπισταμένα τὰ βράκε᾽ ἔλκην ἐπὶ τῶν σφύρων ; Hsch.~σ~1902 στολή· στάσις. πορεία. κίνησις. ἔνδυμα, καὶ περιβόλαιον}
\lp{92 στολὴν γυναικός]}{Ar.~\textit{Thesm.}~851 πάντως ὑπάρχει μοι γυναικεία στολή ; Ar.~\textit{Ran.}~46 ὁρῶν λεοντῆν ἐπὶ κροκωτῷ κειμένην ; Stratt.~fr.~75~K.-A. (Zen. vulg.~2.93) γαλῆ χιτώνιον· ὁμοία ἡ παροιμία αὕτη τῆι οὐ πρέπει γαλῆι κροκωτός, ἐπειδὴ γαλῆ κατὰ πρόνοιαν Ἀφροδίτης γυνὴ γενομένη ἐν χιτῶνι κροκωτῶι οὖσα ἐπέδραμε μυΐ. μέμνηται ταύτης Στράτ⟨τ⟩ις ; \textit{idem} Zen.~2.93 (Γαλῇ χιτώντον ; κροκωτόν ; Στράττις), Diogenian.~3.82 (Γαλῇ χιτώνιον· ἐπὶ τῶν τὰ μὴ ἑαυτῶν περιβλήματα ἐνδυόντων pro ὁμοία … Στράτ⟨τ⟩ις) ; Eur.~\textit{Ba.}~828–836 τίνα στολήν; ἦ θῆλυν; ἀλλ᾽ αἰδώς μ᾽ ἔχει. … / οὐκ ἂν δυναίμην θῆλυν ἐνδῦναι στολήν ; Apollod. Bibl.~3.13.8 κρύψασα ἐσθῆτι γυναικείᾳ ὡς παρθένον Λυκομήδει παρέθετο ; Paus.~8.20.3 ταύτην οἷα δὴ παρθένος πλεξάμενος τὴν κόμην καὶ ἐσθῆτα ἐνδὺς γυναικείαν ; Dio Chrys.~4.108 ὥσπερ οἱ τῶν δραμάτων ὑποκριταί ἀπορρίψας μεταξὺ τὴν γυναικείαν στολήν, ἔπειτα στρατιώτου τινὸς ἢ ῥήτορος στολὴν ἁρπάσας ; Hsch.~γ~1015 γύννις· δειλός, ἄνανδρος, γυναικώδης· Ἀριστοφάνης ; Phryn.~\textit{PS} p.~18.23~Bekker (Ar. fr.~869~K.-A.) Αἱ τῶν γυναικῶν παγίδες: τοὺς κόσμους καὶ τὰς ἐσθῆτας, αἷς χρῶνται αἱ γυναῖκες καλλωπιζόμεναι, παγίδας εἶπεν Ἀριστοφάνης ; Ar.~\textit{Thesm.}~250–262 ἀλλ᾽ ἱμάτιον γοῦν χρῆσον ἡμῖν τουτῳὶ / καὶ στρόφιον … τὸν κροκωτὸν πρῶτον ἐνδύου λαβών ; Eur.~\textit{Ba.}~980–981 ἐπὶ τὸν ἐν γυναικομίμῳ στολᾷ / λυσσώδη κατάσκοπον μαινάδων ; Hdt.~4.116.2 καὶ ἐς πόλεμον φοιτῶσαι καὶ στολὴν τὴν αὐτὴν τοῖσι ἀνδράσι φορέουσαι ; J.~\textit{AJ}~4.301 ὡς μήτε γυναῖκα ἀνδρικῇ σκευῇ χρῆσθαι μήτ᾽ ἄνδρα στολῇ γυναικείᾳ ; Hsch.~γ~1016 γυναικὸς ἄρ᾽ ἀντὶ τέτυξο· ἀντὶ γυναικὸς ἄρα γεγέννησο ; Hsch.~θ~490 θηλυδρίας· ὁ τεθηλυμμένος ; Hsch.~θ~491 θηλυδριῶδες· τὸ κατακεκλασμένον ; Suid.~γ~41 ἀπέστειλε νεανίσκους διασκευάσας εἰς Γάλλους, μετ᾽ αὐλητῶν ἐν γυναικείαις στολαῖς ; Suid.~γ~496 Γυναίκισις· ἡ γυναικεία μίμησις. πρὶν τῆς ἑτέρας γυναικίσεως δίκην δοῦναι (cf. γ 495 Γυναικίας δὲ ἀνήρ, ὁ θηλυμανής … καὶ Γυναικιζόμενος, ὁ θηλυνόμενος) ; Suid.~δ~454 νέος μὲν ὢν Βάταλος ἐκλήθη, ὡς καὶ γυναικείᾳ ἐσθῆτι πολλάκις χρησάμενος ; Diogenian.~7.94 Πίθηκος ἐν πορφύρᾳ: οἱ φαῦλοι κἂν καλοῖς περιβληθῶσιν, ὅμως διαφαίνονται πονηροί ; Plu.~\textit{Prov.Alex.}~1.60 Εἰς Μασσαλίαν πλεύσειας: οἱ Μασσαλιῶται θηλύτερον ἕξων, καὶ στολαῖς ποικίλαις καὶ ποδήρεσι χρώμενοι, ἔτι δὲ τὰς κόμας μεμυρισμένας ἀναδεδεμένοι, καὶ διὰ ταύτην τὴν μαλακίαν ἀσχημονοῦντες ; Tosi n°~168 (Meos tam suspicione quam crimine…) Meos tam suspicione quam crimine iudico carere ; Babr.~139 Ὄνος λεοντῆν ἰσχίοις ἐφαπλώσας / ἔφασκεν εἶναι πᾶσι φοβερὸς ἀνθρώποις ; Plu.~\textit{Quaest.\,Graec.}~58 (304c) διὰ τί παρὰ Κῴοις ὁ τοῦ Ἡρακλέους ἱερεὺς ἐν Ἀντιμαχείᾳ γυναικείαν ἐνδεδυμένος ἐσθῆτα … κατάρχεται τῆς θυσίας ; LXX \textit{Deut.}~22.5 οὐδὲ μὴ ἐνδύσηται ἀνὴρ στολὴν γυναικίαν ; Dio Chrys.~21.7 ἀμπείχετο ἐσθῆτα γυναικείαν ; Hsch.~σ~1903 στολμὸν χρωτός· στολισμὸν σώματος ; Poll.~4.117–118 ὁ δὲ κροκωτὸς ἱμάτιον· Διόνυσος δ᾽ αὐτῷ ἐχρῆτο [καὶ μασχαλιστῆρι ἀνθινῷ καὶ θύρσῳ] … γυναικείας δὲ συρτὸς [πορφυροῦς, παράπηχυ λευκὸν τῆς βασιλευούσης] ; Tosi n°~2186 (Εἵματ' ἀνήρ) Εἵματ᾽ ἀνήρ}
\lp{92 ἠμφιεσμένον]}{Parth.~15.2 ἀμφιεσάμενος δὲ γυναικείας ἀμπεχόνας καὶ ὁμοιωθεὶς κόρῃ συνεθήρα αὐτῇ ; Ar.~\textit{Thesm.}~840 τὴν Ὑπερβόλου καθῆσθαι μητέρ᾽ ἠμφιεσμένην ; Ar.~\textit{Eccl.}~879 ἕστηκα καὶ κροκωτὸν ἠμφιεσμένη ; Pl.~\textit{Grg.}~523c ψυχὰς πονηρὰς ἔχοντες ἠμφιεσμένοι εἰσὶ σώματά τε καλὰ καὶ γένη καὶ πλούτους ; Zos.~3.3.5 ἀμφιέσας δὲ ἐσθῆτι γυναικείᾳ διὰ τοῦ στρατοπέδου παρήγαγεν ; Hsch.~α~4013 ἀμφιέννυμι· ἐνδύομαι ; Ar.~\textit{Vesp.}~1172 Δοθιῆνι σκόροδον ἠμφιεσμένῳ ; Ar.~\textit{Eccl.}~1057 ἐξ αἵματος φλύκταιναν ἠμφιεσμένη ; Hsch.~ε~6782 εὖ εἱμένος· εὖ ἠμφιεσμένος}
\lp{93 τὸ πρᾶγμα κομψὸν καὶ σφόδρ᾽ ἐκ τοῦ σοῦ τρόπου]}{\textit{Test.} Suid.~κ~2025 Κομψόν: περίτρανον, περίλαλον, πανοῦργον, ἀπατητικόν, πιθανόν, τεχνικόν … καὶ αὖθις· τὸ πρᾶγμα κομψὸν καὶ σφόδρ᾽ ἐκ τοῦ σοῦ τρόπου}
\lp{93–94 τὸ πρᾶγμα κομψὸν … ἡμέτερος ὁ πυραμοῦς]}{\textit{Test.} Suid.~τ~439 Τεχνάζειν: τὸ πανούργως τι διαπράττεσθαι. τὸ πρᾶγμα κομψὸν καὶ σφόδρ᾽ ἐκ τοῦ σοῦ τρόπου· τοῦ γὰρ τεχνάζειν ἡμέτερος ὁ πυραμοῦς. εἶδος πλακοῦντος ὁ πυραμοῦς· οὗτος δὲ ἐτίθετο τοῖς διαπαννυχίζουσιν ἔπαθλον, καὶ ὁ ἐγρηγορὼς ἕως πρωΐας ὡς νικῶν αὐτὸν ἐλάμβανε. μεταφορικῶς οὖν αὐτῷ ἐχρήσατο ὡς νικῶντος αὐτοῦ πανουργίᾳ πάντας}
\lp{93 κομψόν]}{Ar.~\textit{Eq.}~18 εἴποιμ᾽ ἂν αὐτὸ δῆτα κομψευριπικῶς ; Ar.~\textit{Ran.}~967 οὑμοὶ δὲ Κλειτοφῶν τε καὶ Θηραμένης ὁ κομψός ; Eur.~fr.~16.2~Kannicht μή μοι τὰ κομψὰ ποικίλοι γενοίατο ; Eur.~fr.~188.5~Kannicht ἄλλοις τὰ κομψὰ ταῦτ᾽ ἀφεὶς σοφίσματα ; \textit{idem} Pl.~\textit{Grg.}~486c (ταῦτα ; om. σοφίσματα) ; Σ Ar.~\textit{Eq.}~18 κομψευριπιδικῶς: Ἀντὶ τοῦ εὐριπιδικῶς καὶ δεινῶς, οἷον πανούργως. κομψοὶ δὲ παρὰ τοῖς ἀρχαίοις ἐλέγοντο οἱ πανοῦργοι ; Σ Ar.~\textit{Thesm.}~93 ἀεὶ γὰρ τὸν Εὐριπίδην ὡς πανοῦργον διαβάλλει ; Hsch.~κ~3484 κομψόν· ἔντιμον. σπουδαῖον. γελοιαστήν. τωθαστήν. πανοῦργον. λάλον. στωμύλον. σκωπτικόν ; Moer. κ~2 Hansen κομψοὺς Πλάτων οὐ τοὺς πανούργους ἀλλὰ τοὺς βελτίστους· ἐπαινῶν γὰρ τοὺς Ἀσκληπιάδας κομψοὺς λέγει Πολιτείας γ´ (408 b) ; Antiatt. p.~102.23~Bekker Κομψόν: ἐπὶ τοῦ κομψευομένου φασὶ δεῖν λέγειν, καὶ τοὺς ἐπὶ τῶν ἀσθενούντων λέγοντας, ὅτι κομψότερον διάκεινται, μέμφονται. ἀλλὰ Πλάτων Πολιτείας τετάρτῳ κομψοὺς οὐ τοὺς κεκομψευμένους λέγει, ἀλλὰ τοὺς ἀμείνονας ; AB~273.12 Κομψόν: τὸ πανοῦργον καὶ τετεχνιτευμένον ; Ar.~\textit{Nub.}~649 πρῶτον μὲν εἶναι κομψὸν ἐν συνουσίᾳ ; Ar.~\textit{Nub.}~1030–1031 πρὸς τάδε σ᾽ ὦ κομψοπρεπῆ μοῦσαν ἔχων, / δεῖ σε λέγειν τι καινόν ; Ar.~\textit{Av.}~195 μὴ ᾽γὼ νόημα κομψότερον ἤκουσά πω ; Ar.~\textit{Thesm.}~460 κομψότερον ἔτ᾽ ἢ τὸ πρότερον ; Ar.~fr.~719~K.-A. ῥήματά τε κομψὰ καὶ παίγνι᾽ ἐπιδεικνύναι / πάντ᾽ ἀπ᾽ ἀκροφυσίων κἀπὸ καναβευμάτων ; Cratin.~fr.~342~K.-A. τίς δὲ σύ; κομψός τις ἔροιτο θεατής. / ὑπολεπτολόγος, γνωμιδιώκτης, εὐριπιδαριστοφανίζων ; Alex.~fr.~9.1–2~K.-A. κομψόν γε τοῦτ᾽ ἐστὶν παρ᾽ ὑμῖν, ὦ Σόλων, / ἐν ταῖς Ἀθήναις, δεξιῶς θ᾽ εὑρημένον ; Araros fr.~8~K.-A. „τὰ ϰομψὰ δὴ ταῦτα νω†τ†αλεύματα ; Eur.~\textit{Phoen.}~1407–1408 καί πως νοήσας Ἐτεοκλῆς τὸ Θεσσαλὸν / ἐσήγαγεν σόφισμ᾽ ὁμιλίᾳ χθονός ; Eur.~\textit{Supp.}~426 κομψός γ᾽ ὁ κῆρυξ καὶ παρεργάτης λόγων ; Eur.~\textit{Tro.}~651 ἔσω τε μελάθρων κομψὰ θηλειῶν ἔπη ; Eur.~fr.~473.1~Kannicht φαῦλον ἄκομψον, τὰ μέγιστ᾽ ἀγαθόν ; [Eur.]~\textit{Rh.}~625 τρίβων γὰρ εἶ τὰ κομψὰ καὶ νοεῖν σοφός ; Pl.~\textit{Grg.}~521e καὶ οὐκ ἐθέλων ποιεῖν ἃ σὺ παραινεῖς, τὰ κομψὰ ταῦτα ; Plu.~\textit{Amat.}~24 (769e–770d) τὰ δ᾽ ὑπ᾽ Εὐριπίδου ῥηθέντ᾽ ἐστὶ κομψά· ἔφη γὰρ Ἀγάθωνα τὸν καλὸν ἤδη γενειῶντα περιβάλλων καὶ κατασπαζόμενος, ὅτι τῶν καλῶν καὶ τὸ μετόπωρον ; Σ Ar.~\textit{Nub.}~1030 κομψοπρεπῆ: ἀντὶ τοῦ πανοῦργον ; Hsch.~α~2475 ἄκομψον· ἀπάνουργον, ἁπλοῦν· Ἀρχιλόχοις (Cratin. fr. 15 K.–A.) ; \textit{idem} AB~213.27 (om. ἁπλοῦν· Ἀρχιλόχοις) ; Hsch.~κ~3177 κόβαλος· πανοῦργος. κακοῦργος. στωμύλος, λάλος, ἀφ᾽ οὗ καὶ ὁ κομψός. ἔνιοι μάταιος. ἄλλοι κροταφιστής, ἄσωτος, τωθαστής, ἀπατεών. κακόλαλος ; Hsch.~κ~3482 κομψεία· ποικίλη λαλιά, καὶ πανουργία ; Hsch.~κ~3486 κομψούς· βέλτιστους, πανούργους. σπουδαίους, ἐντίμους. καλούς ; Moer. κ~56 Hansen κομψεία Ἀττικοί· πανουργία Ἕλληνες ; P.Oxy.~2087.23 ἄκομψος· ἄκακος· κομψὸς γὰρ ὁ πανοῦργος ; Synag.~κ~400 κομψεία· ἐλαφρία, ἀστειότης, πιθανολογία, ἀλαζονεία ; Synag.~κ~402 κομψόν· σπουδαῖον, ἔντιμον, εὐπρεπές τε καὶ χαρίεν ; Phot.~κ~926 κομψόν· περίτρανον, περίλαλον, πανοῦργον, ἀπατητικόν, πιθανόν, τεχνικόν· σημαίνει δὲ καὶ τὸ ἀγαθὸν καὶ τὸ σπουδαῖον (cf. κ 927 κομψῶς· ἀντὶ τοῦ καλῶς. Ἀριστοφάνης) ; Suid.~κ~2023 Κομψευριπικῶς: πανούργως, κατὰ Εὐριπίδην ; Macar.~4.90 Κανθάρου σοφώτερος: ἐπὶ τῶν σφόδρα σοφῶν καὶ πανούργων· Κάνθαρος γάρ τις ἐγένετο κακοῦργος καὶ πονηρός, ὃς ἐπὶ προδοσίᾳ κριθεὶς Ἀθήνησιν ἀπέθανεν ; App. Prov.~3.20 Θετταλικὸν σόφισμα: ἐπὶ τῶν πανούργων καὶ εὐμεθόδων ἔργων ; Ar.~\textit{Vesp.}~1317 ἐπὶ τῷ κομᾷς καὶ κομψὸς εἶναι προσποιεῖ ; Ar.~\textit{Lys.}~89 κομψότατα τὴν βληχώ γε παρατετιλμένη ; Cratin.~fr.~182.3~K.-A. οἱ δ᾽ αὖθ᾽ ἡμεῖς, ὡς ὁ παλαιὸς / λόγος, αὐτομάτους ἀγαθοὺς ἰέναι / κομψῶν ἐπὶ δαῖτα θεατῶν ; Eup.~fr.~172.2~K.-A. ἀλλ᾽ ἀκούσαθ᾽ ὡς ἐσμὲν ἅπαντα κομψοὶ / ἄνδρες ; Dionys. Com. fr.~3.1–4~K.-A. ἄγε δὴ Δρόμων, νῦν εἴ τι κομψὸν ἢ σοφὸν / ἢ γλαφυρὸν οἶσθα τῶν σεαυτοῦ πραγμάτων, / φανερὸν ποίησον τοῦτο τῷ διδασκάλῳ. / νῦν τὴν ἀπόδειξιν τῆς τέχνης αἰτῶ σ᾽ ἐγώ ; Eur.~\textit{Cyc.}~315 κομψὸς γενήσῃ καὶ λαλίστατος, Κύκλωψ ; Soph.~\textit{Ant.}~324 κόμψευέ νυν τὴν δόξαν ; Pl.~\textit{R.}~408b πάνυ κομψούς, ἔφη, λέγεις Ἀσκληπιοῦ παῖδας (cf. 405d ἀναγκάζειν τοὺς κομψοὺς Ἀσκληπιάδας) ; Hsch.~α~370 ἀγαυρός· αὐθάδης, κομψός. κακός ; Hsch.~ε~1347 †ἐκδύον· καλόν. κομψόν ; Hsch.~κ~3481 κομψά· τὰ στρογγύλα ; Hsch.~κ~3483 κομψεύεται· κολακεύει. σεμνύνεται. πιθανοποιεῖ ; \textit{idem} Synag.~κ~401 (om. σεμνύνεται) ; Hsch.~κ~3485 κομψότερον· βελτιώτερον. ἐλαφρότερον ; Hsch.~υ~659 ὑποκομψίας· τὸ κομψὸν ὑποφαίνων}
\lp{93 σφόδρ᾽ ἐκ τοῦ σοῦ τρόπου]}{Ar.~\textit{Vesp.}~1002 ἄκων γὰρ αὔτ᾽ ἔδρασα κοὐ τοὐμοῦ τρόπου ; Ar.~\textit{Thesm.}~574 φίλαι γυναῖκες ξυγγενεῖς τοὐμοῦ τρόπου ; \textit{idem} Suid.~ξ~83 (om. φίλαι γυναῖκες ; add. ἀντὶ τοῦ τὰ αὐτὰ πράττουσαι) ; Luc.~\textit{Pro\,imag.}~14 οὐ γὰρ πρὸς τοῦ σοῦ τρόπου ; Eup.~fr.~112~K.-A. εἰς Ἱπποκράτους τε παῖδες ἐκβόλιμοί τινες, / βληχητὰ τέκνα καὶ οὐδαμῶς τοῦ ⟨–⟩ τρόπου ; Eur.~fr.~698~Kannicht δεῖ γάρ με δόξαι πτωχὸν – × – ⏑ –, / εἶναι μὲν ὅσπερ εἰμί, φαίνεσθαι δὲ μή ; Phryn.~\textit{PS} p.~31.14~Bekker Γυναικηρὸς τρόπος: ἀντὶ τοῦ γυναικεῖος ; Greg. Cypr.~1.79 Ἄνθρωπος εὔριπος καὶ ψυχὴ καὶ διάνοια καὶ τρόπος: ἐπὶ τῶν ῥᾷστα μεταβαλλομένων ; Apostol.~12.42c Οἷος ὁ τρόπος, τοιοῦτος ὁ λόγος ; Ar.~fr.~488~K.-A. χρῶμαι γὰρ αὐτοῦ τοῦ στόματος τῷ στρογγύλῳ, / τοὺς νοῦς δ᾽ ἀγοραίους ἧττον ἢ ᾽κεῖνος ποιῶ ; Philostr.~\textit{VA}~8.7 ῥητορικώτερον ἴσως ἀπολελόγημαι τοὐμοῦ τρόπου ; Hsch.~α~6747 ἀπὸ τρόπου· ἔξω τοῦ τρόπου}
\lp{94 τοῦ γὰρ τεχνάζειν]}{Ar.~\textit{Ach.}~385 τί ταῦτα στρέφει τεχνάζεις τε καὶ πορίζεις τριβάς ; Ar.~\textit{Thesm.}~198 τὰς συμφορὰς γὰρ οὐχὶ τοῖς τεχνάσμασιν / φέρειν δίκαιον ἀλλὰ τοῖς παθήμασιν ; Ar.~\textit{Thesm.}~270–271 συσσώσειν ἐμὲ / πάσαις τέχναις, ἤν μοί τι περιπίπτῃ κακόν ; Ar.~\textit{Ran.}~80 κἄλλως ὁ μέν γ᾽ Εὐριπίδης πανοῦργος ὤν ; Ar.~\textit{Ran.}~956–958 νοεῖν ὁρᾶν ξυνιέναι στρέφειν ἐρᾶν τεχνάζειν, / κάχ᾽ ὑποτοπεῖσθαι, περινοεῖν ἅπαντα ; Synag.~κ~166 κατασοφισώμεθα· τεχνασώμεθα, μηχανῇ τινι κακώσωμεν ; Suid.~σ~812 Ἀριστοφάνης λέγει σοφίζῃ, ἀντὶ τοὺ τεχνάζῃ. σοφίας γὰρ καὶ σοφίσματα τὰς τέχνας ἔλεγον ; Zen.~5.68 Πόλλ᾽ οἶδ᾽ ἀλώπηξ, ἀλλ᾽ ἐχῖνος ἓν μέγα: μέμνηται ταύτης Ἀρχίλοχος ἐν ἐπωδῇ … Φησὶ δὲ καὶ Ἴων ὁ τραγικός· Ἀλλ᾽ ἔν τε χέρσῳ τὰς λέοντος ᾔνεσα καὶ τὰς ἐχίνου μᾶλλον οἰζυρὰς τέχνας … Λέγεται δὲ ἡ παροιμία ἐπὶ τῶν πανουργοτάτων ; Ar.~\textit{Thesm.}~65 Ἀγάθωνά μοι δεῦρ᾽ ἐκκάλεσον πάσῃ τέχνῃ ; Ar.~\textit{Ran.}~1520–1521 μέμνησο δ᾽ ὅπως ὁ πανοῦργος ἀνὴρ / καὶ ψευδολόγος καὶ βωμολόχος ; Eur.~\textit{IA}~745 τέχνας πορίζω, πανταχῇ νικώμενος ; Eur.~\textit{Or.}~1560 τεχνάσματ᾽ ἐστὶ ταῦτα καὶ πολὺς γέλως ; Hdt.~3.130.2 κατεφάνη τε τῷ Δαρείῳ τεχνάζειν ἐπιστάμενος ; Hdt.~6.1.2 ὁρέων αὐτὸν τεχνάζοντα ; X.~\textit{An.}~7.6.16 ὅτι ἔξεστι καὶ τὰ ὑμέτερα ἔχοντα παρὰ Σεύθου τεχνάζειν ; D.~40.54 ἐάν τι τοιοῦτον τεχνάζῃ ; Σ Ar.~\textit{Ran.}~106 ἀτεχνῶς: οἷον, ἀδόλως· τέχνας γὰρ ἔλεγον οἱονεὶ τοὺς δόλους ; Hsch.~ε~6288 ἐσοφισάμεθα· ἐτεχνασάμεθα ; Hsch.~κ~2314 κέρδος· σύνεσις. τέχνασμα. ὠφέλεια ; Hsch.~μ~1148 μήδεται· βουλεύεται. τεχνάζεται ; Hsch.~π~149 παλαμήσασθαι· τεχνάσασθαι. καὶ τὰ ὅμοια ; Ammon.~\textit{Diff.}~467 τεχνάσματα καὶ τεχνήματα διαφέρουσι. τεχνάσματα μὲν γὰρ τὰ κακουργήματα, τεχνήματα δὲ αἱ τέχναι ; Synag.~π~570 πόρος· πέρας ἢ ποταμὸς ἢ ὁδὸς ἢ φυσήματα ὑδάτων ἢ μηχανὴ ἢ τεχνάσματα ; Suid.~μ~428 Μεθοδεύει: μετέρχεται, τεχνάζεται, ἀπατᾷ ; Zen.~2.91 Γλαύκου τέχνη ; Diogenian.~3.99 Γύγου δακτύλιος: ἐπὶ τῶν πολυμηχάνων καὶ πανούργων ; Diogenian.~6.70a Μὴ πρὸς ἐμὲ τὰ ποικίλα: τουτέστι, μὴ τὰς τέχνας κατὰ τῶν φίλων ; Diogenian.~8.29 Τὸ Φρυνίχου πάλαισμα: κατὰ τῶν πανούργως σοφιζομένων. Προδότης γὰρ γέγονε τῶν Λακεδαιμονίων ; Ar.~\textit{Vesp.}~192 πονηρὸς εἶ πόρρω τέχνης καὶ παράβολος ; Ar.~fr.~717~K.-A. τί γὰρ ἐπὶ κακότροπον ἔμολε τότε βίον / ἀδικομηχάνῳ τέχνῃ ; Eur.~\textit{Hipp.}~917 τί δὴ τέχνας μὲν μυρίας διδάσκετε ; Soph.~\textit{Phil.}~80 τοιαῦτα φωνεῖν μηδὲ τεχνᾶσθαι κακά ; Agath.~fr.~6 Snell (TrGF~39 F~6) τέχνη τύχην ἔστερξε καὶ τύχη τέχνην ; Hsch.~ε~2578 ἐμώσατο· εὗρεν. ἐτεχνάσατο. ἐζήτησεν ; Hsch.~κ~1421 κατασοφισθεῖσα· χλευασθεῖσα. τεχνασθεῖσα ; Hsch.~μ~1265 μήσατο· εἰργάσατο, ἐτεχνάσατο ; Hsch.~μ~2074 μῶται· ζητεῖ. τεχνάζεται ; Hsch.~π~148 παλαμήσας· τεχνάσας. ἐργάσας ; Hsch.~τ~691 τευμᾶται· τεχνάζει [τιμᾶται] ; Hsch.~τ~705 τευτᾶ⟨ται⟩· τεχνᾶται ; Hsch.~τ~721 τεχνάζω· τεχνάζομαι ; Hsch.~τ~724 τεχνήσομαι· ἐπινοήσω. τεχνάσομαι ; Suid.~π~40 Παλαμᾶσθαι: τεχνάζεσθαι ; Zen.~4.82 ἀπὸ Κόρακος τοῦ Συρακουσίου ῥήτορος, πρώτου διδάξαντος τέχνην ῥητορικήν ; Diogenian.~5.92 Κρὴς πρὸς Αἰγινήτην: ἐπὶ τῶν πανουργίᾳ χρωμένων πρὸς ἀλλήλους λέγεται}
\lp{94 ἡμέτερος ὁ πυραμοῦς]}{\textit{Test.} Suid.~π~3192 Πυραμοῦς: εἶδος πλακοῦντος ἐκ μέλιτος ἑφθοῦ καὶ πυροῦ πεφρυγμένου, ὡς σησαμοῦς ὁ διὰ σησαμῆς. ταῦτα δὲ ἐτίθεσαν ἆθλα τοῖς διαγρυπνηταῖς ἐν τοῖς συμποσίοις, καὶ ὁ διαγρυπνήσας μέχρι τῆς ἕω ταῦτα ἐλάμβανεν ἔπαθλον. ἡμέτερος ὁ πυραμοῦς ; Ar.~\textit{Eq.}~276–277 ἀλλ᾽ ἐὰν μέντοι γε νικᾷς τῇ βοῇ, τήνελλος εἶ· / ἢν δ᾽ ἀναιδείᾳ παρέλθῃ σ᾽, ἡμέτερος ὁ πυραμοῦς ; \textit{idem} Suid.~τ~518 (παρέλθῃς ; add. ὁ δὲ πυραμοῦς εἶδος πλακοῦντος) ; Ar.~\textit{Eq.}~1253 Ἑλλάνιε Ζεῦ σὸν τὸ νικητήριον ; Σ Ar.~\textit{Eq.}~277 Κἂν ἡττηθῇς, φησίν, ὑπ᾽ αὐτοῦ τῇ βοῇ, περιγένῃ δὲ ἐν ἀναιδείᾳ, οὐδὲν ἧττον ἡμετέρα ἡ νίκη. πυραμοῦς δέ, εἶδος πλακοῦντος ἐκ μέλιτος ἑφθοῦ καὶ πυρῶν πεφρυγμένων, ὡς καὶ σησαμοῦς τὸ διὰ σησάμων. ταῦτα δὲ ἐτίθεσαν ἆθλα τοῖς διαγρυπνηταῖς. εἰώθασι γὰρ ἐν τοῖς συμποσίοις ἁμιλλᾶσθαι περὶ ἀγρυπνίας, καὶ ὁ διαγρυπνήσας μέχρι τῆς ἕω ἐλάμβανε τὸν πυραμοῦντα ; Σ Ar.~\textit{Thesm.}~94 εἶδος πλακοῦντος ὁ πυραμοῦς. οὗτος δὲ ἐτίθετο τοῖς διαπαννυχίζουσιν ἔπαθλον, καὶ ὁ ἐγρηγορὼς ἕως πρωΐας ὡς νικῶν αὐτὸν ἐλάμβανε. νῦν οὖν ταύτῃ τῇ μεταφορᾷ ἐχρήσατο ὡς νικῶντος αὐτοῦ τῇ πανουργίᾳ πάντας, ἤτοι τὰς γυναῖκας ; Bauck, \textit{De\,proverbiis\,Ar.,} p.~66 ἡμέτερος ὁ πυραμοῦς ; Hsch.~α~7714 ἀσκὸν λήψεται· ἀσκὸς νικητήριον ἐτίθετο, καὶ μετὰ σάλπιγγος ἔπινον. νικητήριον στέφανον ; Ar.~\textit{Ach.}~1000–1002 κατὰ τὰ πάτρια τοὺς Χοᾶς / πίνειν ὑπὸ τῆς σάλπιγγος· ὃς δ᾽ ἂν ἐκπίῃ / πρώτιστος, ἀσκὸν Κτησιφῶντος λήψεται ; Hsch.~ν~573 νικητήριον· ἔπαθλον ; Hsch.~τ~1427 τρὶς ἕξ· νικητήριος βόλος. καὶ παροιμία· τρὶς ἓξ ἢ τρεῖς κύβοι. ἐλέγετο δὲ ἐπὶ τοῦ εὐστόχου, ἀπὸ τῶν κυβευτῶν}
\lp{94 πυραμοῦς]}{Call.~fr.~227.5–7~Pf. (ap. Ath.~15.668c) ὁ δ᾽ ἀγρυπνήσας … μέχρι τῆς κορώνης· τὸν πυραμοῦντα λήψεται καὶ τὰ κοττάβεια καὶ τῶν παρουσῶν ἣν θέλει χὂν θέλει φιλήσει ; Ath.~14.647b–c καὶ τοῦ πυραμοῦντος καλουμένου, οὐ διαφέρειν λέγων τῆς πυραμίδος καλουμένης· γίνεσθαι γὰρ ταύτην ἐκ πυρῶν πεφωσμένων καὶ μέλιτι δεδευμένων. αὗται δὲ ἆθλα τίθενται ταῖς παννυχίσι τῷ διαγρυπνήσαντι ; Artem.~1.72 κοπταὶ δὲ σησαμίδες καὶ πυραμοὶ ἀγαθὰ πᾶσι, μάλιστα δὲ τοῖς δίκην ἔχουσιν· ἦν γὰρ ὁ πυραμὸς παρὰ τοῖς παλαιοῖς ἐπινίκιος ; Hsch.~π~4391 πυραμοῦς· εἶδος πλακοῦντος, ἐκ πυρῶν πεφρυγμένων καὶ μέλιτι ἀναδεδευμένων ; Hsch.~σ~533 σησαμόεντ᾽ ἄρτον· οὗτος ἔκειτο ἆθλον τοῖς διαγρυπνήσασι, πυραμοῦς καὶ σησάμινος καὶ τοιαῦτά τινα ; Poll.~6.108 τοῖς μέντοι διαπαννυχίσασιν ἆθλα ἦν σησαμοῦς καὶ πυραμοῦς, πέμματα διὰ μέλιτος ἑφθά, τὸ μὲν ἐκ σησάμων, τὸ δ᾽ ἐκ πυρῶν πεφρυγμένων ; EM~533.21–23 καὶ ἆθλα οἱ νικῶντες ἐλάμβανον πλακουντίσκους, πυραμοῦντας ἢ σησαμοῦντας, ἅπερ κοττάβια ἔλεγον ; Ar.~\textit{Thesm.}~204–205 δοκῶν γυναικῶν ἔργα νυκτερείσια / κλέπτειν ὑφαρπάζειν τε θήλειαν Κύπριν ; Ephipp.~fr.~8.3~K.-A. ἴτρια, τραγήμαθ᾽ ἧκε, πυραμοῦς, ἄμης ; Plaut.~\textit{Aul.}~36 qui illam stupravit noctu, Cereris vigiliis ; Eubul.~fr.~1~K.-A. καὶ γὰρ πάλαι πέττει τὰ νικητήρια … ἐξεπήδησ᾽ ἀρτίως πέττουσα τὸν χαρίσιον ; Pl.~\textit{Smp.}~223c–d Ἀγάθωνα δὲ καὶ Ἀριστοφάνη καὶ Σωκράτη ἔτι μόνους ἐγρηγορέναι ; Plu.~\textit{Quaest.\,conv.}~747a ἐκ τούτου πυραμοῦντος … νικητήριον ὀρχήσεως ; Hsch.~π~365 παννυχίς· ἑορτὴ νυκτερινή ; Hsch.~π~4388 πυραμοῦντα· τὴν πυραμίδα ; Poll.~6.111 ὁ δὲ πλεῖστα καταδύσας τῶν ἐπιπολαζόντων τὰ κοτταβεῖα τὸ ἆθλον ἐλάμβανεν· ἐκαλεῖτο γὰρ οὐ τὰ ἀγγεῖα κοτταβεῖα μόνον, ἀλλὰ καὶ τὸ ἆθλον, ἦν δὲ πυραμοῦς καὶ σησαμοῦς καὶ ἄλλα τοιαῦτα ; Ephipp.~fr.~13.3–5~K.-A. μέλι, σησαμίδες, / †βράχος, βρυγμός†, μνοῦς, πυραμίδες, / μῆλον, κάρυον ; Ath.~3.114b ἐστὶ δὲ καὶ ὁ πυραμοῦς ἄρτος διὰ σησάμων πεττόμενος ; Hsch.~σ~529 σησαμῆ· σησαμίς. καὶ πλακοῦς ἐκ σησάμης ; Hsch.~σ~535 σησαμοῦς· πέμμα ἐκ μέλιτος καὶ σησάμης ; Hsch.~χ~199 χαρίσιον· εἶδος πλακούντων (Ar. fr. 211, 2 K.–A.). οἱ δὲ ἄρτον ἀπὸ τῶν λειμμάτων (Eubul. fr. 1, 3 K.–A.) ; Poll.~6.77 τὰ δὲ πλακούντων εἴδη ἄμης ἀμητίσκος, πυραμοῦς, σησαμοῦς, ἔγχυτος}
\lp{95 σίγα]}{Ar.~\textit{Thesm.}~27 σίγα νυν. — σιωπῶ τὸ θύριον ; Ar.~\textit{Thesm.}~39–40 εὔφημος πᾶς ἔστω λαός, / στόμα συγκλῄσας ; Ar.~\textit{Thesm.}~45 βομβάξ. — σίγα. — τί λέγει ; Cratin.~fr.~315~K.-A. ἄκουε, σίγα, πρόσεχε τὸν νοῦν, δεῦρ᾽ ὅρα ; Eur.~\textit{Hec.}~532–533 Σιγᾶτ᾽, Ἀχαιοί, σῖγα πᾶς ἔστω λεώς, / σίγα σιώπα ; Eur.~\textit{Hec.}~1070a σίγα· κρυπτὰν βάσιν αἰσθάνομαι ; Eur.~\textit{Hipp.}~565 σιγήσατ᾽, ὦ γυναῖκες· ἐξειργάσμεθα ; [Eur.]~\textit{Rh.}~730 σῖγα πᾶς ὕφιζ᾽· ἴσως γὰρ ἐς βόλον τις ἔρχεται ; Soph.~\textit{Aj.}~87 σίγα νυν ἑστὼς καὶ μέν᾽ ὡς κυρεῖς ἔχων ; Hor.~\textit{Carm.}~3.1.1–4 Odi profanum volgus et arceo. / favete linguis: carmina non prius / audita Musarum sacerdos / virginibus puerisque canto ; \textit{idem} Tosi n°~2083 (Favete linguis!) ; Hsch.~σ~573 σιγάζοντος· Ξενοφῶν ἐν Ἀναβάσει ; Diogenian.~3.43 Βάχχης τρόπον: ἐπὶ τῶν σιωπηλῶν. Παρόσον αἱ Βάκχαι σιγῶσι ; Greg. Cypr.~3.79 Τοῦ Κροίσου παιδὸς σιγηλότερος ; App. Prov.~3.7 Ἢ λέγε τι σιγῆς κρεῖττον, ἢ σιγὴν ἔχε: ἐπὶ τῶν οὐκ ἄξια λεγόντων ; Ar.~\textit{Pax}~91 σίγα σίγα ; Ar.~\textit{Thesm.}~381 σίγα σιώπα, πρόσεχε τὸν νοῦν· χρέμπτεται γὰρ ἤδη ; Eur.~\textit{Alc.}~703 σίγα· νόμιζε δ᾽, εἰ σὺ τὴν σαυτοῦ φιλεῖς ; Eur.~\textit{HF}~1068–1069 σῖγα σῖγα· / παλίντροπος ἐξεγειρόμενος στρέφεται ; Eur.~\textit{Or.}~140 σῖγα σῖγα, λεπτὸν ἴχνος ἀρβύλης ; \textit{idem} D.L.~7.172 (add. τὸ τῆς Ἠλέκτρας, ἔφη ; om. ἀρβύλης) ; Soph.~\textit{Trach.}~975 σίγα, τέκνον, μὴ κινήσῃς ; Luc.~\textit{Deor.\,conc.}~1 σὺ δὲ κήρυττε, ὦ Ἑρμῆ, τὸ κήρυγμα τὸ ἐκ τοῦ νόμου. «ἄκουε, σίγα» ; Σ Eur.~\textit{Hec.}~533 μεμίμηται τὸ πολλάκις κηρύττειν τὰ αὐτὰ τοὺς κήρυκας ; Hsch.~σ~473 σιγᾷ· σιωπᾷ ; Hsch.~σ~580 σιγηλός· ἄφωνος ; Hsch.~σ~581 σιγημονᾷς· σιγᾷς ; Hsch.~σ~582 σιγηρῶς· σιγηλῶς, ἡσύχως ; Synag.~α~762 ἄκουε· ἀντὶ τοῦ σίγα. Κρατῖνος (fr. 315). οἷον τόπον δίδου διὰ τῆς σιγῆς τῇ ἀκοῇ}
\lp{95 σίγα. — τί δ᾽ ἔστιν;]}{Eur.~\textit{Hipp.}~566 τί δ᾽ ἔστι, Φαίδρα, δεινὸν ἐν δόμοισί σοι ; Plaut.~\textit{Pseud.}~207 Vah tace. — Quid est? ; Soph.~fr.~815~R. ἄκουε, σίγα· τίς ποτ᾽ ἐν δόμοις βοή}
\lp{95 σίγα … Ἁγάθων ἐξέρχεται]}{Ar.~\textit{Ach.}~238–240 σῖγα πᾶς. ἠκούσατ᾽ ἄνδρες ἆρα τῆς εὐφημίας; / οὗτος αὐτός ἐστιν ὃν ζητοῦμεν. ἀλλὰ δεῦρο πᾶς / ἐκποδών· θύσων γὰρ ἁνὴρ ὡς ἔοικ᾽ ἐξέρχεται ; Eur.~\textit{Cyc.}~487–491 σίγα σίγα. καὶ δὴ μεθύων / ἄχαριν κέλαδον μουσιζόμενος / σκαιὸς ἀπῳδὸς καὶ κλαυσόμενος / χωρεῖ πετρίνων ἔξω μελάθρων ; Eur.~\textit{Or.}~1366–1367 [σιγήσατ᾽· ἔξω γάρ τις ἐκβαίνει Φρυγῶν,] ; Soph.~\textit{OC}~111 σίγα· πορεύονται γὰρ οἵδε δή τινες}
\lp{95 τί δ᾽ ἔστιν;]}{Ar.~\textit{Ach.}~189 αἰβοῖ. — τί ἔστιν ; Ar.~\textit{Lys.}~830 τί δ᾽ ἔστιν; εἰπέ μοι τίς ἡ βοή ; Ar.~\textit{Thesm.}~193 Εὐριπίδη— τί ἔστιν ; Ar.~\textit{Thesm.}~582 τί δ᾽ ἔστιν ὦ παῖ; παῖδα γάρ σ᾽ εἰκὸς καλεῖν ; Lysipp.~fr.~1.1~K.-A. (Α.) Ἕρμων. (Ερ.) τί ἔστι; (Α.) πῶς ἔχομεν ; Ar.~\textit{Plu.}~1107 τί δ᾽ ἔστιν; — ὁ Ζεὺς ὦ πόνηρε βούλεται ; Philem.~fr.~117.1~K.-A. (Α.) Σύρα, Σύρα. (Συ.) τί ἐστι; (Α.) πῶς ἡμῖν ἔχεις ; Soph.~\textit{Aj.}~897 τί δ᾽ ἔστιν}
\lp{95 τί δ᾽ ἔστιν; — Ἁγάθων ἐξέρχεται]}{Ar.~\textit{Vesp.}~137 οἴμοι. — τί ἔστι; — Βδελυκλέων ἀνίσταται}
\lp{95 Ἁγάθων ἐξέρχεται]}{Ar.~\textit{Eq.}~234 οἴμοι κακοδαίμων ὁ Παφλαγὼν ἐξέρχεται ; Ar.~\textit{Thesm.}~36–37 ἀλλ᾽ ἐκποδὼν πτήξωμεν, ὡς ἐξέρχεται / θεράπων τις αὐτοῦ πῦρ ἔχων καὶ μυρρίνας ; Ar.~\textit{Thesm.}~70 περίμεν᾽, ὡς ἐξερχεται ; Ar.~\textit{Lys.}~5 πλὴν ἥ γ᾽ ἐμὴ κωμῆτις ἥδ᾽ ἐξέρχεται ; Ar.~\textit{Lys.}~1107 αὐτὴ γάρ, ὡς ἤκουσεν, ἥδ᾽ ἐξέρχεται ; Hsch.~ε~3665 ἔξεισιν· ἐξέρχεται (Eur. Alc. 215)}
\lp{96 καὶ ποῖός ἐστιν;]}{Ar.~\textit{Thesm.}~30–33 ποῖος οὗτος Ἁγάθων; — ἔστιν τις Ἀγάθων — μῶν ὁ μέλας ὁ καρτερός; … μῶν ὁ δασυπώγων ; Eur.~\textit{Ba.}~501–502 καὶ ποῦ ᾽στιν; οὐ γὰρ φανερὸς ὄμμασίν γ᾽ ἐμοῖς. / — παρ᾽ ἐμοί· σὺ δ᾽ ἀσεβὴς αὐτὸς ὢν οὐκ εἰσορᾷς ; Ar.~\textit{Ach.}~963 ὁ ποῖος οὗτος Λάμαχος τὴν ἔγχελυν ; Ar.~\textit{Thesm.}~645 καὶ ποῦ ᾽στιν; — αὖθις ἐς τὸ πρόσθεν οἴχεται ; Ar.~\textit{Ran.}~1123 καὶ ποῖον αὐτοῦ βασανιεῖς; — πολλοὺς πάνυ ; Ar.~\textit{Thesm.}~78 καὶ πῶς; ἐπεὶ νῦν γ᾽ οὔτε τὰ δικαστήρια / μέλλει δικάζειν}
\lp{96 οὑκκυκλούμενος]}{\textit{Test.} Suid.~ε~699 Ἐκκυκλεῖ· ἐκκαλύπτει. καὶ Ἐκκυκλούμενος, ἐκκύκλημα ἔχων ; \textit{idem} Zonar.~678~Tittmann ; Ar.~\textit{Ach.}~408–409 ἀλλ᾽ ἐκκυκλήθητ᾽. — ἀλλ᾽ ἀδύνατον. — ἀλλ᾽ ὅμως. / ἀλλ᾽ ἐκκυκλήσομαι· καταβαίνειν δ᾽ οὐ σχολή ; Ar.~\textit{Thesm.}~265 εἴσω τις ὡς τάχιστά μ᾽ ἐσκυκλησάτω ; Σ Ar.~\textit{Ach.}~408 ἐκκύκλημα δὲ λέγεται μηχάνημα ξύλινον τροχοὺς ἔχον, ὅπερ περιστρεφόμενον τὰ δοκοῦντα ἔνδον ὡς ἐν οἰκίᾳ πράττεσθαι καὶ τοῖς ἔξω ἐδείκνυε, λέγω δὴ τοῖς θεαταῖς ; \textit{idem} Suid.~ε~132 (Ἐγκυκλήθητι· … ἐγκύκλημα ; om. ἔξω … λέγω δὴ ; ἐδείκνυεν) ; Σ Ar.~\textit{Thesm.}~96 ἐπὶ ἐκκυκλήματος γὰρ φαίνεται ; Σ Ar.~\textit{Thesm.}~277 παρεπιγραφή. ἐκκυκλεῖται ἐπὶ τὸ ἔξω τὸ θεσμοφόριον ; Σ Il.~18.477 δαιμονίως τὸν πλάστην αὐτὸς διέπλασεν, ὥσπερ ἐπὶ σκηνῆς ἐκκυκλήσας καὶ δείξας ἡμῖν ἐν φανερῷ τὸ ἐργαστήριον ; Poll.~4.128 [καὶ] τὸ μὲν ἐκκύκλημα ἐπὶ ξύλων [ὑψηλὸν] βάθρον, ᾧ ἐπίκειται θρόνος· δείκνυσι δὲ τὰ ὑπὸ σκηνὴν ἐν ταῖς οἰκίαις ἀπόρρητα πραχθέντα. καὶ τὸ ῥῆμα τοῦ ἔργου καλεῖται ἐκκυκλεῖν ; Ar.~\textit{Ach.}~479 ἁνὴρ ὑβρίζει· κλῇε πηκτὰ δωμάτων ; Ar.~\textit{Eq.}~1249 κυλίνδετ᾽ εἴσω τόνδε τὸν δυσδαίμονα ; Max.Tyr.~17.5 εἰσκυκλεῖς ποιητῶν χορὸν ἀκόλαστον καὶ βακχεύοντα ; Philostr.~\textit{VA}~8.7 καὶ τὸ μὲν ἱπποτροφεῖν αὐτοὺς κἀπὶ ζευγῶν ἐς τὴν ἀγορὰν ἐκκυκλεῖσθαι λευκῶν ; Clem.Al.~\textit{Protr.}~2 οἷον ἐπὶ σκηνῆς τοῦ βίου τοῖς τῆς ἀληθείας ἐκκυκλήσω θεαταῖς ; Clem.Al.~\textit{Protr.}~7 ἤδη δὲ ἐν Ἴωνι τῷ δράματι γυμνῇ τῇ κεφαλῇ ἐκκυκλεῖ τῷ θεάτρῳ τοὺς θεούς ; Σ Ar.~\textit{Nub.}~184 [ὁρᾷ δὲ ὡς φιλοσόφους κομῶντας, στραφέντος τοῦ ἐγκυκλήματος.] ; Σ Soph.~\textit{Aj.}~346 ἐνταῦθα ἐκκύκλημά τι γίνεται ἵνα φανῇ ἐν μέσοις ὁ Αἴας ποιμνίοις ; Σ Eur.~\textit{Hipp.}~171 τοῦτο σεσημείωκεν Ἀριστοφάνης, ὅτι κατὰ τὸ ἀκριβὲς τὸ ἐκκύκλημα τοιοῦτον ἐστὶ τῇ ὑποθέσει. ἐπὶ γὰρ τῆς σκηνῆς δείκνυται τὰ ἔνδον πραττόμενα ; Hsch.~ε~1465 ἐκκυκλεῖ· ἐκκαλύπτει ; \textit{idem} Synag.~ε~204 ; Hsch.~ε~4014 ἐξώστρα· ἐπὶ τῆς σκηνῆς τὸ ἐκκύκλημα (Phryn. trag. 3 F 21a, com. fr. 90 K.–A.) ; AB~219.28 Βῆμα: πλείονα σημαίνει κοινότερον. ἔστι δὲ οὕτω πως καὶ τὸ λογεῖον, ὥσπερ καὶ τὸ ἐν ἐκκλησίᾳ καὶ τὸ ἐν δικαστηρίοις. ἴσως δὲ καὶ τὸ ἐκκύκλημα ; Plu.~\textit{Prof.\,virt.}~80a ὃ λαβόντες ἐκ φιλοσοφίας εὐθὺς εἰς ἀγορὰν ἢ νέων διατριβὴν ἢ βασιλικὸν συμπόσιον ἐκκυκλήσουσιν ; Luc.~\textit{Philops.}~29 θεὸν ἀπὸ μηχανῆς ἐπεισκυκληθῆναί μοι τοῦτον ᾤμην ὑπὸ τῆς Τύχης ; Philostr.~\textit{VA}~6.11 φιλοσοφίας … εὖ κεκοσμημένης, ἣν ἐς τὸ πρόσφορον Ἰνδοὶ στείλαντες ἐφ᾽ ὑψηλῆς τε καὶ θείας μηχανῆς ἐκκυκλοῦσιν ; Ath.~6.270e πλῆθος ἰχθύων καὶ ἄλλων παντοδαπῶν ὄψων παρασκευὴν εἰσκυκλουμένην ; Clem.Al.~\textit{Strom.}~3.4.26 ἡ εἰς μέσον τῶν ἀποστόλων τῆς ζηλοτυπουμένης ἐκκύκλησις γυναικός ; Clem.Al.~\textit{Strom.}~7.14.88 οὐ γὰρ ἐκκυκλεῖν χρὴ τὸ μυστήριον ; Plb.~11.5.8 τῆς τύχης ὥσπερ ἐπίτηδες ἐπὶ τὴν ἐξώστραν ἀναβιβαζούσης τὴν ὑμετέραν ἄγνοιαν ; Hsch.~β~563 βῆμα· … σημαίνει δὲ καὶ τὸ κύκλημα (γράφεται δὲ καὶ ἐκκύκλημα)}
\lp{97 ἀλλ᾽ ἦ]}{Ar.~\textit{Ach.}~424 ἀλλ᾽ ἦ Φιλοκτήτου τὰ τοῦ πτωχοῦ λέγεις ; Eur.~\textit{Alc.}~58 πῶς εἶπας; ἀλλ᾽ ἦ καὶ σοφὸς λέληθας ὤν ; Eur.~\textit{Ba.}~922 ἀλλ᾽ ἦ ποτ᾽ ἦσθα θήρ; τεταύρωσαι γὰρ οὖν ; Ar.~\textit{Vesp.}~8 ἀλλ᾽ ἦ παραφρονεῖς ἐτεὸν ἢ κορυβαντιᾷς ; Ar.~\textit{Lys.}~928 ἀλλ᾽ ἢ τὸ πέος τόδ᾽ Ἡρακλῆς ξενίζεται ; Eur.~\textit{Hipp.}~858–859 ἀλλ᾽ ἦ λέχους μοι καὶ τέκνων ἐπιστολὰς / ἔγραψεν ἡ δύστηνος, ἐξαιτουμένη ; Soph.~\textit{Phil.}~414 πῶς εἶπας; ἀλλ᾽ ἦ χοὖτος οἴχεται θανών ; Eur.~\textit{Alc.}~816 ἀλλ᾽ ἦ πέπονθα δείν᾽ ὑπὸ ξένων ἐμῶν}
\lp{97 ἀλλ᾽ ἦ τυφλὸς μέν εἰμ᾽]}{Ar.~\textit{Eq.}~953 τὸ γοῦν σημεῖον ἕτερον φαίνεται, / ἀλλ᾽ ἢ οὐ καθορῶ}
\lp{97 τυφλός]}{Ar.~fr.~33~K.-A. τυφλότερος λεβηρίδος ; Ar.~\textit{Plu.}~48–49 δῆλον ὁτιὴ καὶ τυφλῷ / γνῶναι δοκεῖ τοῦθ᾽ ; Men.~fr.~108~K.-A. (Σ \textit{Pl.\,Soph.}~241d) καὶ τυφλῷ δῆλον … μέμνηται ταύτης καὶ Μένανδρος ἐν Δεισιδαίμονι ; \textit{idem} Macar.~3.29 (Δῆλόν ἐστι καὶ τυφλῷ· ἐπὶ τῶν προφανῶν), Macar.~4.91 (Καὶ τυφλῷ δῆλον· ἐπὶ τῶν ἄγαν σαφεστάτων) ; Pl.~\textit{R.}~550d καὶ τυφλῷ γε δῆλον ὡς μεταβαίνει ; Plb.~18.4.4 τοῦτο μὲν ἔφησεν ὦ Φαινέα, καὶ τυφλῷ δῆλον ; Liv.~32.34.3 ᾽apparet id quidem᾽ inquit Philippus ᾽etiam caeco᾽, iocatus in valetudinem oculorum Phaeneae ; \textit{idem} Tosi n°~1966 (Apparet id quidem… etiam caeco) ; Hsch.~γ~1003 γυμνότερος λεβηρίδος· Ἀριστοφάνης (fr. 33 K.–A.) φησὶ τυφλότερος λεβηρίδος. ἔστι δὲ λεβηρὶς τὸ τοῦ ὄφεως γῆρας, διὰ τὸ λέπος εἶναι. κενὸν δὲ τελέως ἐστὶ τοῦτο καὶ τυφλόν. τὰς γὰρ ὀπὰς μονὰς ἔχει τῶν ὀφθαλμῶν ; \textit{idem} Zen.~2.95 (τέλεως ἐστὶ δὲ καὶ τοῦτο τυφλόν ; μόνας ; add. Τριχῇ δὲ ἀναγράφουσι τὴν παροιμίαν, καὶ οἱ μὲν τυφλότερος λεβηρίδος, οἱ δὲ κενότερος, οἱ δὲ γυμνότερος) ; Soph.~\textit{OT}~371 τυφλὸς τά τ᾽ ὦτα τόν τε νοῦν τά τ᾽ ὄμματ᾽ εἶ ; \textit{Od.}~8.195 καί κ᾽ ἀλαός τοι, ξεῖνε, διακρίνειε τὸ σῆμα ; Pl.~\textit{R.}~465c δῆλά τε δὴ καὶ ἀγεννῆ καὶ οὐκ ἄξια λέγειν. δῆλα γάρ, ἔφη, καὶ τυφλῷ ; Pl.~\textit{Sph.}~241d πῶς γὰρ οὐ φαίνεται καὶ τὸ λεγόμενον δὴ τοῦτο τυφλῷ ; Plu.~\textit{Quaest.\,conv.}~633c καί ποτε λαβὼν ἀξίωμα μεγάλοις γράμμασι γεγραμμένον ταυτὶ μέν ἔφη καὶ τυφλῷ δῆλα ; Demetr.~\textit{Eloc.}~112 Τὸ δὲ ποιητικὸν ἐν λόγοις ὅτι μὲν μεγαλοπρεπές, καὶ τυφλῷ δῆλόν φασι ; Hsch.~τ~1700 τυφλός· τίθεται καὶ ἀντὶ τοῦ κωφός ; Diogenian.~8.25 Τυφλότερος ἀσπάλακος: φασὶν αὐτὸν τυφλωθῆναι διὰ τὸ φθείρειν τὴν γῆν ; \textit{idem} Hsch.~τ~1701 (σπάλακος ; τοῦτο τὸ ζῷον οὐκ ἔχει ὀφθαλμούς pro φασὶν … γῆν) ; Plu.~\textit{Boiss.}~27 Τυφλῷ κάτοπτρον χαρίζῃ ; App. Prov.~5.12 Φαλακρὸς κτένα, Εὐνοῦχος παλλακήν, Κωφὸς αὐλητήν, Κάτοπρον ὁ τύφλος, Ὁ ἠπειρώπης κώπην, Ὁ κυβερνήτης ἄροτρον: πάντα ὅμοια ταῦτα. Λέγεται ἐπὶ τῶν ἀχρήστως τι ἑαυτοῖς ἐμπορευομένων ; Mant. Prov.~2.100 Τίς γὰρ κατόπτρῳ καὶ τυφλῷ κοινωνία ; \textit{idem} Tosi n°~1900 ; Eust.~1587.48–49 Ὅτι οὐδὲν μὲν θεάτρου μέτον τυφλῷ ἀνδρὶ ὥσπερ οὐδὲ κατόπτρου κατὰ τὴν παροιμίαν ; Erasm.~\textit{Adag.}~1.8.93 Vel caeco appareat. Proverbialis hyperbole frequens apud auctores, Vel caeco appareat, de re vehementer perspicua ; Plu.~\textit{Non\,posse}~1098f ὁ γὰρ οὕτω διαθεὶς ἑτέρους, ὥστε καὶ χαίρειν καὶ γάνυσθαι καὶ ποθεῖν ἅψασθαι καὶ προσαγορεῦσαι, δῆλός ἐστι καὶ τυφλῷ μεγάλας ἔχων ἐν ἑαυτῷ καὶ καρπούμενος ἡδονάς ; Luc.~\textit{Ind.}~19 οὐ μᾶλλον ἢ φαλακρὸς ἄν τις πρίαιτο κτένας ἢ κάτοπτρον ὁ τυφλὸς ἢ ὁ κωφὸς αὐλητὴν ἢ παλλακὴν ὁ εὐνοῦχος ἢ ὁ ἠπειρώτης κώπην ἢ ὁ κυβερνήτης ἄροτρον ; Euseb.~\textit{PE}~4.21.1 ἀλλὰ γὰρ καὶ τυφλῷ, φασί, δῆλον ; Quint.~12.7.9 caecis hoc, ut aiunt, satis clarum est ; Plu.~\textit{Boiss.}~44 Τυφλῷ διανεύεις}
\lp{97 τυφλὸς μέν εἰμ᾽]}{Ar.~\textit{Av.}~1214 ὑγιαίνεις μέν ; Eur.~\textit{Ion}~520 εὖ φρονεῖς μέν; ἤ σ᾽ ἔμηνε θεοῦ τις, ὦ ξένε, βλάβη ; Eur.~\textit{Med.}~676 θέμις μὲν ἡμᾶς χρησμὸν εἰδέναι θεοῦ ; Eur.~\textit{Alc.}~146 ἐλπὶς μὲν οὐκέτ᾽ ἐστὶ σῴζεσθαι βίον ; Eur.~\textit{Hipp.}~316 ἁγνὰς μέν, ὦ παῖ, χεῖρας αἵματος φορεῖς ; Eur.~\textit{Med.}~1129 τί φῄς; φρονεῖς μὲν ὀρθὰ κοὐ μαίνῃ, γύναι ; Pl.~\textit{Chrm.}~153c παρεγένου μέν, ἦ δ᾽ ὅς, τῇ μάχῃ ; Pl.~\textit{Men.}~82b Ἕλλην μέν ἐστι καὶ ἑλληνίζει ; Luc.~\textit{DMar.}~2.2 καὶ ἀπ᾽ ἐκείνου τυφλός εἰμί σοι, ὦ Πόσειδον}
\lp{97 τυφλὸς μέν εἰμ᾽ … ἐγὼ γὰρ οὐχ ὁρῶ]}{Luc.~\textit{Vit.\,auct.}~18 οὐχ ὁρῶ ταῦθ᾽ ἅπερ λέγεις τὰ παραδείγματα. εἰκότως· τυφλὸς γὰρ εἶ τῆς ψυχῆς τὸν ὀφθαλμόν}
\lp{97 ἐγὼ γὰρ οὐχ ὁρῶ]}{Ar.~\textit{Nub.}~208 ἐπεὶ δικαστὰς οὐχ ὁρῶ καθημένους ; Ar.~\textit{Thesm.}~19 διὰ τὴν χοάνην οὖν μήτ᾽ ἀκούω μήθ᾽ ὁρῶ ; D.~21.158 ἐγὼ μὲν γὰρ οὐχ ὁρῶ}
\lp{98 ἄνδρ᾽ οὐδέν᾽]}{Ar.~\textit{Thesm.}~141–143 τίς δ᾽ αὐτὸς ὦ παῖ; πότερον ὡς ἀνὴρ τρέφει; / καὶ ποῦ πέος; ποῦ χλαῖνα; ποῦ Λακωνικαί; / ἀλλ᾽ ὡς γυνὴ δῆτ᾽· εἶτα ποῦ τὰ τιτθία ; Pl.~\textit{Smp.}~193b–c καὶ μή μοι ὑπολάβῃ Ἐρυξίμαχος, κωμῳδῶν τὸν λόγον, ὡς Παυσανίαν καὶ Ἀγάθωνα λέγω—ἴσως μὲν γὰρ καὶ οὗτοι τούτων τυγχάνουσιν ὄντες καί εἰσιν ἀμφότεροι τὴν φύσιν ἄρρενες ; D.L.~6.27 ἄνδρας μέν, εἶπεν, οὐδαμοῦ, παῖδας δ᾽ ἐν Λακεδαίμονι ; D.L.~6.41 λύχνον μεθ᾽ ἡμέραν ἅψας περιῄει λέγων ἄνθρωπον ζητῶ ; Phryn.~\textit{PS} p.~31.22~Bekker Γυναικίζειν: γυναικῶν τρόπῳ διάγειν. καὶ γυναικίζοντες οἱ γυναικιζόμενοι ; Diocl.~fr.~4~K.-A. Γυναικισμός: Διοκλῆς Βάκχαις γυναικίζειν φησὶ καὶ γυναικηρόν ; Hdt.~7.210.2 ὅτι πολλοὶ μὲν ἄνθρωποι εἶεν, ὀλίγοι δὲ ἄνδρες ; D.L.~6.32 ἀνθρώπους ἐκάλεσα, οὐ καθάρματα ; Hsch.~γ~1015 γύννις· δειλός, ἄνανδρος, γυναικώδης· Ἀριστοφάνης (Thesm. 136). γύννιδας· μαλακούς ; Phryn.~\textit{PS} p.~19.4~Bekker Ἀῤῥενοῦσθαι: τὰ ἀῤῥένων, ἀλλὰ μὴ τὰ γυναικῶν πράττειν ; Phot.~α~1764 ἀνδρογύνων ἄθυρμα· Εὔπολις ἐν Ἀστρατεύτοις (fr. 3 Dem.). ἀνδρογύνους δὲ ἔλεγον τοὺς ἄνδρας μὲν τὸ σῶμα φύντας, εἰς γυναῖκας δὲ σφᾶς αὐτοὺς ἀφέντας καὶ τὰς τούτων ἐπιτηδεύσεις ἐπιτηδεύοντας ; cf. α 1763 ἀνδρογύνων· ἀσθενῶν, γυναικῶν καρδίας ἐχόντων ; App. Prov.~4.81 Τὰ Κινησίου δρᾷ: ἐπὶ τῶν μαλακῶν ; Hsch.~α~4745 ἀνδρόγυνος· ὁ ἑρμαφρόδιτος (Hdt. 4, 67, 2). καὶ ὁ ἀσθενής ἤγουν ὁ ἀνίσχυρος (Prov. 18, 8) ; Hsch.~θ~491 θηλυδριῶδες· τὸ κατακεκλασμένον (Ar. Thesm. 131)}
\lp{98 ἄνδρ᾽ οὐδέν᾽ … Κυρήνην δ᾽ ὁρῶ]}{Ar.~\textit{Eccl.}~166–168 γυναῖκας αὖ δύστηνε τοὺς ἄνδρας λέγεις; / — δι᾽ Ἐπίγονόν γ᾽ ἐκεῖνον· ἐπιβλέψασα γὰρ / ἐκεῖσε πρὸς γυναῖκας ᾠόμην λέγειν}
\lp{98 Κυρήνην]}{Ar.~\textit{Ran.}~1325–1328 τοιαυτὶ μέντοι σὺ ποιῶν / τολμᾷς τἀμὰ μέλη ψέγειν, / ἀνὰ τὸ δωδεκαμήχανον / Κυρήνης μελοποιῶν ; Pl. Com. fr.~143~K.-A. Ξενοκλῆς ὁ δωδεκαμήχανος, / ὁ Καρκίνου παῖς τοῦ θαλαττίου ; Eur.~fr.~765b~Kannicht ἀνὰ τὸ δωδεκαμήχανον ἄστρον ; Suid.~π~253 Πάξαμος, λόγιος. Ὀψαρτυτικὰ κατὰ στοιχεῖον, Βοιωτικὰ ἐν βιβλίοις β΄, Δωδεκάτεχνον, ἔστι δὲ περὶ αἰσχρῶν σχημάτων ; Σ Ar.~\textit{Ran.}~1328 Κυρήνη τις ἑταίρα ἐπίσημος, δωδεκαμήχανος ἐπικαλουμένη, διὰ τὸ τοσαῦτα σχήματα ἐν τῇ συνουσίᾳ ποιεῖν. ἔστι δὲ (παρὰ τὰ) ἐξ Ὑψιπύλης (Εὐριπίδου « ἀνὰ τὸ δωδεκαμήχανον ἄντρον. ») ; \textit{idem} Suid.~δ~1442 (add. Δωδεκαμήχανον· ἐπὶ τῶν παντοδαποῖς καὶ ποικίλοις χρωμένων ἤθεσι ; γέγονεν ; ἄστρον ; add. καὶ Δωδεκαμήχανος πόρνη, ιβ΄ σχήμασι χρωμένη, Κυρήνη τοὔνομα), Apostol.~6.41 (Δωδεκαμήχανος Κυρήνη· ἐπὶ τῶν παντοδαποῖς καὶ ποικίλοις χρωμένων ἤθεσιν tantum) ; Σ Ar.~\textit{Thesm.}~98 Κυρήνην: Τὴν ἑταίραν, περὶ ἧς πολλάκις φησί. τοιοῦτον δὲ τὸν Ἀγάθωνα καὶ οἱ ἄλλοι κωμῳδοῦσιν ; Hsch.~κ~4670 Κυρήνη· πόρνη τις ἥτις ἐκαλεῖτο δωδεκαμήχανος, διὰ τὸ τοσαῦτα σχήματα ἀφροδισίων ποιεῖν ; \textit{idem} Hsch.~δ~2706 (δωδεκαμήχανος· πόρνη τις ἐλέγετο … σχήματα ποιεῖν συνουσίας) ; Suid.~α~4261 Ἀστυάνασσα, Ἑλένης τῆς Μενελάου θεράπαινα· ἥτις πρώτη τὰς ἐν τῆ συνουσίᾳ κατακλίσεις εὗρε καὶ ἔγραψε περὶ σχημάτων συνουσιαστικῶν· ἣν ὕστερον παρεζήλωσαν Φιλαινὶς καὶ Ἐλεφαντίνη ; Ar.~\textit{Lys.}~231–232 οὐ στήσομαι λέαιν᾽ ἐπὶ τυροκνήστιδος ; Ath.~5.220f ἃ οὔτε Νικὼ ἡ Σαμία ἢ Καλλιστράτη ἡ Λεσβία ἢ Φιλαινὶς ἡ Λευκαδία ; Ath.~8.335b ἐγὼ Φιλαινὶς ἡ ᾽ πίβωτος ἀνθρώποις ἐνταῦθα γήρᾳ τῷ μακρῷ κεκοίμημαι ; Suid.~κ~396 Ξενοκλῆς γὰρ ὁ Καρκίνου δοκεῖ μηχανὰς καὶ τερατείας εἰσάγειν ἐν τοῖς δράμασιν· ὃν Πλάτων δωδεκαμήχανόν φησιν ; Ar.~\textit{Pax}~789–790 ναννοφυεῖς σφυράδων ἀποκνίσματα μηχανοδίφας ; Hsch.~μ~1953 μῦς λευκός· ἐπὶ τῶν θηλυδριῶν... κατωφερῶν (Philem. fr. 65, 1 K.–A.)}
\lp{98 Κυρήνην δ᾽ ὁρῶ]}{Ar.~\textit{Ach.}~118 ἐγᾦδ᾽ ὅς ἐστι, Κλεισθένης ὁ Σιβυρτίου ; Ar.~\textit{Nub.}~355 καὶ νῦν γ᾽ ὅτι Κλεισθένη εἶδον, ὁρᾷς, διὰ τοῦτ᾽ ἐγένοντο γυναῖκες ; Ar.~\textit{Thesm.}~235 ὁρᾷς σεαυτόν; — οὐ μὰ Δί᾽ ἀλλὰ Κλεισθένη ; Ar.~\textit{Nub.}~51–52 ἡ δ᾽ αὖ μύρου κρόκου καταγλωττισμάτων, / δαπάνης λαφυγμοῦ Κωλιάδος Γενετυλλίδος ; Ar.~\textit{Lys.}~42–48 αἳ καθήμεθ᾽ ἐξηνθισμέναι, / κροκωτοφοροῦσαι καὶ κεκαλλωπισμέναι … / τὰ κροκωτίδια καὶ τὰ μύρα χαἰ περιβαρίδες / χἤγχουσα καὶ τὰ διαφανῆ χιτώνια ; Ar.~\textit{Lys.}~149–151 εἰ γὰρ καθοίμεθ᾽ ἔνδον ἐντετριμμέναι, / κἀν τοῖς χιτωνίοισι τοῖς Ἀμοργίνοις / γυμναὶ παρίοιμεν δέλτα παρατετιλμέναι ; Ar.~\textit{Thesm.}~35 καὶ μὴν βεβίνηκας σύ γ᾽, ἀλλ᾽ οὐκ οἶσθ᾽ ἴσως ; Ar.~\textit{Thesm.}~261 φέρ᾽ ἔγκυκλον· — τουτὶ λάβ᾽ ἀπὸ τῆς κλινίδος ; Ar.~\textit{Thesm.}~1217–1218 τὴν γραῦν ἐρωτᾷς, ἣ ᾽φερεν τὰς πηκτίδας}
\lp{99 σίγα· μελῳδεῖν νῦν παρασκευάζεται]}{Ar.~\textit{Av.}~225–226 οὗτος. — τί ἔστιν; — οὐ σιωπήσει; — τί δαί; / — οὕποψ μελῳδεῖν αὖ παρασκευάζεται ; Soph.~\textit{OC}~111 σίγα. πορεύονται γὰρ οἵδε δή τινες}
\lp{99 μελῳδεῖν]}{Ar.~\textit{Av.}~1381 παῦσαι μελῳδῶν, ἀλλ᾽ ὅ τι λέγεις εἰπέ μοι ; Ar.~\textit{Thesm.}~67–69 καὶ γὰρ μελοποιεῖν ἄρχεται· χειμῶνος οὖν / ὄντος κατακάμπτειν τὰς στροφὰς οὐ ῥᾴδιον, / ἢν μὴ προίῃ θύρασι πρὸς τὸν ἥλιον ; Ar.~\textit{Thesm.}~130–133 ὡς ἡδὺ τὸ μέλος ὦ πότνιαι Γενετυλλίδες / καὶ θηλυδριῶδες καὶ κατεγλωττισμένον / καὶ μανδαλωτόν ; Zen.~5.40 Οὐδὲν πρὸς τὸν Διόνυσον: ἐπὶ τῶν τὰ μὴ προςήκοντα τοῖς ὑποκειμένοις λεγόντων ἡ παροιμία εἴρηται ; Ar.~\textit{Av.}~227–228 ἐποποῖ ποποποποποποποῖ, / ἰὼ ἰὼ ἰτὼ ἰτὼ ἰτὼ ἰτώ ; Ar.~\textit{Ran.}~312–313 οὗτος. — τί ἔστιν; — οὐ κατήκουσας; — τίνος; / — αὐλῶν πνοῆς ; Eur.~\textit{Hel.}~1110 σὲ τὰν ἀοιδοτάταν ὄρνιθα μελῳδὸν ; Pl.~\textit{Lg.}~790e καὶ οὐ σιγὴν ἀλλά τινα μελῳδίαν ; Pl.~\textit{Lg.}~812d ἄλλα δὲ τοῦ τὴν μελῳδίαν συνθέντος ποιητοῦ ; Hsch.~μ~770 μελῳδία· ἡδυφωνία ; Hsch.~μ~771 μελῳδοῦσι· μετ᾽ ᾠδῆς ὑμνοῦσιν ; Suid.~α~4402 Ἀοιδός: ὁ μελῳδός. καὶ ὁ ποιητής ; Suid.~ε~1562 Ἐξάρξατε: ἄρξασθε μελῳδίας, ὑμνήσατε ; Zen.~3.77 Ἐν γὰρ ἀμούσοις Καὶ κορυδὸς [κύκνου] φθέγγετ᾽ [ἀοιδότερος.] ; Eur.~\textit{IA}~1045 μελῳδοῖς Θέτιν ἀχήμασι τόν τ᾽ Αἰακίδαν ; Eur.~\textit{IT}~1104–1105 ἔνθα κύκνος μελῳ- / δὸς Μούσας θεραπεύει ; Eur.~fr.~188.1–3~Kannicht ἀλλ᾽ ἐμοὶ πιθοῦ· / παῦσαι ματάζων, καὶ πόνων εὐμουσίαν / ἄσκει· τοιαῦτ᾽ ἄειδε, καὶ δόξεις φρονεῖν ; [Eur.]~\textit{Rh.}~923 Μοῦσαι μεγίστην εἰς ἔριν μελῳδίας ; Hsch.~λ~978 λιγυρίζει· μελῳδεῖ ; Hsch.~χ~637 χορεύει· μελῳδεῖ. βακχεύει. ὀρχεῖται ; Suid.~ν~478 Νόμος: ὁ κιθαρῳδικὸς τρόπος τῆς μελῳδίας}
\lp{99 παρασκευάζεται]}{Aesch.~\textit{PV}~920 τοῖον παλαιστὴν νῦν παρασκευάζεται ; Ar.~\textit{Ran.}~848 τυφὼς γὰρ ἐκβαίνειν παρασκευάζεται ; Hsch.~π~668 παρασκευάζεται· εὐτρεπίζεται, ἑτοιμάζεται ; Ar.~\textit{Pax}~227 ἡμᾶς δὲ δὴ τί δρᾶν παρασκευάζεται ; Ar.~\textit{Ran.}~784 τί δῆθ᾽ ὁ Πλούτων δρᾶν παρασκευάζεται ; Aesch.~\textit{Ag.}~353 θεοὺς προσειπεῖν εὖ παρασκευάζομαι ; Hsch.~γ~961 γρυτεύεται· παρασκευάζεται}
\lp{99–100 μινυρισμός (παρεπιγραφή)]}{Σ Ar.~\textit{Av.}~222 αὐλεῖ: Τοῦτο παρεπιγέγραπται, δηλοῦν ὅτι μιμεῖταί τις τὴν ἀηδόνα ὡς ἔτι ἔνδον οὖσαν ἐν τῇ λόχμῃ ; Σ Ar.~\textit{Ran.}~1264 (διαύλιον προσαυλεῖ τις.) τοῦτο παρεπιγραφή, ὥσπερ καὶ ἄλλα πολλάκις. φασὶ δὲ διαύλιον λέγεσθαι ὅταν ἡσυχίας πάντων γενομένης ἔνδον ὁ αὐλητὴς ᾄσῃ ; Σ Ar.~\textit{Thesm.}~100 μεταξὺ δὲ τῶν δυοῖν ἀξιοῦσί τινες γράφειν μινυρισμός, ὡς πολλὰ τοιαῦτα παρεπιγράφεται ; Σ Ar.~\textit{Thesm.}~129 ὀλολύζει. παρεπιγραφὴ τοῦτο· εἰπόντος γὰρ τοῦ Ἀγάθωνος ὁ κηδεστὴς τοῦ Εὐριπίδου ὀλολύζει εἰς τὴν θηλύτητα αὐτοῦ καὶ μόνον οὐχὶ γυναικὸς εὐχήν ; Aesch.~\textit{Eum.}~117 (μυγμός.) ; Σ Ar.~\textit{Ran.}~269 μετὰ δὲ τὸν μδʹ κῶλον ἰαμβικὸν ἐφθημιμερὲς τὸ, αὐλεῖ τις ἔνδον· ὃ κατὰ παρεπιγραφήν φασι κεῖσθαι ; Σ Ar.~\textit{Thesm.}~277 παρεπιγραφή. ἐκκυκλεῖται ἐπὶ τὸ ἔξω τὸ θεσμοφόριον ; Σ Eur.~\textit{Or.}~1384 τινὲς τοῦτο παρεπιγραφὴν εἶναι ὡς εἰς τὰ κωμικὰ δράματα}
\lp{100 Μύρμηκος ἀτραπούς, ἢ τί διαμινύρεται;]}{\textit{Test.} Suid.~μ~1445 Μύρμηξ· μυρμήκων ἀτραποὺς ἢ τί διαμινυρίζεται. περὶ Ἀγάθωνός φησι τοῦ ποιητοῦ, ὡς λεπτὰ καὶ ἀγκύλα ἀνακρουομένου μέλη· τοιαῦται γὰρ αἱ τῶν μυρμήκων ὁδοί. ὁ δὲ Ἀγάθων ὑποκριτικὰ μέλη ἐποίει καὶ αὐτὸς ὑποκρίνεται. διὸ καὶ χορικὰ λέγεται}
\lp{100 Μύρμηκος]}{Suid.~φ~393 Φιλόξενος, Εὐλυτίδου, Κυθήριος, λυρικός. ἔγραψε διθυράμβους κδʹ· τελευτᾷ δὲ ἐν Ἐφέσῳ. οὗτος ἀνδραποδισθέντων τῶν Κυθήρων ὑπὸ Λακεδαιμονίων ἠγοράσθη ὑπὸ Ἀγεσύλου τινὸς καὶ ὑπ᾽ αὐτοῦ ἐτράφη καὶ Μύρμηξ ἐκαλεῖτο ; Ar.~\textit{Ran.}~1504–1506 καὶ τουτὶ τοῖσι πορισταῖς / Μύρμηκί θ᾽ ὁμοῦ καὶ Νικομάχῳ ; Eup.~fr.~416~K.-A. (Hsch.~ο~925) ὄνου γνάθος· Εὔπολις παίζει εἰς πολυφαγίαν. ἔστι δὲ καὶ τόπος οὕτω καλούμενος ⟨ἐν Λακωνικῇ⟩ ; \textit{idem} Diogenian.~6.100 (Ὄνου γνάθος· εἰς πολυφάγον) ; Harp. s.v.~Μελίτη κεκλῆσθαι δέ φησι τὸν δῆμον Φιλόχορος ἐν γ ἀπὸ Μελίτης θυγατρὸς κατὰ μὲν Ἡσίοδον Μύρμηκος ; App. Prov.~1.68 Γελοιότερον ἀπεργάζῃ τῶν τὰς κέγχρους ἀποτορεύειν ἐπιχειρούντων: ἢ καθάπερ τὸν Μυρμηκίδην ἀντιπραττόμενον τῇ Φειδίου τέχνῃ ; Posidipp. Com. frr.~19–20~K.-A. (Μύρμηξ) Ποσιδίππου ἐκ Μύρμηκος (Stob. 4.51.16 et 4.35.21) ; \textit{Od.}~13.408 πὰρ Κόρακος πέτρῃ ἐπί τε κρήνῃ Ἀρεθούσῃ ; Hdt.~7.183.2 περὶ τὸ ἕρμα τὸ μεταξὺ ἐὸν Σκιάθου τε καὶ Μαγνησίης, καλεόμενον δὲ Μύρμηκα ; X.~\textit{Hell.}~5.4.50 ἀντετάξαντο ἐπὶ Γραὸς στήθει ; Ael.~\textit{VH}~1.17 ταῦτα ἄρα ἐστὶ τὰ θαυμαζόμενα Μυρμηκίδου τοῦ Μιλησίου καὶ Καλλικράτους τοῦ Λακεδαιμονίου μικρὰ ἔργα ; Hsch.~μ~1902 μύρμηκας· ἔξω τοῦ ζῴου καὶ οἱ πυκτικοὶ ἱμάντες. καὶ τῶν θαλασσίων πετρῶν εἶδος. καὶ ζῷα ὑπόπτερα, ὡς κύνες τὸ μέγεθος ; Hsch.~μ~1906 μυρμηδόνες· οἱ μύρμηκες ὑπὸ Δωριέων ; Hsch.~μ~1909 μύρμοι· μύρμηκες ; Suid.~μ~1443 Μυρμηκανθρώποις Φερεκράτης γράφει ; Zen.~4.23 Ἢ τρὶς ἓξ ἢ τρεῖς κύβοι: ἡ παροιμία παρὰ Φερεκράτει ἐν τοῖς Μυρμηκανθρώποις}
\lp{100 Μύρμηκος ἀτραπούς]}{\textit{Test.} Hsch.~μ~1904 Μύρμηκος ἀτραπούς· Ἀθήνησιν ἐν Σκαμβωνιδῶν ἔστι Μύρμηκος ἀτραπός, ἀπὸ ἥρωος Μύρμηκος ὀνομαζόμενη ; Phot.~μ~606 Μύρμηκος ἀτραπός· ἐν Θεσμοφοριαζούσαις (100) Ἀριστοφάνης <***> ἀντιφωνεῖ ἤ τι τοιοῦτο διὰ τὴν βραχύτητα· ἔστιν Ἀθήνησι Μύρμηκος ἀτραπὸς ἐν Σκαμβωνιδῶν ἀπὸ Μύρμηκος τοῦ Μελανίππου τοῦ Κύκλωπος τοῦ Ζευξίππου ; Pherecr.~fr.~155.21–23~K.-A. Μιλήσιός τις πυρρίας / κακά μοι παρέσχεν οὗτος, ἅπαντας οὓς λέγω / παρελήλυθ᾽ ᾄδων ἐκτραπέλους μυρμηκιάς ; Plaut.~\textit{Men.}~888 atque eccum incedit. move formicinum gradum ; \textit{AP}~11.78.3–4 ὄντως μυρμήκων τρυπήματα λοξὰ καὶ ὀρθά, / γράμματα τῶν λυρικῶν Λύδια καὶ Φρύγια ; Arist.~\textit{HA}~9.38 Ἡ μὲν οὖν τῶν μυρμήκων ἐργασία πᾶσίν ἐστιν ἐπιπολῆς ἰδεῖν, καὶ ὡς ἀεὶ μίαν ἀτραπὸν πάντες βαδίζουσι ; Plu.~\textit{De\,soll.\,an.}~968a–b λέγουσι δ᾽ οὖν οὐκ εὐθεῖαν εἶναι τὴν ἀπὸ τῆς ὀπῆς κάθοδον οὐδ᾽ εὔπορον ἄλλῳ θηρίῳ διεξελθεῖν, ἀλλὰ καμπαῖς καὶ στρεβλότησι κεκλασμένας ὑποπορεύσεις καὶ ἀνατρήσεις ἐχούσαις εἰς τρεῖς κοιλότητας ἀποτελευτῶσαν ; Ael.~\textit{NA}~6.43 μυρμήκων δὲ ἐν γεωρυχίᾳ ποικίλας ἀτραποὺς καὶ ἑλιγμοὺς καὶ περιόδους οὔπω ἴσασι. σοφίᾳ δὲ ἄρα τὴν ὑπόγειον οἰκονομίαν τήνδε ἀπεργάζονται σκολιωτάτην ; Ael.~\textit{NA}~16.15 ἐν αὐτοῖς δὲ περιόδους τινὰς καὶ ὡς εἰπεῖν σύριγγας Αἰγυπτίας ἢ λαβυρίνθους Κρητικοὺς … διατρήσαντες … οὐκ εὐθυτενῆ καὶ ῥᾴδια παρελθεῖν ἀλλ᾽ ἑλιγμοῖς καὶ διατρήσεσι λοξά ; Σ Ar.~\textit{Thesm.}~100 ὡς λεπτὰ καὶ ἀγκύλα ἀνακρουομένου μέλη τοῦ Ἀγάθωνος· τοιαῦται γὰρ αἱ τῶν μυρμήκων ὁδοί ; Hsch.~μ~1903 μυρμηκιά· τάσσεται δὲ καὶ ἐπὶ διδασκαλείου καὶ συμφοιτήσεως. ἦν δὲ καὶ λόγος, ὡς ὁ πτύσας εἰς μυρμηκιὰν οἰδεῖ τὰ χείλη, ὡς ὁ Δεινολόχος (fr. 12 K.–A.) ; \textit{idem} Tosi n°~707 (Ὁ πτύσας εἰς μυρμηκιὰν οἰδεῖ τὰ χείλη tantum), Erasm.~\textit{Adag.}~4.6.80 (Qui inspuerit in cavernam formicarum, \&c. … Quod qui inspuerit in agmen formicarum, huic intumescunt labra) ; Hsch.~μ~1905 Μυρμήκων ὁδοί· Ἀθήνησι τόπος. καὶ αἱ μονόκωλοι τρίβοι. ἀπὸ τῆς τοῦ ζῴου ὁμοιότητος κατὰ τὴν ὁδὸν γινομένου ; Zen.~1.11 Ἀγαθῶν μυρμηκιά: ἐπὶ πλήθους εὐδαιμονίας ; \textit{idem} Diogenian.~1.10 (Ἀγαθῶν θάλασσα· ἐπὶ πλήθους ἀγαθῶν. Ὡς καὶ τὸ, Ἀγαθῶν σωρός, καὶ, Ἀγαθῶν μυρμηκία. Ἐπὶ πλήθους γὰρ εὐδαιμονίας καὶ τοῦτο), Suid.~α~123 (Ἀγαθῶν ἀγαθίδες· τάττεται ἡ παροιμία παρὰ τοῖς κωμικοῖς ἐπὶ τῶν πολλῶν ἀγαθῶν. καὶ Ἀγαθῶν θάλασσα, ἐπὶ πλήθους ἀγαθῶν. καὶ Ἀγαθῶν μυρμηκίαι), Erasm.~\textit{Adag.}~1.3.32 (Bonorum myrmecia, de immensis opibus. Nam Myrmecia Graece significat cavernam, aut agmen formicarum) ; Ar.~\textit{Thesm.}~1175 σὺ δ᾽ ὦ Τερηδὼν ἐπαναφύσα Περσικόν ; [Aesch.]~\textit{PV}~452–453 κατώρυχες δ᾽ ἔναιον ὥστ᾽ ἀήσυροι / μύρμηκες ἄντρων ἐν μυχοῖς ἀνηλίοις ; Σ Ar.~\textit{Vesp.}~352 Κράτης δὲ μύρμηκα φησὶ τὸν σέρφον … ἀφ᾽ οὗ ἡ παροιμία ἔνεστι κἀν μύρμηκι κἀν σέρφῳ χολή … — σέρφῳ: παροιμία «οὐδὲ μύρμηκι ὁδός ἐστι.» R ; Hsch.~μ~1907 μυρμηκιών· ξυνοικία τῶν μυρμήκων ; Phot.~μ~605 μυρμηκιᾶν· τοὺς εἰς μυρμηκιὰν λαλήσαντας ἑλκοῦσθαι τὴν γλῶττάν φασι ; Ar.~\textit{Vesp.}~352 πάντα πέφαρκται κοὐκ ἔστιν ὀπῆς οὐδ᾽ εἰ σέρφῳ διαδῦναι ; Ar.~\textit{Thesm.}~427 ἐδίδαξε θριπήδεστ᾽ ἔχειν σφραγίδια ; Apollod. Epit.~1.15 ὁ δὲ ἐξάψας μύρμηκος λίνον καὶ τρήσας τὸν κοχλίαν εἴασε δι᾽ αὐτοῦ διελθεῖν ; Ael.~\textit{NA}~6.50 ὁρᾷ ἐξ ἀτραποῦ τινος ἑτέρας νεκρὸν μύρμηκα μύρμηκας ἄλλους κομίζοντας ; Zen.~3.70 Ἔνεστι κἀν μύρμηκι χολή· παρεγγυᾷ μηδὲ τῶν μικρῶν καταφρονεῖν ; \textit{idem} Phot.~ε~901 Theodoridis}
\lp{100 Μύρμηκος ἀτραπούς (ἐν Σκαμβωνιδῶν)]}{Plu.~\textit{Prov.Alex.}~1.18 Τὸ σκαμβὸν ξύλον οὐδέποτε ὀρθοῦν: ἐπὶ τῶν πειρωμένων διάστροφα ξύλα κατευθύνειν ; \textit{idem} Phot.~τ~396 (ὀρθόν· αὕτη δημώδης ἐστὶ καὶ φανερά· τάττεται δὲ καὶ κατὰ τῶν διάστροφα ξύλα ἢ πράγματα πειρωμένων κατευθύνειν, μηδὲν ὠφελούντων), Suid.~σ~530 (Σκαμβὸν ξυλὸν οὐδέποτ᾽ ὀρθόν· ζήτει ἐν τῷ τὸ σκαμβόν), Macar.~6.17 (Ξύλον ἀγκύλον οὐδέποτ᾽ ὀρθόν), Macar.~7.69 (Σκαμβὸν ξύλον οὐδέποτ᾽ ὀρθόν· πρὸς τοὺς ἀκάμπτους καὶ σκληρούς) ; Ar.~\textit{Vesp.}~81 Νικόστρατος δ᾽ αὖ φησιν ὁ Σκαμβωνίδης ; Plu.~\textit{Alc.}~22.3 Θεσσαλὸς Κίμωνος Λακιάδης Ἀλκιβιάδην Κλεινίου Σκαμβωνίδην εἰσήγγειλεν ; Hsch.~ρ~48 ῥαιβόν· ἐπικαμπές, τὸ μὴ ὀρθόν, καμπύλον, στρεβλόν, σκαμβόν ; Hsch.~σ~831 σκαμβάλυξ· σκαμβός, στρεβλός ; Harp. s.v.~Σκαμβωνίδαι Σκαμβωνίδαι· Λυκοῦργος ἐν τῇ διαδικασίᾳ Κροκωνιδῶν. ἔστι δὲ δῆμος τῆς Λεοντίδος}
\lp{100 ἀτραπούς]}{Hsch.~α~8138 ἀτραπός· ὁδὸς τετριμμένη, μὴ ἔχουσα ἐκτροπάς, ἀλλ᾽ εὐθεῖα ; Macar.~6.21 Ὁδοῦ παρούσης τὴν ἀτραπὸν ζητεῖς: ἐπὶ τῶν προδήλων ; \textit{idem} Apostol.~12.34 (ἐπὶ τῶν καταλιμπανόντων τὰ δῆλα καὶ ζητούντων τὰ ἀφανῆ), App. Prov.~4.12 (παρούσης, ; μὴ ζήτει) ; Ar.~fr.~47~K.-A. ὅρμου παρόντος τὴν ἀτραπὸν κατερρύην ; \textit{idem} Phot.~ο~37 (add. ὁδοῦ παρούσης τὴν ἀτραπὸν ζητεῖς· λεγόμενον. Ἀριστοφάνης Ἀναγύρῳ), Suid.~ο~48 (add. Ὁδοῦ παρούσης τὴν ἀτραπὸν ζητεῖς· τὸ λεγόμενον. Ἀριστοφάνης Ἀναγύρῳ ; παρόντος,) ; Synag.~α~1062 ἀτραπούς· ὁδοὺς τετριμμένας ; \textit{idem} Suid.~α~4380 (add. Ἀτραπιτός· ἡ ὁδός. καὶ) ; Ar.~\textit{Nub.}~76 μίαν ηὗρον ἀτραπὸν δαιμονίως ὑπερφυᾶ ; Ar.~\textit{Av.}~21–22 οὐ γάρ ἐστ᾽ ἐνταῦθά τις / ὁδός. — οὐδὲ μὰ Δί᾽ ἐνταῦθά γ᾽ ἀτραπὸς οὐδαμοῦ ; Ar.~\textit{Ran.}~123 ἀλλ᾽ ἔστιν ἀτραπὸς ξύντομος τετριμμένη ; X.~\textit{Cyr.}~1.3.4 ὑμεῖς δὲ εἰς μὲν τὸ αὐτὸ ἡμῖν σπεύδετε, πολλοὺς δέ τινας ἑλιγμοὺς ἄνω καὶ κάτω πλανώμενοι μόλις ἀφικνεῖσθε ; Hsch.~α~8137 ἀτραπιτοῖς· ὁδοῖς ; Plaut.~\textit{Men.}~247 In scirpo nodum quaeris ; \textit{idem} Ter.~\textit{An.}~941 (nodum in scirpo quaeris)}
\lp{100 ἀτραπούς (image musicale)]}{Ar.~\textit{Nub.}~969–971 εἰ δέ τις αὐτῶν βωμολοχεύσαιτ᾽ ἢ κάμψειέν τινα καμπήν, / οἵας οἱ νῦν τὰς κατὰ Φρῦνιν ταύτας τὰς δυσκολοκάμπτους, / ἐπετρίβετο τυπτόμενος πολλὰς ὡς τὰς Μούσας ἀφανίζων ; Ar.~\textit{Thesm.}~53 κάμπτει δὲ νέας ἁψῖδας ἐπῶν ; Ar.~\textit{Thesm.}~68 ὄντος κατακάμπτειν τὰς στροφὰς οὐ ῥᾴδιον ; Pherecr.~fr.~31~K.-A. ἀτραπίζοντες τὰς ἁρμονίας διὰ πασῶν ; Plu.~\textit{Quaest.\,conv.}~3.1.1 (645d–646a) κατηγοροῦντα τοῦ καλοῦ Ἀγάθωνος, ὃν πρῶτον εἰς τραγῳδίαν φασὶν ἐμβαλεῖν καὶ ὑπομῖξαι τὸ χρωματικόν, ὅτε τοὺς Μυσοὺς ἐδίδασκεν ; Ar.~\textit{Nub.}~333 κυκλίων τε χορῶν ᾀσματοκάμπτας ἄνδρας μετεωροφένακας ; Ar.~\textit{Av.}~1374 πέτομαι δ᾽ ὁδὸν ἄλλοτ᾽ ἐπ᾽ ἄλλαν μελέων ; Pherecr.~fr.~155.8–9,~14–15,~26–28~K.-A. Κινησίας δὲ μ᾽ ὁ κατάρατος Ἀττικός, / ἐξαρμονίους καμπὰς ποιῶν ἐν ταῖς στροφαῖς … Φρῦνις δ᾽ ἴδιον στρόβιλον ἐμβαλών τινα / κάμπτων με καὶ στρέφων ὅλην διέφθορεν, … ἐξαρμονίους ὑπερβολαίους, τ᾽ ἀνοσίους / καὶ νιγλάρους, ὥσπερ τε τὰς ῥαφάνους ὅλην / καμπῶν με κατεμέστωσε ; Pi.~\textit{O.}~1.110 ἐπίκουρον εὑρὼν ὁδὸν λόγων ; \textit{h.Merc.}~451 τῇσι χοροί τε μέλουσι καὶ ἀγλαὸς οἶμος ἀοιδῆς ; Plu.~\textit{De\,def.\,or.}~418c (Emp. fr.~24~DK) ἀλλ᾽ ἵνα μὴ τὸ Ἐμπεδόκλειον ποιεῖν δόξω „κορυφὰς ἑτέρας ἑτέρῃσι / προσάπτων μύθων, μήτε λέγειν ἀτραπὸν μίαν“ ; Hsch.~α~8136 ἀτραπίζειν· βαδίζειν. ὁδοιπορεῖν (Pherecr. fr. 31 K.–A.) ; Zen.~2.15 Ἄειδε τὰ Τέλληνος: οὗτος ὁ Τέλλην ἐγένετο αὐλητὴς καὶ μελῶν ἀνυποτάκτων ποιητής ; Diogenian.~7.40 Οὐ πατρικὰ αὐλεῖ μέλη: ἐπὶ τῶν οὐχ ὁμοίως τοῖς πατράσι βιούντων ; App. Prov.~4.94 Τὸν αὐλητὴν αὐλεῖν: αὐλητὴς ἐγένετο μὴ πάνυ τοῖς αὐλητικοῖς ἐμμένων νόμοις, ἀλλὰ παρακινῶν· ὅθεν ἡ παροιμία ; Ar.~\textit{Pax}~733 ἣν ἔχομεν ὁδὸν λόγων εἴπωμεν ὅσα τε νοῦς ἔχει ; Hdt.~1.95.1 ἐπιστάμενος περὶ Κύρου καὶ τριφασίας ἄλλας λόγων ὁδοὺς φῆναι}
\lp{100 διαμινύρεται]}{Ar.~\textit{Vesp.}~219–220 λύχνους ἔχοντες καὶ μινυρίζοντες μέλη / ἀρχαῖα μελισιδωνοφρυνιχήρατα ; Ar.~\textit{Av.}~1414 ὅδ᾽ αὖ μινυρίζων δεῦρό τις προσέρχεται ; Ar.~\textit{Eccl.}~880 ἀργός, μινυρομένη τι πρὸς ἐμαυτὴν μέλος ; Phryn. Com. fr.~74~K.-A. καὶ νιγλάρους θρηνεῖν, ἐν οἷσι Λάμπρος ἐναπέθνῃσκεν, ἄνθρωπος ὢν ὑδατοπότης, μινυρὸς ὑπερσοφιστής, Μουσῶν σκελετός, ἀηδόνων ἠπίαλος, ὕμνος Ἅιδου ; Aesch.~\textit{Ag.}~16 ὅταν δ᾽ ἀείδειν ἢ μινύρεσθαι δοκῶ ; Soph.~\textit{OC}~671–672 ἁ λίγεια μινύρεται / θαμίζουσα μάλιστ᾽ ἀηδών ; Pl.~\textit{R.}~411a διὰ τῶν ὤτων ὥσπερ διὰ χώνης ἃς νυνδὴ ἡμεῖς ἐλέγομεν τὰς γλυκείας τε καὶ μαλακὰς καὶ θρηνώδεις ἁρμονίας, καὶ μινυρίζων τε καὶ γεγανωμένος ὑπὸ τῆς ᾠδῆς ; Ammon.~\textit{Diff.}~321 μινυρίζειν καὶ μινύρεσθαι διαφέρει. μινυρίζειν μὲν γάρ ἐστι ἠρέμα προσᾴδειν, μινύρεσθαι δὲ τὸ θρηνεῖν· τὸ δ᾽ αὐτὸ καὶ κινύρεσθαί φασιν ; Hsch.~μ~1410 μινύρεται· κλαίει, θρηνεῖ, ὀδύρεται ἢ εὐφώνως λέγει ; \textit{idem} Synag.~μ~228 (om. ὀδύρεται ἢ), Phot.~μ~463 (om. ὀδύρεται ἢ), Suid.~μ~1102 (add. ᾄδει ; om. ὀδύρεται ἢ ; add. Σοφοκλῆς· ἔνθα λίγεια μινύρεται θαμίζουσα μάλιστ᾽ ἀηδὼν) ; Aesch.~\textit{Ag.}~1165 δυσαλγεῖ τύχᾳ μινυρὰ [κακὰ] θρεομένας ; \textit{Il.}~5.889 μή τί μοι ἀλλοπρόσαλλε παρεζόμενος μινύριζε ; \textit{Od.}~4.719 οἴκτρ᾽ ὀλοφυρομένη· περὶ δὲ δμῳαὶ μινύριζον ; Theoc.~\textit{Ep.}~4.11 ξουθαὶ δ᾽ ἀδονίδες μινυρίσμασιν ἀνταχεῦσι ; \textit{AP}~9.70.1 Τραυλὰ μινυρομένα, Πανδιονὶ παρθένε, φωνᾷ ; Plu.~\textit{Num.}~4.7 ποιηταῖς δὲ καὶ λυρικοῖς μινυρίζουσιν ; Σ Ar.~\textit{Eccl.}~880 μινυρομένη : Ἀντὶ τοῦ ἠρέμα ᾄδουσα ; Σ Od.~4.719 μινύριζον] ποία λέξις; ἡσύχως ἔκλαιον καὶ μικρῶς· μινυὸν γὰρ τὸ μικρόν ; Σ Soph.~\textit{OC}~671 [μινύρεται] ᾄδει ; Hsch.~μ~1411 μινυρίζει· ὀλίγῃ φωνῇ καὶ οἰκτρᾷ χρῆται ; Hsch.~μ~1412 μινυρομένη· θρηνοῦσα. καὶ τὰ ὅμοια ; Hsch.~μ~1413 μινυρόν· μικρόν, ὀλίγον. τῇ φωνῇ μινυρίζον ; Suid.~μ~1103 Μινυρίζω: θρηνῶ· ἢ τὸ ἠρέμα ᾄδω. παρὰ τὸ μεῖον. Πλάτων ἐν τρίτῳ Πολιτείας· μινυρίζων τε καὶ γεγανωμένος ὑπὸ τῆς ᾠδῆς. ἀντὶ τοῦ οἰκτρᾷ φωνῇ χρώμενος καὶ θρηνῶν. ἢ ὀλίγῃ φωνῇ χρώμενος. Ἀριστοφάνης· μινυρίζοντες μέλη ἀρχαιομελισιδωνοφρυνιχήρατα ; \textit{idem} Phot.~μ~464 (μινυρίζω· θρηνῶ ἢ τὸ ἠρέμα ᾄδω … Πλάτων ἐν γ´ Πολιτείας· „μινυρίζων τε καὶ γεγανωμένως ὑπὸ τῆς ᾠδῆς“) ; Timoth.~\textit{PMG}~802.3 μακάριος ἦσθα, Τιμόθε᾽, ὅτε κᾶρυξ / εἶπε· νικᾶι Τιμόθεος / Μιλήσιος τὸν Κάμωνος τὸν ἰωνοκάμπταν ; Theoc.~13.12 οὐδ᾽ ὁπόκ᾽ ὀρτάλιχοι μινυροὶ ποτὶ κοῖτον ὁρῷεν ; Lyr. adesp.~PMG~1036.3 ἀνά⟨ρ⟩θροις μινυρίσμ[ασιν ; Philox. (Δεῖπνον) ap. Ath.~4.146f–147e καὶ νῶτος ἐσῆλθε καὶ ὀσφὺς καὶ μινυρίγματα θερμά ; Arist.~\textit{HA}~618b31 οὐ γὰρ μινυρίζει οὐδὲ λέληκεν ; Plu.~\textit{De\,aud.}~9 (42c–e) οὐδ᾽ οἴεσθαι δεῖν ἐκ σχολῆς ἀπιέναι φιλοσόφου μινυρίζοντα καὶ γεγανωμένον ; Plu.~\textit{Quomodo\,adul.}~56f οὐκ Ἀπόλλωνες μὲν ἂν μινυρίσωσι ; [Plu.]~\textit{De\,mus.}~1138c δῆλον οὖν ὅτι οἱ παλαιοὶ οὐ δι᾽ ἄγνοιαν ἀλλὰ διὰ προαίρεσιν ἀπείχοντο τῶν κεκλασμένων μελῶν ; Clem.Al.~\textit{Strom.}~2.8.37 τερετίσματα καὶ μινυρίσματα ἀναπλάσσοντες ; Σ Ar.~\textit{Vesp.}~219 μινυρίζοντες: ᾄδοντες ; Hsch.~κ~2713 κιναθίζειν· ἰδιάζειν, ἀποθησαυρίζειν κατὰ μικρὸν συλλέγοντα. ἔνιοι μινυρίζειν. καὶ κινεῖν ; Hsch.~μ~1414 μινύρονται· προφωνοῦσι, προλέγουσι}
\lp{100 ἢ τί διαμινύρεται;]}{Ar.~\textit{Lys.}~128 ποιήσετ᾽ ἢ οὐ ποιήσετ᾽; ἢ τί μέλλετε}
